\documentclass[11pt]{article}
\usepackage{xcolor}
\definecolor{repositorylink}{HTML}{0366d6}

\usepackage[T1]{fontenc}
\usepackage{lmodern}
\usepackage[letterpaper,margin=1in]{geometry}
\usepackage{microtype}
\usepackage{mathtools}
\usepackage{amssymb}
\usepackage{mathrsfs}
\usepackage{amsthm}
\usepackage{graphicx}
\usepackage{academicons}
\IfFileExists{fontawesome5.sty}{%
	\usepackage{fontawesome5}%
	\newcommand{\GitHubIcon}{\mbox{{\small\faGithub}}\nobreakspace}%
}{%
	\newcommand{\GitHubIcon}{}%
}
\usepackage{flafter}
\usepackage[font=small,labelfont=bf,skip=6pt]{caption}
\usepackage{needspace}
\usepackage{booktabs}
\usepackage{tabularx}
\usepackage{array}
\usepackage{enumitem}
\usepackage{tikz}
\usetikzlibrary{arrows.meta,positioning,shapes.geometric}
\usepackage{aliascnt}
\usepackage[hidelinks]{hyperref}
\hypersetup{
	pdftitle={A Full Complexity Dichotomy for Complex-Valued Boolean Holant Problems},
	pdfauthor={Chenghua Liu}
}
\usepackage[capitalise,nameinlink,noabbrev]{cleveref}
\usepackage{authblk}

\renewcommand\Authfont{\normalsize}
\renewcommand\Affilfont{\small}
\setlength{\affilsep}{0.5em}
\setlength{\emergencystretch}{2em}
\setlist[itemize]{leftmargin=1.6em,itemsep=0.25em,topsep=0.3em}
\setlist[enumerate]{leftmargin=1.8em,itemsep=0.3em,topsep=0.3em}
\Urlmuskip=0mu plus 1mu\relax
\allowdisplaybreaks[2]

\theoremstyle{plain}
\newtheorem{theorem}{Theorem}[section]

\newaliascnt{lemma}{theorem}
\newtheorem{lemma}[lemma]{Lemma}
\aliascntresetthe{lemma}

\newaliascnt{proposition}{theorem}
\newtheorem{proposition}[proposition]{Proposition}
\aliascntresetthe{proposition}

\newaliascnt{corollary}{theorem}
\newtheorem{corollary}[corollary]{Corollary}
\aliascntresetthe{corollary}

\newaliascnt{claim}{theorem}
\newtheorem{claim}[claim]{Claim}
\aliascntresetthe{claim}

\theoremstyle{definition}
\newaliascnt{definition}{theorem}
\newtheorem{definition}[definition]{Definition}
\aliascntresetthe{definition}

\newaliascnt{procedure}{theorem}
\newtheorem{procedure}[procedure]{Decision Procedure}
\aliascntresetthe{procedure}

\theoremstyle{remark}
\newaliascnt{remark}{theorem}
\newtheorem{remark}[remark]{Remark}
\aliascntresetthe{remark}

\newenvironment{proofstrategy}
{\par\noindent\textit{Proof strategy.}\ }
{\par}

\crefname{theorem}{Theorem}{Theorems}
\Crefname{theorem}{Theorem}{Theorems}
\crefname{lemma}{Lemma}{Lemmas}
\Crefname{lemma}{Lemma}{Lemmas}
\crefname{proposition}{Proposition}{Propositions}
\Crefname{proposition}{Proposition}{Propositions}
\crefname{corollary}{Corollary}{Corollaries}
\Crefname{corollary}{Corollary}{Corollaries}
\crefname{claim}{Claim}{Claims}
\Crefname{claim}{Claim}{Claims}
\crefname{definition}{Definition}{Definitions}
\Crefname{definition}{Definition}{Definitions}
\crefname{procedure}{Procedure}{Procedures}
\Crefname{procedure}{Procedure}{Procedures}
\crefname{remark}{Remark}{Remarks}
\Crefname{remark}{Remark}{Remarks}

\newcommand{\AlgNums}{\overline{\mathbb Q}}
\newcommand{\Ftwo}{\mathbb F_2}
\newcommand{\cF}{\mathcal F}
\newcommand{\cG}{\mathcal G}
\newcommand{\cH}{\mathcal H}
\newcommand{\cA}{\mathscr A}
\newcommand{\cP}{\mathscr P}
\newcommand{\cL}{\mathscr L}
\newcommand{\cM}{\mathscr M}
\newcommand{\cT}{\mathscr T}
\newcommand{\cC}{\mathscr C}
\newcommand{\cE}{\mathscr E}
\newcommand{\angles}[1]{\left\langle #1\right\rangle}
\newcommand{\Holant}{\operatorname{Holant}}
\newcommand{\KHolant}{\operatorname{K\text{-}Holant}}
\newcommand{\EOTract}{\mathsf{EOTract}}
\newcommand{\Tract}{\mathsf{Tract}}
\newcommand{\FP}{\mathsf{FP}}
\newcommand{\SharpP}{\#\mathsf P}
\newcommand{\PGL}{\operatorname{PGL}}
\newcommand{\GL}{\operatorname{GL}}
\newcommand{\Mat}{\operatorname{Mat}}
\newcommand{\diag}{\operatorname{diag}}
\newcommand{\Span}{\operatorname{span}}
\newcommand{\rank}{\operatorname{rank}}
\newcommand{\wt}{\operatorname{wt}}
\newcommand{\supp}{\operatorname{supp}}
\newcommand{\arity}{\operatorname{arity}}
\newcommand{\irr}{\operatorname{Irr}}
\newcommand{\EO}{\mathrm{EO}}
\newcommand{\GE}{\mathrm{GE}}
\newcommand{\Eq}{\mathrm{EQ}}
\newcommand{\Neq}{\mathrm{NEQ}}
\newcommand{\Id}{I}
\newcommand{\trans}{\mathsf T}
\newcommand{\eps}{\varepsilon}
\newcommand{\one}{\mathbf 1}
\newcommand{\zero}{\mathbf 0}
\newcommand{\bzero}{\boldsymbol 0}
\newcommand{\bone}{\boldsymbol 1}
\newcommand{\DeltaZero}{\Delta_0}
\newcommand{\DeltaOne}{\Delta_1}
\newcommand{\RMcore}{\mathrm{RM}_8}
\newcommand{\leT}{\leq_{\mathrm T}}
\newcommand{\equivT}{\equiv_{\mathrm T}}
\newcommand{\lemref}[1]{Lemma~\ref{#1}}
\DeclarePairedDelimiter{\abs}{\lvert}{\rvert}
\DeclarePairedDelimiter{\norm}{\lVert}{\rVert}

\newcolumntype{L}[1]{>{\raggedright\arraybackslash}p{#1}}

\begin{document}
	
	\title{A Full Complexity Dichotomy for Complex-Valued\\[0.15em]
		Boolean Holant Problems}
	\author[1,*]{Chenghua Liu}
	\author[2,\textdagger]{Boning Meng}
	\author[1,\textdaggerdbl]{Juqiu Wang}
	
	\affil[1]{Institute of Software, Chinese Academy of Sciences, Beijing, China}
	\affil[2]{University of Regensburg, Regensburg, Germany}
	
	\affil[ ]{%
		\textsuperscript{*}\texttt{liuch.russell@gmail.com}\quad
		\textsuperscript{\textdagger}\texttt{mengboning2013@gmail.com}\quad
		\textsuperscript{\textdaggerdbl}\texttt{wangjuqiu104@gmail.com}%
	}
	
	\date{}
	
	\maketitle
	
	\begin{abstract}
		We prove a complexity dichotomy for Boolean Holant problems defined by
		arbitrary finite sets of algebraic complex-valued signatures.  The tractable
		cases are characterized by an explicit, decidable criterion.
	\end{abstract}
	
	\section*{Acknowledgments}
	Chenghua Liu is sincerely grateful to Zhengfeng Ji and Mingsheng Ying
	for their continued encouragement and generous support.
	Boning Meng and Juqiu Wang warmly thank Mingji Xia for his
	invaluable guidance and dedicated mentorship.
	The authors used OpenAI's ChatGPT in preparing this manuscript, including
	exploratory proof development, exact symbolic and finite-case checks,
	language editing, and \LaTeX{} preparation. All companion verification code
	was written by OpenAI Codex and subsequently reviewed by the authors.
	All mathematical claims, computer-assisted certificates, and the final
	presentation were reviewed and finalized by the authors, who take full
	responsibility for the content.
	
	\section*{Code Availability}
	The verification code for the finite calculations in
	Appendix~\ref{sec:exact-finite-certificates} is available in the
	\path{certificates/} directory of the
	\href{https://github.com/liuchliuch/Full-Holant-Dichotomy}%
	{\textcolor{repositorylink}{\GitHubIcon\textit{GitHub repository:
	Full-Holant-Dichotomy}}}. The complete suite is run from the repository
	root with \texttt{python3 certificates/run\_all\_exact.py}.
	The \path{README.md} in that directory gives requirements and reproduction
	commands; \path{theorem-to-script-map.md} identifies the statements checked
	by each program and distinguishes supplementary checks from calculations
	used in the proof.
	
	\setcounter{tocdepth}{2}
	\tableofcontents
	
	\section{Introduction}
	\label{sec:introduction}
	
	A Holant problem asks for the value of a graph tensor-network partition
	function: each vertex is assigned a local function, called a signature,
	of its incident Boolean edge variables.  One sums the product of these
	functions over all edge assignments.  The framework
	includes weighted graph homomorphisms, counting CSPs, Eulerian orientations,
	and vertex models.  A \emph{holographic transformation} is
	an invertible simultaneous change of basis on the local tensors and wires
	\cite{Valiant2008,CaiLuXia2009}.
	
	We prove a dichotomy for finite algebraic complex-valued Boolean Holant
	problems on general graphs.  Our proof treats unrestricted all-even
	signature sets, in which every positive-arity signature has even arity and
	no auxiliary binary signature is prescribed, and combines this analysis
	with known dichotomies.  We first review the graph-homomorphism,
	counting-CSP, and Holant classifications that precede this result.
	
	\subsection{From counting CSP to the full Holant dichotomy}
	\label{subsec:intro-history}
	
	Weighted graph homomorphisms
	(\(\#\mathrm{GH}\)) and counting constraint satisfaction problems
	(\(\#\mathrm{CSP}\)) provide earlier examples of counting dichotomies on
	general graphs.  Their classifications developed in
	three directions: from Boolean to arbitrary finite domains, from symmetric to
	arbitrary local functions, and from Boolean-valued through nonnegative,
	signed-real, and complex weights.  Cai and Chen's general-domain complex
	classification covers both problems \cite{CaiChen2017}.
	\Cref{tab:intro-gh-csp-history} records these classifications.  A dash
	indicates that no separate reference is listed for that case; the case is
	covered by a more general entry in the table.
	
	\begin{table}[!htbp]
		\centering
		\caption{Development of general-graph dichotomies for weighted graph
			homomorphisms and counting CSP.}
		\label{tab:intro-gh-csp-history}
		\small
		\setlength{\tabcolsep}{3.5pt}
		\renewcommand{\arraystretch}{1.16}
		\begin{tabularx}{\textwidth}{@{}l
				>{\raggedright\arraybackslash}X
				>{\raggedright\arraybackslash}X
				>{\raggedright\arraybackslash}X
				>{\raggedright\arraybackslash}X@{}}
			\toprule
			Weight range & Symmetric \(\#\mathrm{GH}\) & \(\#\mathrm{GH}\) &
			Boolean-domain \(\#\mathrm{CSP}\) &
			General-domain \(\#\mathrm{CSP}\) \\
			\midrule
			\(\{0,1\}\)
			& \cite{DyerGreenhill2000}
			& \cite{DyerGoldbergPaterson2007}
			& \cite{CreignouHermann1996}
			& \cite{Bulatov2013,DyerRicherby2013} \\
			\(\mathbb Q_{\ge0}\)
			& -- & -- & --
			& \cite{BulatovDyerGoldbergJalseniusJerrumRicherby2012} \\
			\(\mathbb R_{\ge0}\)
			& \cite{BulatovGrohe2005}
			& \cite{CaiChen2019}
			& \cite{DyerGoldbergJerrum2009}
			& \cite{CaiChenLu2016} \\
			\(\mathbb R\)
			& \cite{GoldbergGroheJerrumThurley2010}
			& --
			& \cite{BulatovDyerGoldbergJalseniusRicherby2009}
			& -- \\
			\(\mathbb C\)
			& \cite{CaiChenLu2013}
			& \cite{CaiChen2017}
			& \cite{CaiLuXia2014}
			& classification \cite{CaiChen2017};
			decidability \cite{LiuMeng2026CSPDecidability} \\
			\bottomrule
		\end{tabularx}
	\end{table}
	
	The last cell distinguishes classification from uniform recognition.  Cai and
	Chen's 2017 theorem classified every fixed algebraic-complex language, but its
	tractability side was expressed by three conditions over an unbounded generated
	family of partial sums.  In 2026, Liu and Meng gave a total exact
	algorithm deciding those conditions from the finite input language, proved
	that Block Orthogonality alone implies the other two, and also decided the
	corresponding degree-multiple counting-CSP criterion
	\cite{LiuMeng2026CSPDecidability}.  This makes the tractability criterion
	decidable from a finite, exactly encoded algebraic complex-valued language
	on any finite domain.
	
	For Holant on larger domains, classifications are available under further
	restrictions on the signatures.  A domain-three theorem treated one
	symmetric ternary signature with all unary signatures available
	\cite{CaiLuXia2013Domain3}; its journal version appeared in 2026
	\cite{CaiLuXia2026Domain3}.  Subsequent classifications covered special
	edge-colouring signatures and restricted domain-three and domain-four
	settings \cite{CaiGuoWilliams2016EdgeColoring,LiuFanCai2025}.  In 2025,
	Cai and Ihm proved a dichotomy for real-valued symmetric signatures on
	domain three with all unary signatures available, allowing an arbitrary
	finite signature set \cite{CaiIhm2025}.
	
	On the Boolean domain, the Holant framework was introduced in 2009
	\cite{CaiLuXia2009}.  Classifications first treated symmetric signatures or
	assumed access to specified unary signatures.  The latter assumptions were
	reduced from all unary signatures (\(\Holant^*\)), to four distinguished
	unaries (\(\Holant^+\)), to the two unary functions fixing a bit to zero or
	one (\(\Holant^c\)).  Later work treated signature sets containing any
	nonzero odd-arity signature (\(\Holant^{\mathrm{odd}}\)).  The symmetric
	and arbitrary-signature classifications are separated in
	\cref{tab:intro-boolean-holant-history}; the dash convention is the same as in
	\cref{tab:intro-gh-csp-history}.
	
	\begin{table}[!htbp]
		\centering
		\caption{Development of Boolean-domain Holant dichotomies on general
			graphs.}
		\label{tab:intro-boolean-holant-history}
		\scriptsize
		\setlength{\tabcolsep}{3pt}
		\renewcommand{\arraystretch}{1.18}
		\begin{tabularx}{\textwidth}{@{}l
				>{\raggedright\arraybackslash}X
				>{\raggedright\arraybackslash}X
				>{\raggedright\arraybackslash}X@{}}
			\multicolumn{4}{@{}l}{\textit{(a) Symmetric signature sets}}\\[0.2em]
			\toprule
			Weight range & \(\Holant^*\) & \(\Holant^c\) & Unrestricted \(\Holant\) \\
			\midrule
			\(\mathbb R_{\ge0}\) & -- & -- & -- \\
			\(\mathbb R\) & -- & \cite{CaiLuXia2009}
			& \cite{HuangLu2016} \\
			\(\mathbb C\) & \cite{CaiLuXia2009}
			& \cite{CaiHuangLu2012}
			& \cite{CaiGuoWilliams2016} \\
			\bottomrule
		\end{tabularx}
		
		\medskip
		\begin{tabularx}{\textwidth}{@{}l
				>{\raggedright\arraybackslash}X
				>{\raggedright\arraybackslash}X
				>{\raggedright\arraybackslash}X
				>{\raggedright\arraybackslash}X
				>{\raggedright\arraybackslash}X@{}}
			\multicolumn{6}{@{}l}{\textit{(b) Arbitrary signature sets}}\\[0.2em]
			\toprule
			Weight range & \(\Holant^*\) & \(\Holant^+\) & \(\Holant^c\) &
			\(\Holant^{\mathrm{odd}}\) & Unrestricted \(\Holant\) \\
			\midrule
			\(\mathbb R_{\ge0}\) & -- & -- & -- & --
			& \cite{LinWang2018} \\
			\(\mathbb R\) & -- & --
			& \cite{CaiLuXia2018RealHolantc}
			& \cite{CaiFuShao2020Entanglement}
			& \cite{ShaoCai2020} \\
			\(\mathbb C\)
			& \cite{CaiLuXia2011}
			& \cite{Backens2017}
			& \cite{Backens2021,Backens2025}
			& \cite{MengWangXiaZheng2025,GuanShaoShi2026}
			& This paper; with binary disequality: \cite{Liu2026DEQ2Holant} \\
			\bottomrule
		\end{tabularx}
	\end{table}
	
	Classifications without auxiliary unary signatures were obtained for
	symmetric complex-valued signatures \cite{CaiGuoWilliams2016}, arbitrary
	nonnegative signatures \cite{LinWang2018}, and arbitrary real-valued
	signatures \cite{ShaoCai2020}.  These results address restrictions on
	symmetry or weights that are absent from the present theorem.  Other
	dichotomies treat a single ternary signature on 3-regular graphs,
	the six-vertex and eight-vertex models, and weighted Eulerian orientations
	with arrow-reversal symmetry
	\cite{YangHuangFu2024,CaiFuXia2018,CaiFu2023,CaiFuShao2020ARS}.
	
	In 2025, Meng, Wang, and Xia proved an \(\FP^{\mathsf{NP}}\) versus
	\(\SharpP\)-hard dichotomy for weighted Eulerian orientations
	\cite{MengWangXia2025EO}.  Meng, Wang, Xia, and Zheng combined this result
	with tensor decomposition to obtain the same classification for complex
	Boolean Holant with a nonzero odd-arity signature
	\cite{MengWangXiaZheng2025}.  In 2026, Guan, Shao, and Shi gave a
	polynomial-time algorithm based on linear programming for the former
	\(\FP^{\mathsf{NP}}\) cases, yielding \(\FP\) versus \(\SharpP\)-hard
	dichotomies for both problems \cite{GuanShaoShi2026}.
	
	In 2026, Liu and Meng classified complex Boolean Holant with binary
	disequality \(X(a,b)=\one[a\ne b]\) available.  Their theorem allows an
	arbitrary finite algebraic complex-valued signature set and gives an
	explicit decidable tractability criterion \cite{Liu2026DEQ2Holant}.
	Their theorem has an equivalent formulation for \(K\)-Holant with binary
	equality \(I(a,b)=\one[a=b]\) available, obtained by replacing equality and
	disequality wires \cite[Theorem~1.2]{Liu2026DEQ2Holant}.  We use this
	formulation below; \(K\)-Holant is defined in \cref{subsec:holographic}.
	
	Xia also proposed a framework that uses finite groups associated with
	binary signatures to organize Boolean tensor-network classifications.
	It analyzes several branches under the assumption that all realizable
	nonzero quaternary signatures decompose \cite{Xia2026Framework}.
	
	Our proof handles the unrestricted all-even case by combining Liu and
	Meng's binary-equality dichotomy and their dichotomy for tensor-prime
	even-parity quaternary signatures with both endpoint values nonzero
	\cite{Liu2026DEQ2Holant} with a height decomposition based on Hamming
	imbalance.  Together with the odd-arity dichotomy, this gives the full
	classification in the unrestricted complex cell of
	\cref{tab:intro-boolean-holant-history}.
	
	\paragraph{Independent work.}
	We have been informed of at least the following related work, which was
	carried out independently of the present paper and brought to our attention
	before its public release.
	Yuan Huang informed us of work on complex-valued Holant problems on
	4-regular graphs, for which a manuscript is being prepared for posting
	on arXiv.
	Jiayi Zheng informed us of unpublished work on the complex-valued Holant
	dichotomy in the case where the realizable binary signatures form an odd
	dihedral group.
	Mingji Xia informed us of work concerning the twisted Klein-group case
	and the \(C_1\) case, under the condition that all quaternary signatures
	are decomposable.
	The results and proofs in the present paper were developed independently
	of these works.
	
	\subsection{Our result: the full dichotomy}
	\label{subsec:intro-main}
	
	The predicate \(\Tract(\cF)\), defined in \cref{sec:predicate}, is an
	equivalent formulation of the seven conditions in the odd-arity dichotomy
	of Meng, Wang, Xia, and Zheng \cite[Theorem~30]{MengWangXiaZheng2025}.
	Our result applies this criterion without an odd-arity assumption.
	Each condition concerns the signature set as a whole.  In each holographic
	alternative, one basis must work simultaneously for the edge and every
	member of \(\cF\).
	
	\begin{theorem}[Full dichotomy theorem]
		\label{thm:main}
		Let \(\cF\) be a finite set of algebraic complex-valued Boolean signatures,
		given in the exact algebraic representation fixed in
		\cref{subsec:complexity-conventions}.  The predicate \(\Tract(\cF)\)
		is decidable, and the following complexity classification holds:
		\[
		\begin{aligned}
			\Tract(\cF)
			&\quad\Longrightarrow\quad \Holant(\cF)\in\FP,\\
			\neg\Tract(\cF)
			&\quad\Longrightarrow\quad
			\Holant(\cF)\text{ is }\SharpP\text{-hard}
		\end{aligned}
		\]
		under polynomial-time Turing reductions.
	\end{theorem}
	
	The theorem allows arbitrary mixed arities, zero and tensor-decomposable signatures,
	and singular binary factors.  Recognition is effective for the fixed exact
	algebraic encoding, while evaluation and all oracle reductions are
	polynomial in the input-instance size.
	
	\subsection{Proof organization}
	\label{subsec:proof-sketch}
	
	In the standard \(K\)-basis, ordinary equality edges become disequality
	\(X\)-edges.  Each such edge contributes one zero and one one, so the
	\emph{half-charge}---half the difference between the numbers of ones and
	zeros---is conserved.  Cases supported only on one side of the central layer,
	or on a single layer, reduce to the known weighted Eulerian-orientation
	classification.  Before constructing the height parameter,
	\cref{sec:mountain-terminals} proves the binary, unary, EO, and equality
	reductions used in the later cases.  In the remaining
	two-sided, multilayer case, \cref{sec:factor-charge} shows that tensor-factor
	extraction and the charge argument either reach an already classified branch
	or produce a \emph{mountain}.  If \(d\) is the gcd of the nonzero
	half-charges supported by the current signature set, this is an arity-\(2d\)
	tensor with nonzero all-zero and all-one coefficients and all remaining
	support on the central weight-\(d\) layer.  This height \(d\) indexes the rest
	of the proof.
	
	The first two heights are treated separately.  At height one the mountain
	is binary.  The singular and infinite-order subcases fall into classifications
	already available.  In the remaining finite-order subcase, the symmetry split
	either exposes a nonsingular weighted equality, which can be normalized to
	\(I\) before applying the binary-equality dichotomy
	\cite[Theorem~1.2]{Liu2026DEQ2Holant}, stated here as
	\cref{thm:eq2-kholant-dichotomy}, or reduces to an odd cyclic or odd dihedral
	form.  Those two forms are resolved within the same height-one analysis,
	giving \cref{cor:height-one-complete}.
	
	At height two the mountain is an eight-vertex signature whose two endpoint
	values are nonzero.  If it does not factor nontrivially,
	\cref{thm:quaternary-kholant-dichotomy} applies; if it does factor, one obtains
	a nonsingular weighted equality and returns to the binary-equality case.
	Together the two branches give \cref{cor:eight-vertex-complete-closure}.
	
	Height three is settled by the loop and two-copy contractions in
	\cref{thm:height-three-complete}.  For all heights at least four,
	\cref{thm:height-at-least-four} uses one relative port normalization and
	a global phase recovery.  At heights at least five, the resulting
	matching-product structure gives a hard two-copy output directly.
	Height four is the exceptional eight-port case: a classification of the possible matchings
	leads to the Reed--Muller core and a Turing reduction adjoining binary
	equality.  A final normalization argument transports each common tractable
	representation or hardness reduction back to the original coordinates.
	
	The appendix records the exact finite certificates and their direct-use
	consequences for the odd-dihedral and high-charge arguments.
	
	\subsection{Related work}
	\label{subsec:related-work}
	
	Beyond dichotomies for exact counting on general graphs, related work
	studies restricted graph classes, stronger lower bounds under
	\(\#\mathrm{ETH}\), and algebraic criteria for equality of Holant values.
	
	Restricting the input graph to be planar can make additional counting
	problems tractable.  For Boolean counting CSP, Cai, Lu, and Xia classified
	real-valued symmetric constraints \cite{CaiLuXia2017PlanarCSP}.
	Guo and Williams extended the classification to symmetric constraints with
	complex weights \cite{GuoWilliams2020PlanarCSP}, and Cai and Fu removed the
	symmetry assumption \cite{CaiFu2022PlanarCSP}.  In this planar model, the
	incidence graph is embedded in the plane and the argument order follows
	the cyclic order around each constraint vertex.  Every case that is hard
	on general graphs but tractable in this planar setting admits a holographic
	reduction to matchgates.  These reductions evaluate the partition function
	through weighted perfect-matching sums computed by the FKT algorithm.
	
	Graph homomorphism problems are a subclass of counting CSP, so these
	Boolean-domain results also classify planar graph homomorphisms with
	complex symmetric weight matrices.  For larger domains, Cai and Maran
	proved a dichotomy for real symmetric \(3\times3\) matrices in 2023
	\cite{CaiMaran2023PlanarGH}.  Their 2024 preprint classified nonnegative
	real symmetric \(4\times4\) matrices of full rank
	\cite{CaiMaran2024PlanarGH}.  In a 2026 preprint, Cai, Maran, and Young
	obtained a dichotomy for nonnegative real symmetric matrices of arbitrary
	finite size whose diagonal entries are pairwise distinct
	\cite{CaiMaranYoung2026}.
	
	For planar Holant with symmetric complex-valued Boolean signatures,
	Cai, Fu, Guo, and Williams gave a classification that includes tractable
	cases beyond holographic reductions to the FKT algorithm
	\cite{CaiFuGuoWilliams2022PlanarHolant}.  In 2026, Cai, Fan, Shao, and Tang
	classified the eight-vertex model into three cases: tractable on general
	graphs, hard on general graphs but tractable on planar graphs, and hard
	even on planar graphs \cite{CaiFanShaoTang2026}.
	
	Beyond planar graphs, classifications are also known for other graph
	restrictions.
	Thilikos and Wiederrecht characterized the minor-closed graph classes on
	which perfect matchings can be counted in polynomial time
	\cite{ThilikosWiederrecht2024}.  A class is minor-closed if it is closed
	under vertex deletion, edge deletion, and edge contraction.
	Meng and Pan's ESA 2026 paper gives dichotomies for complex-weighted
	Boolean counting CSP when the incidence graph excludes a fixed clique as
	a minor, including each fixed variable-occurrence bound \(D\ge3\)
	\cite{MengPan2026Minor}.  The incidence graph records variables and
	constraints as its two types of vertices.  Their classifications show how
	the tractable cases depend on the excluded clique and the occurrence bound.
	
	Another direction strengthens the running-time lower bounds within the
	hard cases.  Under the counting exponential-time hypothesis
	(\(\#\mathrm{ETH}\)), which rules out counting satisfying assignments of
	3-CNF formulas in time subexponential in the number of variables,
	Brand, Dell, and Roth
	proved dichotomies for the Tutte polynomial and unweighted Boolean counting
	CSP \cite{BrandDellRoth2019}.  Liu extended the Boolean counting-CSP
	classification to algebraic complex weights, including bounded-occurrence
	versions \cite{Liu2023CSPExponential}.
	
	For bipartite Holant, Cai and Young proved that two pairs of finite
	signature sets give equal values on all corresponding labelled graphs
	if and only if their orbit closures under holographic transformations
	intersect \cite{CaiYoung2026}.
	
	\section{Preliminaries and the Tractable Boundary}
	\label{sec:preliminaries}
	
	\subsection{Algebraic notation}
	\label{subsec:foundational-algebra}
	
	We use standard algebraic notation.  All rings are commutative with identity.
	The fields \(\mathbb Q,\mathbb R,\mathbb C\) denote the rationals, reals,
	and complex numbers, and \(F^\times=F\setminus\{0\}\) for a field \(F\).
	We also use \(\mathbb F_3=\mathbb Z/3\mathbb Z\).
	For a finite set \(V\), \(F^V\) denotes the functions \(V\to F\).
	We write \([m]=\{1,\ldots,m\}\), with \([0]=\varnothing\), and
	\(\Mat_{m\times n}(F)\) for matrices over \(F\), abbreviated to
	\(\Mat_n(F)\) when \(m=n\).  Transposes are denoted by \(M^{\mathsf T}\);
	they do not include complex conjugation.  The group of invertible matrices
	is \(\GL_n(F)\), and \(I_n\) is its identity.
	
	\subsection{Complexity and exact-input conventions}
	\label{subsec:complexity-conventions}
	
	A \emph{counting problem} is a function from finite encoded instances to exact
	numbers.  The class \(\FP\) consists of functions computable exactly by a
	deterministic algorithm in time polynomial in the input length.  The class
	\(\SharpP\) consists of the functions that count accepting computation paths
	of a nondeterministic polynomial-time Turing machine.  A (possibly
	complex-valued) counting problem is \emph{\(\SharpP\)-hard} if every function
	in \(\SharpP\) has a polynomial-time Turing reduction to it.  We use
	\emph{tractable} to mean ``in \(\FP\)'' and \emph{hard} to mean
	``\(\SharpP\)-hard'' under this convention.
	
	An \emph{oracle algorithm} for a problem \(B\) may request exact values of
	\(B\) on instances that it constructs.  For counting problems \(A,B\),
	\(A\leT B\) means that such an algorithm computes \(A\), makes only
	polynomially many calls to \(B\), constructs instances of polynomial size,
	and otherwise uses polynomial time.  We write \(A\equivT B\) when both
	directions hold.  Known nonzero algebraic rescalings and exact algebraic
	postprocessing are allowed.
	The notation \(\FP^{\mathsf{NP}}\) permits an \(\mathsf{NP}\) decision
	oracle.  Recognition of a finite exact signature list is required to halt
	with the correct answer; it is not asserted to run in polynomial time in
	that list.  In every tractable alternative, any shared basis, sign or class
	must work simultaneously for the entire signature set.
	Throughout,
	\[
	\AlgNums:=\{\alpha\in\mathbb C:p(\alpha)=0
	\text{ for some }0\ne p\in\mathbb Q[t]\}
	\qquad\text{and}\qquad
	\Ftwo:=\mathbb Z/2\mathbb Z
	\]
	are respectively the algebraic complex numbers and the two-element field.
	We fix an \emph{exact algebraic encoding} of \(\alpha\in\AlgNums\) by its
	monic irreducible polynomial over \(\mathbb Q\) and a rectangle in
	\(\mathbb C\) whose vertices have rational real and imaginary parts,
	containing \(\alpha\) and no other complex root of that polynomial;
	the rectangle is called isolating.  A
	\emph{number field} is a finite extension \(L/\mathbb Q\), and
	\([L:\mathbb Q]\) is its dimension as a \(\mathbb Q\)-vector space.
	Equality, zero tests, conjugation, and field arithmetic are exact.  The signature set is fixed when a Holant problem is evaluated;
	running time is measured in the encoded instance size.  
	
	For later comparison, a weighted Boolean \emph{counting constraint
		satisfaction problem} (counting CSP) has Boolean variables, local constraint
	functions, and partition function equal to the sum, over all variable
	assignments, of the product of all constraint values.  The notation
	\(\#\mathrm{CSP}^q\) used below imposes the additional degree-multiple
	condition specified in \cref{subsec:holant-definitions}.
	
	\subsection{Tensor and group conventions}
	\label{subsec:algebra-group-terminology}
	
	For a nonzero vector, tensor or matrix \(v\) over \(F\), write
	\([v]=F^\times v\) for its projective class.  In particular,
	\(\PGL_n(F)=\GL_n(F)/(F^\times I_n)\).
	A projective line consists of the classes in a two-dimensional linear span.
	We use the standard terms \emph{Segre variety} for the projective classes of
	nonzero simple tensors, \emph{ruling} for a linear family with all but one
	tensor factor fixed, and \emph{secant} for a line through two distinct points.
	A pencil is a family \(q_s=A-sB\); its parameter range is specified at each use.
	
	Tensor products use the convention \(V^{\otimes0}=F\), with empty product
	\(1\).  A nonzero tensor has rank one when it is a product of one vector
	on each of its specified factors.  For two factors this is matrix rank one.
	We say \(a\) divides \(f\), with cofactor \(b\), when
	\(f=a\otimes b\) up to port permutation and nonzero scalar redistribution.
	Every tensor is taken to divide the zero tensor.  A proper tensor factor
	has positive arity strictly smaller than its parent.
	
	The notation \(\Span S\) denotes linear span.  An affine subspace is a
	translate \(a+L\) of a linear subspace; its direction is \(L\), and an
	affine plane has dimension two.  For nonempty \(S\),
	\(\operatorname{aff}(S)=a+\Span\{x-a:x\in S\}\), where \(a\in S\).
	For \(x,y\in\Ftwo^V\), the dot product \(x\cdot y=\sum_{v\in V}x_vy_v\)
	is taken in \(\Ftwo\).  The notation \(x_{\widehat{pq}}\) deletes
	coordinates \(p,q\) and keeps the inherited order.
	
	We use \(H\le G\) for a subgroup, \(H\triangleleft G\) for a normal
	subgroup, and \([G:H]\) for its index.  An involution has order two.
	Our finite-group notation is
	\[
	C_m,\qquad
	D_{2m}=\langle r,s\mid r^m=s^2=1,\ srs=r^{-1}\rangle,\qquad
	V_4=C_2\times C_2,
	\]
	where \(C_m\) is cyclic of order \(m\) and \(D_{2m}\) has order \(2m\).
	The symmetric and alternating groups are \(S_m\) and \(A_m\).
	For \(S\subseteq G\), the centralizer is
	\(C_G(S)=\{g\in G:gs=sg\text{ for every }s\in S\}\), and the center is
	\(Z(G)=C_G(G)\).  If \(\zeta\) is a primitive \(m\)-th root of unity,
	then \([\mathbb Q(\zeta):\mathbb Q]=\varphi(m)\), where \(\varphi\) is
	Euler's totient.
	
	Angle brackets denote generated subgroup in a group context, linear span
	in a vector-space context, and the tensor closure \(\angles{\cC}\) for a
	signature class, as defined in \cref{subsec:holant-definitions}.
	The Frobenius pairing is defined in \cref{subsec:reduction-conventions}.
	
	\subsection{Graphs, matchings, and Walsh coordinates}
	\label{subsec:graph-code-walsh-terminology}
	
	Graphs may have parallel edges and loops; each loop contributes two to
	\(\deg_G(v)\).  A perfect matching of a finite set \(V\) is a partition
	into pairs.  The mate of a vertex is the other endpoint of its matching
	edge.  An incidence graph joins two families according to the stated
	incidence relation, and an alternating path alternates two specified edge
	types.  We write \(K_n\) for the complete graph and \(\binom Vk\) for the
	\(k\)-element subsets of \(V\).
	
	A binary linear code is a subspace of \(\Ftwo^V\), and an affine code is
	a translate of one.  An information set consists of coordinates on which
	restriction is a bijection onto the corresponding binary coordinate space.
	For a represented binary matroid, independence means linear independence
	of the representing vectors; its basis-exchange graph joins bases differing
	by one exchanged element.
	
	For \(x\in\Ftwo^V\), put \(\wt(x)=|\{v:x_v=1\}|\).
	A Hamming layer fixes \(\wt(x)\); parity means \(\wt(x)\bmod2\).
	The operations \(\oplus\) and \(\bigoplus\) denote coordinatewise addition
	in \(\Ftwo\), while \(\bar x\) is the bitwise complement.
	A portwise bit flip by \(b\in\Ftwo^V\) sends \(f\) to
	\(f^{\oplus b}(x)=f(x\oplus b)\).  Bit-flip/port-permutation orbits use
	these flips and permutations of argument positions.  Binary words are
	ordered lexicographically with \(0<1\).
	
	For an integer \(r\ge0\) and \(f:\Ftwo^r\to\AlgNums\), viewed as the lexicographically ordered column
	array \((f(x))_{x\in\Ftwo^r}\),
	its support is \(\supp(f)=\{x:f(x)\ne0\}\).  We fix
	the unnormalized Walsh matrix and the normalized Walsh coordinates
	\begin{equation}
		P=\begin{pmatrix}1&1\\1&-1\end{pmatrix},
		\qquad
		\widetilde f=(P^{-1})^{\otimes r}f.
		\label{eq:Walsh-coordinates}
	\end{equation}
	Thus inversion is \(f=P^{\otimes r}\widetilde f\).  In coordinates,
	\[
	f(x)=\sum_{y\in\Ftwo^r}(-1)^{x\cdot y}\widetilde f(y).
	\]
	Consequently, if \(\widetilde f\) has even-parity support then
	\(f(\bar x)=f(x)\), and if it has odd-parity support then
	\(f(\bar x)=-f(x)\).  An overline on a
	complex scalar means complex conjugation.  At ambient arity \(r\), we write \(\mathbf1_{\rm even}\) and
	\(\mathbf1_{\rm odd}\) for the indicator tensors of even and odd parity;
	the arity is inferred from context.
	References below to \emph{Fourier inversion} mean precisely this inverse
	Walsh formula.
	
	Given a finite set \(V\), a perfect matching \(M\) of \(V\), and
	\(\epsilon\in\Ftwo^M\), the corresponding \emph{matching code} is
	\[
	C(M,\epsilon)=
	\{x\in\Ftwo^V:x_u\oplus x_v=\epsilon_{\{u,v\}}
	\text{ for every }\{u,v\}\in M\}.
	\]
	After choosing an orientation of every edge of \(M\), write \(\vec M\) for
	the resulting set of ordered pairs.  A nonzero \emph{matching product} on
	\(V\) is a function \(f:\Ftwo^V\to\AlgNums\) of the form
	\[
	f(x)=\lambda\prod_{(u,v)\in\vec M}b_{uv}(x_u,x_v),
	\qquad
	\lambda\in\AlgNums^\times,
	\quad b_{uv}:\Ftwo^2\to\AlgNums.
	\]
	For the all-one choice of \(\epsilon\), the set \(C(M,\epsilon)\) is
	\[
	\{x\in\Ftwo^V:x_u+x_v=1\text{ for every }\{u,v\}\in M\},
	\]
	where this sum is over the integers.  It has one independent binary choice
	per matching edge.  Any coefficient function on these choices is only a
	coordinate representation; support of this form does not by itself imply
	a matching-product factorization.
	
	\subsection{Gadgets and allowed operations}
	\label{subsec:gadget-proof-vocabulary}
	
	A gadget is a finite tensor network over the allowed signatures, with
	labelled external ports.  Its output is the tensor obtained by summing over
	all internal indices.  The native edge is equality in ordinary coordinates
	and disequality in \(K\)-coordinates.  A \emph{direct} construction uses
	one fixed finite gadget.  Connecting a binary signature at a port, joining
	two copies, and contracting two ports are specified by their wiring or
	matrix formulas when used.
	
	A slice fixes some arguments of a tensor.  Such a restriction, a linear
	combination or a Walsh-coordinate calculation does not alone give a gadget.
	Pinning requires an available point-mass unary.  Interpolation instead
	gives an oracle reduction by solving a linear system from exact evaluations.
	The convention for adjoining signatures under these reductions is stated
	in \cref{subsec:reduction-conventions}; it does not identify general
	Turing reductions with direct gadgets.
	
	A signature set is \emph{preprocessed} after removing zero and nullary
	signatures as in \cref{subsec:zero-nullary}.  Tensor-factor extraction and
	global changes of basis are used only under the hypotheses of the results
	stated below.
	
	\subsection{Signatures, tensors, and Holant instances}
	\label{subsec:holant-definitions}
	
	Recall that \(\AlgNums\) is the field of algebraic complex numbers and
	\(\Ftwo\) is the binary field.  For an integer \(r\ge0\), an arity-\(r\)
	Boolean signature is a
	function \(f:\{0,1\}^r\to\AlgNums\), identified with its lexicographically
	ordered column tensor in \((\AlgNums^2)^{\otimes r}\).  Its \emph{ports} are
	its labelled argument positions, and their order is part of the data.  We
	write \(\arity(f)=r\) and \(\supp(f)=\{x:f(x)\ne0\}\), and use
	\(\one[\mathsf P]\) for the indicator of a proposition \(\mathsf P\).
	A signature set is \emph{all-even} when each positive-arity member has even
	arity.  The terms nullary, unary, binary, ternary, quaternary, six-ary, and
	eight-ary mean arity \(0,1,2,3,4,6,8\), respectively; an \(r\)-port tensor
	has arity \(r\).
	For an ordered binary signature \(B\), we identify \(B\) with the matrix
	\((B(a,b))_{a,b\in\Ftwo}\), with port \(1\) indexing rows and port \(2\)
	indexing columns.
	All binary-vector arithmetic is over \(\Ftwo\), with \(\oplus\) used when
	bitwise addition should be explicit.  The symbols \(\zero^r\) and
	\(\one^r\) denote the all-zero and all-one words, respectively; at \(r=0\),
	both denote the unique empty word.  More generally,
	\(\zero^V,\one^V\in\Ftwo^V\) are the constant words on a finite index set
	\(V\).  The values \(f(\zero^r)\) and \(f(\one^r)\) are the two
	\emph{endpoint coefficients}.  For \(x\in\Ftwo^V\) and \(A\subseteq V\),
	\(x_{\widehat A}\) denotes the restriction of \(x\) to \(V\setminus A\), in
	the inherited order when \(V\) is ordered.
	For scalars \(a,b\), we use
	\[
	\diag(a,b)=\begin{pmatrix}a&0\\0&b\end{pmatrix}.
	\]
	
	For an integer \(r\ge0\), let \(S_r\) denote the symmetric group on
	\([r]\).  For an arity-\(r\) signature \(f\) and \(\tau\in S_r\),
	\((\tau f)(x)=f(x_{\tau^{-1}(1)},\ldots,x_{\tau^{-1}(r)})\).  For signatures
	\(f:\Ftwo^r\to\AlgNums\) and \(g:\Ftwo^s\to\AlgNums\) on disjoint ordered
	ports, and \(x\in\Ftwo^r,y\in\Ftwo^s\), define
	\((f\otimes g)(x,y)=f(x)g(y)\).
	The signature \(f\) is \emph{symmetric} when \(\tau f=f\) for every
	\(\tau\in S_r\), equivalently when \(f(x)\) depends only on \(\wt(x)\).
	For \(T\subseteq[r]\), \(\one_T\in\{0,1\}^r\) denotes the incidence
	vector of \(T\); in particular \(\one_\varnothing=\zero^r\).  We write
	\(e_j=\one_{\{j\}}\) for the \(j\)-th standard basis vector whenever the
	ambient arity is clear.
	A nonzero positive-arity signature is \emph{tensor-prime} if, up to port permutation and nonzero scalar, it cannot
	be written as \(a\otimes b\) with both \(a\) and \(b\) of positive arity.
	A nonzero positive-arity signature that is not tensor-prime is
	\emph{tensor-decomposable}.
	Prime factorization is unique up to these operations and factor order
	\cite[full version, Lemma~2.13]{CaiFuShao2020ARS}.  For any class \(\cC\) of
	signatures, its tensor closure \(\angles{\cC}\) consists of tensor products
	of members of \(\cC\), up to port permutation and nonzero scalar.  By
	convention, it also contains zero and the arity-zero empty product.
	
	For an integer \(r\ge0\) and \(a\in\{0,1\}^r\), define the
	\emph{point-mass signature}
	\(\delta_a:\{0,1\}^r\to\AlgNums\) at \(a\) by
	\[
	\delta_a(x)=\one[x=a].
	\]
	For an integer \(r\ge1\) and \(a\in\{0,1\}^r\), a
	\emph{weighted generalized equality} is a signature
	\begin{equation*}
		c_0\delta_a+c_1\delta_{\bar a},
		\qquad c_0,c_1\in\AlgNums^\times,
	\end{equation*}
	where \(\bar a\) is the bitwise complement.  It is \emph{pure} when
	\(a=\zero^r\).  Portwise bit flips make it pure, but are used only when
	actually available or inside a stated global transform.
	We abbreviate generalized equality by \(\GE\).  For
	\(p,q\in\AlgNums\), in arity two,
	\(p\delta_{00}+q\delta_{11}=\diag(p,q)\) is a \emph{weighted binary
		equality}, and \(p\delta_{01}+q\delta_{10}\) is a \emph{weighted binary
		disequality}; either is nonsingular when \(pq\ne0\).
	
	For a signature set \(\cF\), a signature grid over \(\cF\) is a finite
	multigraph \(G=(V,E)\), with loops allowed,
	together with a signature \(f_v\in\cF\) satisfying
	\(\arity(f_v)=\deg_G(v)\) and an ordering, of that length, of the incident
	edge-ends at every vertex.  An edge assignment
	\(\sigma:E\to\{0,1\}\) therefore induces an ordered word \(x_v^\sigma\)
	at \(v\).  For the resulting grid \(\Omega\),
	\begin{equation*}
		\Holant_\Omega
		=
		\sum_{\sigma:E\to\{0,1\}}
		\prod_{v\in V}
		f_v(x_v^\sigma).
	\end{equation*}
	The problem \(\Holant(\cF)\) asks for this value on an input grid over
	\(\cF\).
	Throughout, comma-separated signature or signature-set arguments mean adjoining
	their union; for example, \(\Holant(\cF,a,b)\) abbreviates
	\(\Holant(\cF\cup\{a,b\})\).
	
	We use the unary pinnings
	\[
	\DeltaZero=[1,0]^{\mathsf T},
	\qquad
	\DeltaOne=[0,1]^{\mathsf T},
	\]
	and the binary equality and disequality tensors
	\[
	I=I_2=\Eq_2=
	\begin{pmatrix}1&0\\0&1\end{pmatrix},
	\qquad
	X=\Neq_2=
	\begin{pmatrix}0&1\\1&0\end{pmatrix}.
	\]
	We also fix the sign and skew matrices
	\[
	Z=\begin{pmatrix}1&0\\0&-1\end{pmatrix},
	\qquad
	Y=ZX=\begin{pmatrix}0&1\\-1&0\end{pmatrix}.
	\]
	More generally, for \(r\ge1\), let
	\[
	\Eq_r(x_1,\ldots,x_r)=\one[x_1=\cdots=x_r].
	\]
	For signature sets \(\cE\) and \(\cF\), the bipartite problem
	\(\Holant(\cE\mid\cF)\) is defined on a bipartite graph
	\(H=(U,V,E)\): vertices in the left part \(U\) are assigned signatures
	from \(\cE\), and vertices in the right part \(V\) are assigned signatures
	from \(\cF\).
	Subdividing every ordinary edge by an \(\Eq_2\)-vertex gives
	\[
	\Holant(\cF)=\Holant(\Eq_2\mid\cF).
	\]
	
	For an integer \(q\ge1\), the counting-CSP problem
	\(\#\mathrm{CSP}^q(\cF)\) restricts every variable degree to be divisible
	by \(q\); equivalently, its variable side uses \(\Eq_{kq}\), \(k\ge1\).
	We use the convention that the variable list contains only variables that
	occur in at least one constraint.  If an input encoding explicitly lists
	\(z\) isolated Boolean variables, they may first be deleted and the final
	partition value multiplied by \(2^z\).
	The notation \(\#\mathrm{CSP}^q(\cF,\cH)\) means
	\(\#\mathrm{CSP}^q(\cF\cup\cH)\).
	
	\subsection{Holographic transformations and the
		\texorpdfstring{\(K\)}{K}-basis}
	\label{subsec:holographic}
	
	We use the polynomial-time Turing-reduction and exact-postprocessing
	conventions of \cref{subsec:complexity-conventions}.
	
	For tensors or matrices \(A\) and \(B\), write
	\[
	A\doteq B
	\]
	when
	\[
	A=\lambda B
	\qquad\text{for some }\lambda\in\AlgNums^\times.
	\]
	In every gadget or reduction identity, this scalar is known, tracked, and
	restored by exact algebraic postprocessing.
	
	For \(M\in\GL_2(\AlgNums)\) and an arity-\(r\) right-side signature \(f\),
	write
	\[
	Mf=M^{\otimes r}f,
	\qquad
	M\cF=\{Mf:f\in\cF\}.
	\]
	For an arity-\(r\) left-side signature \(g\), juxtaposition denotes the dual
	action \(gM=(M^{\mathsf T})^{\otimes r}g\).  In particular,
	\((\Eq_2)M\) has matrix \(M^{\mathsf T}M\), and the holographic identity is
	\begin{equation*}
		\Holant(\Eq_2\mid\cF)
		=
		\Holant((\Eq_2)M\mid M^{-1}\cF),
		\qquad
		(\Eq_2)M=M^{\mathsf T}M.
	\end{equation*}
	This exact identity makes the transformed edge a genuine weighted binary,
	not a componentwise scalar.
	
	Write \(Z_X(\Omega)\) for the partition function of a grid whose internal
	edges carry \(X\).  The general orthogonal similitude group of \(X\) is
	\begin{equation*}
		GO(X)
		=
		\{D\in\GL_2(\AlgNums):D^{\mathsf T}XD=\lambda X
		\text{ for some }\lambda\in\AlgNums^\times\}
	\end{equation*}
	for legal global transformations preserving the native disequality edge
	\(X\) projectively.  If \(E(\Omega)\) is the internal edge set,
	\(D\Omega\) applies \(D\) at every vertex, and
	\(D^{\mathsf T}XD=\lambda X\), then
	\begin{equation}
		Z_X(D\Omega)=\lambda^{|E(\Omega)|}Z_X(\Omega).
		\label{eq:GOX-global-scalar}
	\end{equation}
	Thus the two problems differ only by the known edge-count scalar.  The
	matrix acts globally and is not thereby a local gadget.  For \(m\ge1\), let
	\[
	\mu_m=\{z\in\AlgNums:z^m=1\}.
	\]
	We use the projective linear group \(\PGL_2(\AlgNums)\) defined in
	\cref{subsec:algebra-group-terminology}.
	Thus, for a nonsingular \(2\)-by-\(2\) matrix \(T\), the general projective
	notation \([T]\) from \cref{subsec:algebra-group-terminology} identifies
	matrices differing by a nonzero algebraic scalar.  The
	\emph{projective order} of \([T]\) is the least \(m\ge1\) for which
	\(T^m=\lambda I\) for some \(\lambda\in\AlgNums^\times\), and is infinite
	if no such \(m\) exists.
	\begin{proposition}[Finite projective binary groups]
		\label{prop:finite-pgl2-classification}
		Every finite subgroup of \(\PGL_2(\mathbb C)\) is conjugate to exactly one
		of the following types:
		\[
		C_m,\qquad D_{2m},\qquad A_4,\qquad S_4,\qquad A_5,
		\]
		where \(m\ge1\) for \(C_m\), \(m\ge2\) for \(D_{2m}\),
		\(|C_m|=m\), \(|D_{2m}|=2m\), and
		\(|A_4|=12\), \(|S_4|=24\), \(|A_5|=60\).
	\end{proposition}
	
	\begin{proof}
		This is \cite[Proposition~2.1]{Liu2026DEQ2Holant}.
	\end{proof}
	
	\begin{lemma}[Finite projective-group facts]
		\label{lem:finite-pgl2-facts}
		The following properties will be used below.
		\begin{enumerate}
			\item In
			\[
			D_{2m}=\langle r,s\mid r^m=s^2=1,\ srs=r^{-1}\rangle,
			\]
			the rotations \(\angles r\cong C_m\) form a normal index-two
			subgroup, and every element outside it is a reflection and an
			involution.  If \(m\) is odd, the rotation group is the unique
			index-two subgroup.  If \(m\) is even, the only nonidentity
			rotational involution is the central element \(r^{m/2}\).
			\item The group \(A_4\) has no index-two subgroup, and the centralizer
			of an involution is \(V_4\).
			\item The unique index-two subgroup of \(S_4\) is \(A_4\).  Every
			involution is a transposition or a double transposition; precisely
			the double transpositions belong to \(A_4\).
			\item The group \(A_5\) is simple and has no index-two subgroup.  The
			centralizer of each involution is \(V_4\).
			\item The projective centralizer of the distinguished involution is
			\[
			C_{\PGL_2(\mathbb C)}([Z])
			=
			\left\{
			\left[\begin{pmatrix}a&0\\0&b\end{pmatrix}\right]:a,b\in\mathbb C,\ ab\ne0
			\right\}
			\sqcup
			\left\{
			\left[\begin{pmatrix}0&a\\b&0\end{pmatrix}\right]:a,b\in\mathbb C,\ ab\ne0
			\right\}.
			\]
			Thus \([Z]\) is the only nonidentity diagonal involution, while
			every nonsingular anti-diagonal class is an involution.
		\end{enumerate}
	\end{lemma}
	
	\begin{proof}
		This is \cite[Lemma~2.2]{Liu2026DEQ2Holant}.
	\end{proof}
	
	The conjugacy classification above is group-theoretic only.  Its
	conjugating matrix need not lie in \(GO(X)\), so it does not by itself
	define a legal global K-Holant normalization or provide an actual binary
	gadget.  Later applications retain the distinguished classes \([X]\) and
	\([Z]\).
	
	Fix the K-coordinate matrix
	\begin{equation}
		K=\frac1{\sqrt2}
		\begin{pmatrix}1&1\\ i&-i\end{pmatrix},
		\qquad
		K^{-1}=\frac1{\sqrt2}
		\begin{pmatrix}1&-i\\1&i\end{pmatrix}.
		\label{eq:K-matrix}
	\end{equation}
	Since \(K^{\mathsf T}K=X\), putting
	\(\widehat{\cF}=K^{-1}\cF\) gives
	\begin{equation}
		\Holant(\cF)
		\equivT
		\Holant(X\mid\widehat{\cF})
		=:\KHolant(\widehat{\cF}).
		\label{eq:K-equivalence}
	\end{equation}
	For an original-coordinate signature \(f\), or an original-coordinate
	signature set \(\cF\), a hat records passage to \(K\)-coordinates:
	\[
	\widehat f=K^{-1}f,
	\qquad
	\widehat{\cF}=K^{-1}\cF.
	\]
	Once we work with a signature set already in \(K\)-coordinates, we suppress
	these hats and denote its member signatures simply by \(f\).  Unless stated
	otherwise, signatures and signature sets appearing in a K-Holant problem
	are understood to be in \(K\)-coordinates, with native edge \(X\).  For a signature \(f\) already in \(K\)-coordinates, a tilde denotes
	the Walsh coordinates fixed in \cref{eq:Walsh-coordinates}; using them in
	a calculation does not implement a global transformation.  For
	\(s\in\AlgNums\), the identities
	\begin{equation*}
		P^{-1}\!\left[\frac12
		\begin{pmatrix}1+s&1-s\\1-s&1+s\end{pmatrix}\right]P
		=\diag(1,s),
		\qquad
		P^{\mathsf T}XP=2Z
	\end{equation*}
	explain its use in the odd cyclic and dihedral calculations.
	
	\subsection{Gadgets, interpolation, and factor reductions}
	\label{subsec:reduction-conventions}
	
	\paragraph{Equivalence convention.}
	Throughout a proof, \(\mathsf H(\cG)\) denotes the relevant Holant or
	\(K\)-Holant problem, and \(\cG\) denotes the current finite signature set up
	to polynomial-time Turing equivalence.  Whenever a fixed gadget construction,
	polynomial interpolation, or a cited reduction supplies a finite set \(S\)
	such that
	\begin{equation}
		\mathsf H(\cG\cup S)\leT \mathsf H(\cG),
		\label{eq:finite-equivalent-enlargement}
	\end{equation}
	the reverse reduction is immediate by restriction.  We therefore replace
	\(\cG\) by \(\cG\cup S\), keep the notation \(\cG\), and thereafter write
	\(g\in\cG\) for \(g\in S\), without recording its construction history.
	Every such enlargement is finite.  Signatures that differ by a known nonzero
	scalar are identified, with the scalar restored in the final value.  Formal
	slices, linear combinations, Walsh-coordinate identities, and local changes
	of basis do not add signatures unless one of the stated reductions supplies
	them.  A whole-set \(GO(X)\)-normalization is instead a holographic
	equivalence and is governed by \cref{lem:normalization-ledger}.
	This is the problem-level form of the gadget and interpolation rules used by
	Liu and Meng \cite[Section~2.5]{Liu2026DEQ2Holant} and by
	\cite[Section~2.3]{MengWangXiaZheng2025}; the factor instance is supplied by
	\cref{thm:external-factor}.
	
	A \emph{direct gadget} over \(\cG\) is a finite \(X\)-edge network of
	\(\cG\)-vertices with ordered external ports.  If it realizes \(h\), then
	\begin{equation}
		\KHolant(\cG\cup\{h\})
		\leT
		\KHolant(\cG).
		\label{eq:direct-gadget-reduction}
	\end{equation}
	A fixed gadget has constant size.  A gadget realizing \(\rho h\), for
	\(\rho\in\AlgNums^\times\), implements the same Turing reduction after an instance with
	\(N\) occurrences is rescaled by \(\rho^{-N}\).
	Thus a direct gadget is one instance of the equivalence convention above.
	In arguments using local topology or port order, we specify a direct gadget
	rather than relying only on membership obtained by a Turing reduction.
	
	Let \(f\) have arity \(r\).  For an ordered pair of distinct ports
	\(p=(i,j)\), a binary kernel
	\(C:\Ftwo^2\to\AlgNums\), and a residual word \(z\) indexed by
	\([r]\setminus\{i,j\}\) in increasing port order, let
	\(z^{i\leftarrow a,j\leftarrow b}\) be the arity-\(r\) word obtained by
	inserting \(a,b\) at ports \(i,j\).  Define the two-port contraction by
	\begin{equation*}
		(\partial_p^C f)(z)
		=
		\sum_{a,b\in\Ftwo}C(a,b)
		f\bigl(z^{i\leftarrow a,j\leftarrow b}\bigr).
	\end{equation*}
	We also write
	\(f_{ij}^{ab}(z)=f(z^{i\leftarrow a,j\leftarrow b})\) for the residual
	slice.  If an ordered binary \(B\) is placed on the ordered port pair
	\(e=(u,v)\), we write \(B_e\).  The notation \(B_e\mid f\) means tensor
	divisibility on this pair in the sense of
	\cref{subsec:algebra-group-terminology}.
	In the notation of Shao--Cai
	\cite[full version, Section~2.4]{ShaoCai2020}, their
	\(\widehat\partial_{ij}\widehat f\) and \(\partial_{ij}f\) are,
	respectively, our \(\partial_{(i,j)}^X\widehat f\) and
	\(\partial_{(i,j)}^I f\).  Port order fixes the ordered pair \((i,j)\).
	For \(U,V\in\Mat_2(\AlgNums)\), the Frobenius pairing is bilinear:
	\begin{equation*}
		\langle U,V\rangle_{\mathrm{Frob}}
		=\operatorname{tr}(U^{\mathsf T}V)
		=\sum_{a,b\in\Ftwo}U(a,b)V(a,b),
	\end{equation*}
	without conjugation.  Two matrices are \emph{orthogonal} for this pairing
	when their Frobenius pairing is zero.  Suppose a directly realized
	subnetwork with ordered boundary tensor \(B\) closes ports \(p\).  Its two attachment
	\(X\)-edges complement both boundary bits, so
	\begin{equation*}
		C(a,b)
		=\sum_{u,v\in\Ftwo}X(a,u)B(u,v)X(v,b)
		=(XBX)(a,b)
		=B(1-a,1-b).
	\end{equation*}
	Exchanging the two ordered boundary inputs replaces \(B\) by
	\(B^{\mathsf T}\), and the same attachment therefore gives
	\(XB^{\mathsf T}X\).  Likewise, two inward-ordered binary chains \(L\) and
	\(R\) joined by one native \(X\)-edge have external tensor
	\(LXR^{\mathsf T}\).  These formulas explain where transposes arise in
	direct gadget constructions and fix the ordered-port convention used
	below.  Every contraction so constructed is directly realized; a contraction
	with arbitrary \(C\), or a linear combination of contractions, is only a
	tensor calculation unless a separate reduction makes it available.
	
	For any ordered binary signature \(B\), its \emph{transfer} is
	\[
	T_B=BX.
	\]
	Series composition through one native edge multiplies transfers.  Reversing
	the two ports induces
	\[
	T^\#=XT^{\mathsf T}X.
	\]
	Thus \((ST)^\#=T^\#S^\#\) and \((T^\#)^\#=T\).  For nonsingular
	\(B\), the projective order of its transfer is the order of \([BX]\).
	Global normalizations act on the whole problem, not by adjoining individual
	signatures; after applying \cref{lem:normalization-ledger}, the transformed
	finite set is again denoted by the same letter.
	
	Let \(G\le\PGL_2(\AlgNums)\) be finite, and suppose each class has
	a directly realized ordered-binary representative.  Choose one
	nonsingular ordered binary \(B_g\) for each class \(g\), normalized by
	\(R_g=B_gX\), \([R_g]=g\), and \(R_1=I\), where \(1=[I]\).  The matrix of
	every representative of class \(g\) is \(\rho B_g\) for a unique nonzero
	scalar \(\rho\).  Since \(G\) is finite, the equivalence convention lets us
	adjoin these representatives once and write \(B_g\in\cG\).  Displayed
	projective identities are read up to their known nonzero scalars; a scalar is
	shown explicitly whenever its zero/nonzero status or a coefficient ratio is
	used.
	
	Interpolation is used only through a stated polynomial-time oracle reduction;
	linear combinations are not gadgets.  The tensor-factor extraction result is
	\cref{thm:external-factor}, with its common-one-sided exception kept explicit.
	
	\subsection{Zero and nullary preprocessing}
	\label{subsec:zero-nullary}
	
	An arity-zero (or \emph{nullary}) signature is a scalar.  Let \(\cF^\circ\)
	be obtained from \(\cF\) by deleting all identically zero and nullary
	signatures.  An instance containing a zero-signature vertex has value zero.
	Otherwise, deleting its nullary vertices gives \(\Omega^\circ\), and
	\[
	\Holant_\Omega=
	\left(\prod_{\substack{c\in\cF\\\arity(c)=0}}
	c^{m_c(\Omega)}\right)\Holant_{\Omega^\circ},
	\]
	where \(m_c(\Omega)\) is the multiplicity of \(c\) and the empty residual
	instance has value \(1\).  Thus all structural arguments assume, without
	further mention, that signatures are nonzero and have positive arity; the
	evaluation algorithm restores the known nullary factors exactly.
	
	All algebraic arithmetic and recognition use the exact-input conventions of
	\cref{subsec:complexity-conventions}.
	
	\subsection{Charge and Hamming layers}
	\label{subsec:charge}
	
	All charge terminology in this subsection is in \(K\)-coordinates.  Recall
	the Hamming weight \(\wt(x)=|\{j\in[r]:x_j=1\}|\) from
	\cref{subsec:graph-code-walsh-terminology}.  For an integer \(r\ge0\), define
	the doubled-charge function \(q_r:\Ftwo^r\to\mathbb Z\) by
	\[
	q_r(x)=2\wt(x)-r.
	\]
	For an even integer \(r\), we also use the integer half-charge
	\[
	\chi_r(x)=\frac{q_r(x)}2=\wt(x)-\frac r2.
	\]
	In the all-even analysis, gcd, divisibility, height, and interval statements
	use \(\chi_r\); \emph{doubled charge} explicitly means \(q_r\).
	Put
	\[
	HW_r^{\ge}=\{x\in\Ftwo^r:q_r(x)\ge0\},
	\quad
	HW_r^{\le}=\{x\in\Ftwo^r:q_r(x)\le0\},
	\quad
	HW_r^{=}=\{x\in\Ftwo^r:q_r(x)=0\},
	\]
	and define \(HW_r^{>},HW_r^{<}\) by strict inequalities.  The unsubscripted
	symbols use the ambient arity:
	\[
	HW^{\ge},\qquad HW^{\le},\qquad HW^=,\qquad HW^>,\qquad HW^<.
	\]
	A
	signature has one of these properties when its support lies in the
	corresponding set.  A signature is \emph{single-layer} when its support lies
	in one Hamming layer; a signature set is single-layer when each member is so
	supported, with member-dependent layers allowed, and is \emph{multilayer}
	otherwise.  A signature set is \emph{common
		one-sided} when all members are \(HW^{\ge}\) or all are \(HW^{\le}\);
	otherwise it is \emph{two-sided}, equivalently both charge signs occur among
	support words of its members.  Here \(\EO\) abbreviates Eulerian orientation.  An
	arity-\(2d\) \emph{EO signature} is supported in \(HW_{2d}^=\); this central
	support condition is the local balance constraint for an Eulerian
	orientation.
	
	\begin{definition}[Mountain]
		\label{def:mountain}
		For an integer \(d\ge1\), an arity-\(2d\) signature \(M\) in
		\(K\)-coordinates is a mountain of height \(d\) if
		\[
		\supp(M)\subseteq HW_{2d}^=\cup\{\zero^{2d},\one^{2d}\}
		\qquad\text{and}\qquad
		M(\zero^{2d})M(\one^{2d})\ne0.
		\]
		Equivalently,
		\[
		\{\chi_{2d}(x):M(x)\ne0\}\subseteq\{-d,0,d\},
		\qquad
		\{-d,d\}\subseteq\{\chi_{2d}(x):M(x)\ne0\}.
		\]
		Thus the height is both half the arity and the absolute integer half-charge
		of either endpoint.  No restriction is placed on the central EO layer.
	\end{definition}
	
	For an \(X\)-edge grid \(\Omega\) with vertex set \(V(\Omega)\), every edge
	contributes one \(0\)-bit and one \(1\)-bit.  Hence each contributing
	assignment \(x\), with induced local word \(x_v\) at vertex \(v\), satisfies
	\begin{equation*}
		\sum_{v\in V(\Omega)}q_{\arity(f_v)}(x_v)=0.
	\end{equation*}
	Thus \(X\)-contraction preserves external total charge.
	
	\begin{remark}[Notation relative to prior work]
		\label{rem:mwxz-notation}
		We follow \cite[Sections~2.1--2.3]{MengWangXiaZheng2025} in writing
		\(\widehat{\cF}=K^{-1}\cF\), \(HW^{\ge}\), \(HW^{\le}\), \(HW^=\),
		and \emph{EO signature}.  Their \(\neq_2\) and \(=_r\) are our \(X\) and
		\(\Eq_r\), respectively; subscripts in \(HW_r^\bullet\) record arity.  Thus
		\(q_r(x)=|\{j\in[r]:x_j=1\}|-|\{j\in[r]:x_j=0\}|\), while the height arguments use its
		integer half \(\chi_r=q_r/2\).
		
		Our \(\Span\) is linear span; the affine span denoted
		\(\operatorname{Span}(f)\) in \cite{MengWangXiaZheng2025} is written
		\(\operatorname{aff}(\supp(f))\).
	\end{remark}
	
	\subsection{Bell signatures and the Reed--Muller core}
	\label{subsec:pauli-notation}
	
	For \(d,\ell,u,v\in\Ftwo\), define the ordered Bell signature
	\begin{equation*}
		B_{d,\ell}(u,v)
		=
		\one[u\oplus v=d](-1)^{\ell u}.
	\end{equation*}
	Write \(B_\gamma=B_{d,\ell}\) for
	\(\gamma=(d,\ell)^{\mathsf T}\).  Port order matters
	for \(B_{1,1}\), which is negated under reversal.
	
	With truth-table order \(00,01,10,11\), these four signatures are
	\begin{equation}
		\begin{aligned}
			B_{0,0}&=(1,0,0,1)=I,&
			B_{0,1}&=(1,0,0,-1)=Z,\\
			B_{1,0}&=(0,1,1,0)=X,&
			B_{1,1}&=(0,1,-1,0)=Y.
		\end{aligned}
		\label{eq:bell-matrix-dictionary}
	\end{equation}
	On an even finite port set, a \emph{Bell product} is a nonzero scalar
	multiple of a tensor product of Bell signatures on the oriented pairs of
	a perfect matching.
	Reversing a pair carrying \(B_{1,1}\) changes the scalar sign; on the empty
	port set a Bell product is a nonzero scalar.  A binary proportional to a
	Bell signature is written \(B\doteq B_\gamma\).
	Our real skew matrix \(Y\) differs from the conventional complex Pauli
	\(Y\)-matrix by a nonzero scalar.
	
	\begin{remark}[Bell notation in Shao--Cai]
		The tuple in \cref{eq:bell-matrix-dictionary} is Shao--Cai's
		\((=^+_2,=^-_2,\neq^+_2,\neq^-_2)\) \cite[Section~2]{ShaoCai2020};
		\(Y^{\mathsf T}=-Y\).  Their holographic matrix named \(Z\) is our \(K\),
		whereas our \(Z\) is only the diagonal sign matrix.
	\end{remark}
	
	Index the eight ports by the points of \(\Ftwo^3\) in lexicographic order.
	The first-order Reed--Muller code is
	\[
	RM(1,3)
	=
	\{(\ell(t))_{t\in\Ftwo^3}\mid
	\ell\colon\Ftwo^3\to\Ftwo\text{ is affine}\}.
	\]
	Explicitly, its elements are the truth tables of
	\[
	\ell(t_1,t_2,t_3)
	=a+u_1t_1+u_2t_2+u_3t_3,
	\qquad
	a,u_1,u_2,u_3\in\Ftwo,
	\]
	with all arithmetic over \(\Ftwo\).  Evaluating such an \(\ell\) on the
	eight inputs \(t\in\Ftwo^3\), in the fixed order above, gives one binary
	string of length eight.  The \(2^4=16\) choices of
	\((a,u_1,u_2,u_3)\) give exactly the sixteen codewords of \(RM(1,3)\).
	For \(x\in\Ftwo^8\), we write
	\begin{equation}
		\RMcore(x)=\one[x\in RM(1,3)].
		\label{eq:RM-core}
	\end{equation}
	
	\begin{remark}[The Shao--Cai \(f_8\) core]
		\label{rem:sc-f8-core}
		With the preceding port order, \(\RMcore\) is
		the signature \(f_8=\widehat f_8\) of
		\cite[full version, Section~8]{ShaoCai2020}.  Its Bell contraction identities
		are established in \cite[Theorems~13--14]{CaiFuShao2020Entanglement};
		\cref{thm:q8-normalized-bell-lift} proves the reduction adjoining equality
		to the current signature set.
	\end{remark}
	
	\subsection{Exact finite verification}
	\label{subsec:certificate-terminology}
	
	The appendix states the finite calculations used in the proof and describes
	how the supplied programs reconstruct and check the candidates from their
	defining equations.  All arithmetic is exact.  For an integer matrix, a
	nonzero \(k\)-by-\(k\) minor, or rank \(k\) modulo a prime, gives rational
	rank at least \(k\).  For a matrix with \(n\) columns, \(s\) independent
	rational kernel vectors give rank at most \(n-s\).  We use modular
	arithmetic only with these stated rank bounds, not as a substitute for a
	characteristic-zero identity.
	
	The six-port calculation also uses exact satisfiability checks for real
	polynomial systems after splitting complex variables into real and imaginary
	parts.  The appendix records the normalization and coordinate conventions,
	software requirements and which conclusions depend on the solver.  The correspondence between results and programs is
	in \texttt{certificates/theorem-to-script-map.md}.
	
	\subsection{Tractable signature sets and effective recognition}
	\label{sec:predicate}
	
	\subsubsection{The standard signature classes}
	\label{subsec:standard-classes}
	
	We include the zero signature in every class below.  The affine class
	\(\cA\) consists of the zero signature and the nonzero arity-\(r\) signatures
	\begin{equation}
		f(x)=\lambda\,\one[Ax=b]\,i^{Q(x)},
		\label{eq:affine-class}
	\end{equation}
	where \(r,m\ge0\) are integers,
	\(\lambda\in\AlgNums^\times\),
	\(A\in\Mat_{m\times r}(\Ftwo)\), \(b\in\Ftwo^m\), and
	\(Q:\Ftwo^r\to\mathbb Z_4\), where
	\(\mathbb Z_4=\mathbb Z/4\mathbb Z\), is represented by a multilinear
	polynomial over \(\mathbb Z_4\) of degree at most two whose mixed
	coefficients are even.
	Equivalently, the phase is a product of fourth-root phases of affine linear
	forms interpreted as \(\{0,1\}\)-valued integers.
	Here a \emph{mixed coefficient} is the coefficient of \(x_ix_j\) for
	\(i\ne j\), and an affine linear form over \(\Ftwo\) is
	\(a_0+\sum_{j=1}^r a_jx_j\) with all coefficients in \(\Ftwo\).
	
	The product class \(\cP\) consists of zero and all signatures of the form
	\(\lambda\prod_{j=1}^m h_j\), where \(m\ge0\) is an integer,
	\(0\ne\lambda\in\AlgNums\), and every factor
	\(h_j\) is a unary, binary equality, or binary disequality constraint applied
	to the signature variables.  Empty products and repeated variables are
	allowed; an inconsistent product denotes the zero signature.
	
	For every integer \(q\ge2\), put
	\begin{equation*}
		\rho_q=e^{\pi i/(2q)},\qquad T_q=\diag(1,\rho_q).
	\end{equation*}
	For \(r\in\mathbb Z\), define the arity-wise affine twist \(\cA_q^r\) by
	\begin{equation*}
		f\in\cA_q^r
		\quad\Longleftrightarrow\quad (T_q^{-r})^{\otimes\arity(f)}f\in\cA.
	\end{equation*}
	Since \(T_q^q=\diag(1,i)\) preserves \(\cA\), we take \(r\bmod q\);
	in particular, \(T_2=\diag(1,e^{\pi i/4})\).
	The local-affine class \(\cL\) consists of the signatures \(f\) such that
	\begin{equation*}
		\left(\bigotimes_{j=1}^{\arity(f)}T_2^{\alpha_j}\right)f\in\cA
		\qquad\text{for every }\alpha\in\supp(f).
	\end{equation*}
	
	Let \(\cT\) be the set of all unary and binary signatures, and let
	\begin{equation*}
		\cM=\{f:\text{\(f\) is a Boolean signature and }
		f(x)=0\text{ for every \(x\) with }\wt(x)>1\}.
	\end{equation*}
	
	\begin{definition}[Common transformability and \(K\)-presentation]
		\label{def:common-transformability}
		\label{def:common-K-presentation}
		Let \(\cC\) be closed under nonzero scalars and port permutations.  A
		finite signature set \(\cF\) is \(\cC\)-transformable if one
		\(M\in\GL_2(\AlgNums)\) satisfies
		\begin{equation*}
			M^{\mathsf T}M\in\cC,\qquad M^{-1}\cF\subseteq\cC.
		\end{equation*}
		Equivalently, for \(\widehat{\cF}=K^{-1}\cF\), one matrix
		\(U=K^{-1}M\) gives the common \(K\)-presentation
		\begin{equation*}
			U^{\mathsf T}XU\in\cC,\qquad U^{-1}\widehat{\cF}\subseteq\cC.
		\end{equation*}
		In both displays the edge is included and the quantifier order is one basis
		for the entire signature set, not one basis per signature.  We use this definition
		for \(\cC\in\{\cA,\cP,\cL\}\) and for every fixed affine twist
		\(\cA_q^r\) (in particular \(\cA_2^1\)).
	\end{definition}
	
	\subsubsection{The Eulerian-orientation predicate}
	\label{subsec:eo-predicate}
	
	For an arity-\(r\) signature \(g\), define its central restriction
	\begin{equation*}
		g|_{\EO}(x)=g(x)\one[q_r(x)=0].
	\end{equation*}
	It is zero for odd \(r\), and zero restrictions are deleted.  For a
	signature set \(\cG\), set
	\begin{equation*}
		\cG|_{\EO}=\{g|_{\EO}:g\in\cG,\ g|_{\EO}\ne0\}.
	\end{equation*}
	
	Let \(g\ne0\) have arity \(k\) and be supported on the single Hamming
	layer of weight \(d\).  Its \emph{EO completion} is the ordered tensor
	\begin{equation}
		g_{\to\EO}
		=
		\begin{cases}
			g\otimes\DeltaZero^{\otimes(2d-k)},
			& 2d\ge k,\\
			g\otimes\DeltaOne^{\otimes(k-2d)},
			& 2d<k.
		\end{cases}
		\label{eq:eo-completion}
	\end{equation}
	The original ports precede the added unaries; a zeroth power is the nullary
	unit.  The completed signature has even arity and central support, hence is
	EO\@.  For a signature set of nonzero single-layer signatures, write
	\begin{equation*}
		\cG_{\to\EO}=\{g_{\to\EO}:g\in\cG\}.
	\end{equation*}
	
	For arity-\(2d\) EO signature \(g\), an oriented perfect pairing
	\(\pi=((a_1,b_1),\ldots,(a_d,b_d))\) has unordered pairs partitioning
	\([2d]\).  For \(z\in\Ftwo^d\), let \(y\in\Ftwo^{2d}\) be the unique word
	with \(y_{a_j}=z_j\) and \(y_{b_j}=1-z_j\); the pairing induces
	\begin{equation*}
		g^\pi(z_1,\ldots,z_d)=g(y).
	\end{equation*}
	
	\begin{definition}[The predicate \(\EOTract\)]
		\label{def:eo-tract}
		For a finite EO signature set \(\cG\), \(\EOTract(\cG)\) holds if there exist a
		single sign \(\varepsilon\in\{+1,-1\}\) and a single class
		\(\cC\in\{\cA,\cP\}\), both shared by the entire signature set, such that:
		\begin{enumerate}
			\item for every \(g\in\cG\) and all
			\(\alpha,\beta,\gamma\in\supp(g)\), if
			\(\delta=\alpha\oplus\beta\oplus\gamma\), then
			\(\delta\in\supp(g)\) or
			\(\varepsilon q_{\arity(g)}(\delta)>0\);
			\item for every \(g\in\cG\) and every oriented perfect pairing \(\pi\)
			of its ports, \(g^\pi\in\cC\).
		\end{enumerate}
		We adopt the vacuous convention \(\EOTract(\varnothing)=\mathrm{true}\).
	\end{definition}
	
	This is the opposite-pair restriction of
	\cite{MengWangXiaZheng2025,GuanShaoShi2026}: substitute
	\(y_{b_j}=1-y_{a_j}\), with the converse lifted, for a pairing with
	\(d\) pairs, by
	\(\prod_{j=1}^dX(y_{a_j},y_{b_j})\).  Both \(\cA,\cP\) are closed under it.
	
	\subsubsection{The tractability predicate}
	\label{subsec:tract-predicate}
	
	For an arbitrary finite signature set \(\cF\), define
	\(\Tract(\cF):=\Tract(\cF^\circ)\), where \(\cF^\circ\)
	is defined in \cref{subsec:zero-nullary}.
	For a preprocessed signature set \(\cF\), put
	\(\widehat{\cF}=K^{-1}\cF\).  Define \(\Tract(\cF)\) to be the disjunction
	of the following seven alternatives.
	
	\begin{enumerate}[label=\textnormal{(\arabic*)}]
		\item \label{case:one-sided}
		The signature set \(\widehat{\cF}\) is common one-sided, in the explicit
		sense that
		\[
		[\,\forall\widehat f\in\widehat{\cF},\
		\supp(\widehat f)\subseteq HW_{\arity(\widehat f)}^{\ge}\,]
		\quad\text{or}\quad
		[\,\forall\widehat f\in\widehat{\cF},\
		\supp(\widehat f)\subseteq HW_{\arity(\widehat f)}^{\le}\,],
		\]
		and \(\EOTract(\widehat{\cF}|_{\EO})\) holds.
		
		\item \label{case:single-weight}
		Every \(\widehat f\in\widehat{\cF}\) is supported on one Hamming
		weight \(d_f\).  Across the signature set, some support word has negative
		charge and some (possibly other) support word has positive charge, and
		\(\EOTract(\widehat{\cF}_{\to\EO})\) holds, with the completion
		operation defined in \cref{eq:eo-completion}.
		
		\item \label{case:arity-two}
		\(\cF\subseteq\angles{\cT}\).
		
		\item \label{case:matching}
		\(\cF\subseteq\angles{K\cM}\), or
		\(\cF\subseteq\angles{KX\cM}\).
		
		\item \label{case:affine}
		\(\cF\) is \(\cA\)-transformable.
		
		\item \label{case:product}
		\(\cF\) is \(\cP\)-transformable.
		
		\item \label{case:local-affine}
		\(\cF\) is \(\cL\)-transformable.
	\end{enumerate}
	
	\begin{remark}
		For every signature \(f\), the matrix symbols in \cref{case:matching} act
		tensorwise at the arity of
		each signature.  In particular,
		\[
		\begin{aligned}
			f\in\angles{K\cM}
			&\Longleftrightarrow
			\text{ every tensor-prime factor of }K^{-1}f\text{ lies in }\cM,\\
			f\in\angles{KX\cM}
			&\Longleftrightarrow
			\text{ every tensor-prime factor of }XK^{-1}f\text{ lies in }\cM.
		\end{aligned}
		\]
	\end{remark}
	
	\begin{lemma}[Effective algebraic transformability witnesses]
		\label{lem:algebraic-transform-witness}
		Let \(E\in\Mat_2(\AlgNums)\), let \(\cF\) be a finite algebraic
		signature set, and let \(\cC\in\{\cA,\cP,\cL\}\).  Suppose there is a matrix
		\(U\in\GL_2(\mathbb C)\) such that
		\[
		U^{\mathsf T}EU\in\cC,
		\qquad
		U^{-1}\cF\subseteq\cC,
		\]
		where membership in \(\cC\) allows complex values for the free scalar and
		unary parameters in its normal forms.  Then such a witness exists in
		\(\GL_2(\AlgNums)\).  Moreover, for exact algebraic input, the existence of a
		witness is decidable and an algebraic witness can be produced.
	\end{lemma}
	
	\begin{proof}
		This is \cite[Lemma~2.15]{Liu2026DEQ2Holant}.
	\end{proof}
	
	\subsubsection{Dichotomies and tensor decomposition}
	\label{subsec:dichotomy-decomposition-tools}
	
	\begin{theorem}[Weighted Eulerian orientations
		{\cite{MengWangXia2025EO}; \cite[Theorems~1.1 and~2.17]{GuanShaoShi2026}}]
		\label{thm:external-eo}
		For every finite algebraic EO signature set \(\cG\), if \(\EOTract(\cG)\) holds,
		then \(\KHolant(\cG)\in\FP\); otherwise
		\(\KHolant(\cG)\) is \(\SharpP\)-hard.  The polynomial-time side uses the
		LP algorithm of Guan, Shao, and Shi.
	\end{theorem}
	
	\begin{theorem}[Single-weight K-Holant
		{\cite[full version, Theorem~31; conference version, Theorem~28]{MengWangXiaZheng2025}}]
		\label{thm:external-single-weight}
		Let \(\cG\) be a finite preprocessed algebraic signature set in which every
		signature is supported on one Hamming layer.  Then \(\KHolant(\cG)\) is
		\(\SharpP\)-hard unless either \(\cG\) is common one-sided with
		\(\EOTract(\cG|_{\EO})\), or both positive and negative charges occur and
		\(\EOTract(\cG_{\to\EO})\).  Both exceptions are in \(\FP\) by
		\cref{thm:external-eo}.
	\end{theorem}
	
	\begin{theorem}[Odd-arity Holant
		{\cite[full version, Theorem~33; conference version, Theorem~30]{MengWangXiaZheng2025}}]
		\label{thm:external-odd}
		Let \(\cF\) be a finite preprocessed algebraic signature set containing a
		nonzero odd-arity signature.  Then \(\Holant(\cF)\in\FP\) if
		\(\Tract(\cF)\) holds, and is \(\SharpP\)-hard otherwise.
	\end{theorem}
	
	Its first two alternatives were originally in \(\FP^{\mathsf{NP}}\); the
	stated \(\FP\) conclusion follows from
	\cite[Theorem~1.1 and Corollary~1.2]{GuanShaoShi2026}.
	
	\begin{theorem}[Tensor-factor extraction
		{\cite[full version, Theorem~32 and Remark~39; conference version,
			Theorem~29]{MengWangXiaZheng2025}}]
		\label{thm:external-factor}
		Let \(\Lambda\) be a finite signature set and let
		\(h=a\otimes b\) be a nonzero tensor product.  Unless
		\(\Lambda\cup\{h\}\) is common one-sided,
		\[
		\KHolant(\Lambda,a,b)\equivT\KHolant(\Lambda,h).
		\]
		In the common-one-sided exception, \(\KHolant(\Lambda,h)\) is in \(\FP\)
		when \(\EOTract((\Lambda\cup\{h\})|_{\EO})\) holds and is
		\(\SharpP\)-hard otherwise, by charge conservation and
		\cref{thm:external-eo}.
	\end{theorem}
	
	\begin{theorem}[Quaternary eight-vertex boundary
		{\cite[Theorem~3.1]{CaiFu2023}}]
		\label{thm:external-eight-vertex}
		Let \(M\) be a quaternary signature satisfying
		\[
		\supp(M)\subseteq\{x\in\{0,1\}^4:\wt(x)=2\}\cup\{0000,1111\},\qquad
		M(0000)M(1111)\ne0.
		\]
		Then \(\Holant(X\mid M)\) is \(\SharpP\)-hard unless there exist
		\(\cC\in\{\cA,\cP,\cL\}\) and \(U\in\GL_2(\AlgNums)\) such that
		\[
		U^{\mathsf T}XU\in\cC,\qquad U^{-1}M\in\cC.
		\]
		Every exception is in \(\FP\).  The source permits \(U\) over \(\mathbb C\);
		\cref{lem:algebraic-transform-witness} makes it algebraic.
	\end{theorem}
	
	\begin{theorem}[Even-degree counting CSP
		{\cite[full version, Theorem~4.1]{CaiLuXia2018RealHolantc}}]
		\label{thm:external-csp2}
		Let \(\cG\) be a finite algebraic Boolean signature set.  Then
		\(\#\mathrm{CSP}^2(\cG)\) is \(\SharpP\)-hard unless one common alternative
		holds:
		\[
		\cG\subseteq\cA,\qquad \cG\subseteq\cP,\qquad
		\cG\subseteq\cL,\qquad\text{or}\qquad \cG\subseteq\cA_2^1.
		\]
		Each exceptional case is computable in polynomial time.
	\end{theorem}
	
	The source assumes that both pinnings are freely available.  In the present
	formulation, \(\Eq_4\) is available on the variable side of
	\(\#\mathrm{CSP}^2\), and \cref{cor:higher-equality-unary} supplies those
	pinnings.  Any ordinary-Holant oracle instance arising in this use is
	returned to bipartite CSP form by merging adjacent even-arity equality
	vertices and subdividing every remaining \(\cG\)--\(\cG\) edge with
	\(\Eq_2\).  The class denoted \(\cA^\alpha\) there is our \(\cA_2^1\);
	the inverse-twist convention does not change the class because
	\(T_2^2=\diag(1,i)\) preserves \(\cA\).
	
	\begin{theorem}[Degree-multiple counting CSP
		{\cite[full version, Theorem~5.3; proof in
			Section~5.5]{CaiFuShao2020Entanglement}}]
		\label{thm:external-cspq}
		Let \(q\ge2\) be an integer and let \(\cG\) be a finite algebraic Boolean signature set.  If
		\(\#\mathrm{CSP}^q(X,\cG)\) is not \(\SharpP\)-hard, then either
		\[
		\cG\subseteq\cP\qquad\text{or}\qquad
		\cG\subseteq\cA_q^r\quad\text{for one }r\in\{1,\ldots,q\}.
		\]
		Both alternatives are computable in polynomial time.
	\end{theorem}
	
	\subsubsection{Decidability and polynomial-time evaluation}
	\label{subsec:predicate-algorithms}
	
	\begin{theorem}[Recognition and evaluation]
		\label{thm:predicate-algorithms}
		For every finite algebraic signature set \(\cF\), the predicate \(\Tract(\cF)\),
		with the preprocessing convention of \cref{subsec:zero-nullary}, is
		decidable.  If it holds, then \(\Holant(\cF)\in\FP\).
	\end{theorem}
	
	\begin{proof}
		This is \cite[Theorem~2.22]{Liu2026DEQ2Holant}.
	\end{proof}
	
	\subsection{Problem equivalence and normalization}
	\label{sec:retained-normalization}
	
	\subsubsection{Factor closure}
	\label{subsec:factor-saturation}
	
	For a preprocessed finite signature set \(\Lambda\), choose an exact
	tensor-prime factorization of every member and let \(\Pi(\Lambda)\) be the
	finite set of all prime factors, with their ordered port blocks.  We call
	\(\Lambda\) \emph{factor-closed} when \(\Pi(\Lambda)\subseteq\Lambda\).
	
	\begin{lemma}[Factor closure]
		\label{lem:factor-saturation}
		If \(\Lambda\) is not common one-sided, then
		\begin{equation}
			\KHolant(\Lambda)
			\equivT
			\KHolant(\Lambda\cup\Pi(\Lambda)).
			\label{eq:factor-saturation}
		\end{equation}
		Accordingly, under the equivalence convention we replace \(\Lambda\) by this
		union, keep the same notation, and assume \(\Pi(\Lambda)\subseteq\Lambda\).
		Repeated factorization terminates because every proper positive-arity factor
		has smaller arity than its parent.
	\end{lemma}
	
	\begin{proof}
		This is \cite[Lemma~2.23]{Liu2026DEQ2Holant}.
	\end{proof}
	
	For a factor-closed signature set \(\Lambda\) containing a tensor-prime
	signature of arity at least three, put
	\begin{equation*}
		\nu(\Lambda)
		=\min\{\arity(h):h\in\Lambda
		\text{ is tensor-prime and }\arity(h)\ge3\}.
	\end{equation*}
	This parameter is used only after the branch in which every prime factor has
	arity at most two has been removed.
	
	\subsubsection{Downward inheritance and coordinate transport}
	
	\begin{lemma}[Downward inheritance]
		\label{lem:downward-inheritance}
		If \(\cG\subseteq\cG'\) are preprocessed signature sets and
		\(\Tract(K\cG')\) holds, then \(\Tract(K\cG)\) holds.
	\end{lemma}
	
	\begin{proof}
		This is \cite[Lemma~2.24]{Liu2026DEQ2Holant}.
	\end{proof}
	
	\begin{lemma}[\(GO(X)\)-invariance of tractability]
		\label{lem:GOX-tract-invariance}
		Let \(\cG\) be a finite preprocessed signature set and let
		\(D\in GO(X)\).  Then
		\begin{equation*}
			\Tract(K\cG)
			\quad\Longleftrightarrow\quad
			\Tract\bigl(K(D\cG)\bigr).
		\end{equation*}
	\end{lemma}
	
	\begin{proof}
		This is \cite[Lemma~2.25]{Liu2026DEQ2Holant}.
	\end{proof}
	
	\begin{lemma}[Coordinate transport along equivalent augmentations]
		\label{lem:normalization-ledger}
		Let \(\cG_0\) be a finite preprocessed algebraic signature set, and let
		\(L_0\) be a
		number field containing its entries, their conjugates, and every algebraic
		constant fixed before the construction.  Consider a sequence
		\((A_i,\Lambda_i)\), in which every \(\Lambda_i\) is a finite preprocessed
		algebraic signature set, beginning with
		\[
		A_0=I,
		\qquad
		\Lambda_0=\cG_0,
		\]
		in which each step is one of the following.
		\begin{enumerate}
			\item An augmentation satisfying the following conditions: set
			\(A_{i+1}=A_i\), choose a finite
			\(\Lambda_{i+1}\supseteq\Lambda_i\) with
			\[
			\KHolant(\Lambda_{i+1})\leT\KHolant(\Lambda_i),
			\]
			and require every newly adjoined positive-arity signature \(h\) to
			satisfy
			\[
			h\doteq A_i^{\otimes\arity(h)}h_0
			\quad\text{for a tensor }h_0\text{ over }L_0.
			\]
			\item A global normalization: choose \(D_i\in GO(X)\), and set
			\[
			A_{i+1}=D_iA_i,
			\qquad
			\Lambda_{i+1}=D_i\Lambda_i,
			\]
			recording the known nonzero scalar contributed by each transformed
			native edge as in \cref{eq:GOX-global-scalar}.
		\end{enumerate}
		Then, at every index \(i\),
		\begin{equation}
			A_i\cG_0\subseteq\Lambda_i,
			\qquad
			A_i^{\mathsf T}XA_i=\mu_iX
			\quad(\mu_i\in\AlgNums^\times),
			\qquad
			\KHolant(\Lambda_i)\equivT\KHolant(\cG_0),
			\label{eq:normalization-ledger}
		\end{equation}
		and every positive-arity \(h\in\Lambda_i\) satisfies
		\[
		h\doteq A_i^{\otimes\arity(h)}h_0
		\quad\text{for a tensor }h_0\text{ over }L_0.
		\]
		
		At an index where every nonzero directly realizable ordered binary is
		nonsingular and has finite projective transfer order, let \(G(\Lambda_i)\)
		be the group of its projective transfer classes \([BX]\), and put
		\[
		\overline G_i=[A_i]^{-1}G(\Lambda_i)[A_i].
		\]
		Then \(\overline G_i\le\PGL_2(L_0)\).  On any consecutive sequence
		of such indices, a global normalization leaves \(\overline G_i\) unchanged,
		whereas an augmentation can only enlarge it.  Thus all groups compared in
		the finite-group argument lie in the single fixed field \(L_0\); any generic
		interpolation that leaves this field is used only when it completes the
		complexity classification of the current case.
		
		If \(\Tract(K\Lambda_i)\) holds, then \(\Tract(K\cG_0)\) holds.
		Separately, if one matrix \(V\) is a common \(\cC\) K-presentation of
		\(\Lambda_i\), where \(\cC\) is closed under nonzero scalar multiplication,
		then
		\[
		W=A_i^{-1}V
		\]
		is a common \(\cC\) K-presentation of \(\cG_0\).  In particular, this
		applies to \(\cA,\cP,\cL\).
	\end{lemma}
	
	\begin{proof}
		This is \cite[Lemma~2.26]{Liu2026DEQ2Holant}, applied inductively to the two
		displayed kinds of step.  Restriction supplies the reverse reduction for an
		augmentation.  A normalization conjugates the raw transfer group by
		\([D_i]\), which cancels after pullback by \([A_{i+1}]=[D_iA_i]\).
	\end{proof}
	
	\subsection{Binary-equality and quaternary dichotomies}
	\label{subsec:equality-quaternary-interfaces}
	
	The following dichotomies are consequences of
	the results of Liu and Meng \cite{Liu2026DEQ2Holant}.
	
	\begin{theorem}[Binary-equality \(K\)-Holant dichotomy]
		\label{thm:eq2-kholant-dichotomy}
		Let \(\Lambda\) be a finite algebraic signature set containing
		\(I=\Eq_2\).  If \(\Tract(K\Lambda)\) holds, then
		\(\KHolant(\Lambda)\in\FP\); otherwise
		\(\KHolant(\Lambda)\) is \(\SharpP\)-hard.
		The applicable side is decidable from exact algebraic input.
	\end{theorem}
	
	\begin{proof}
		\cite[Theorem~1.2]{Liu2026DEQ2Holant} is stated with a reduced four-clause
		predicate.
		Since \(I\in\Lambda\), its coordinate image contains
		\(Z=K^{\otimes2}I\in K\Lambda\); hence the binary-anchor lemma
		\cite[Lemma~2.13]{Liu2026DEQ2Holant} identifies that reduced predicate with
		\(\Tract(K\Lambda)\).  The displayed formulation follows.
	\end{proof}
	
	An eight-vertex signature has arity four and vanishes on odd Hamming weight.
	
	\begin{theorem}[Tensor-prime quaternary \(K\)-Holant dichotomy]
		\label{thm:quaternary-kholant-dichotomy}
		Let \(\Lambda\) be a finite algebraic signature set containing a
		tensor-prime eight-vertex signature \(M\) with \(M(0000)M(1111)\ne0\).
		If \(\Tract(K\Lambda)\) holds, then \(\KHolant(\Lambda)\in\FP\); otherwise
		\(\KHolant(\Lambda)\) is \(\SharpP\)-hard.  The applicable side is decidable
		from exact algebraic input.
	\end{theorem}
	
	\begin{proof}
		This is \cite[Theorem~1.3]{Liu2026DEQ2Holant}.
	\end{proof}
	
	\begin{corollary}[Eight-vertex \(K\)-Holant dichotomy]
		\label{cor:eight-vertex-complete-closure}
		Let \(\Gamma\) be a finite algebraic signature set containing an
		eight-vertex signature \(M\) with \(M(0000)M(1111)\ne0\).
		If \(\Tract(K\Gamma)\) holds, then \(\KHolant(\Gamma)\in\FP\); otherwise
		\(\KHolant(\Gamma)\) is \(\SharpP\)-hard.
	\end{corollary}
	
	\begin{proof}
		If \(\Tract(K\Gamma)\) holds,
		\cref{thm:predicate-algorithms,eq:K-equivalence} give the polynomial-time
		conclusion.  Assume \(\neg\Tract(K\Gamma)\).  If \(M\) is tensor-prime,
		\cref{thm:quaternary-kholant-dichotomy} gives hardness.
		
		Suppose instead that \(M\) is tensor-decomposable.  Its nonzero endpoint values
		exclude a \(1+3\) factorization: setting the one-port block to \(1\) and the
		other block to \(000\) would give a nonzero odd-weight entry.  Thus every
		nontrivial factorization has type \(2+2\).  In such a factorization, the
		nonzero endpoints and the odd-weight zeros force each binary factor to be
		diagonal with both diagonal entries nonzero.  Up to a port permutation and a
		nonzero scalar, one factor is therefore
		\[
		E=\diag(p,q),\qquad p,q\in\AlgNums^\times.
		\]
		Because \(M\) contains both charge signs, the common-one-sided exception in
		\cref{lem:factor-saturation} is impossible.  By the equivalence convention,
		we may therefore assume \(E\in\Gamma\).
		Choose algebraic \(a,b\ne0\) with \(a^2p=b^2q=1\), and put
		\(D=\diag(a,b)\).  Then
		\[
		D^{\mathsf T}XD=abX,
		\qquad DED^{\mathsf T}=I,
		\]
		so \(D\in GO(X)\) and \(D\Gamma\) contains \(I\).  The global
		normalization and \cref{lem:GOX-tract-invariance} preserve both complexity and
		the failure of \(\Tract\).  Hence
		\cref{thm:eq2-kholant-dichotomy} gives hardness.  Thus the tensor-prime and
		tensor-decomposable cases follow from the quaternary and binary-equality
		dichotomies, respectively.
	\end{proof}
	\section{Auxiliary Reductions}
	\label{sec:mountain-terminals}
	
	This section collects binary interpolation, reductions using one occurrence
	of an additional signature, and dichotomies for sets containing a nonzero
	unary or a higher equality.  We also prove hardness when a set satisfying
	the stated charge and nontractability assumptions contains an EO signature
	outside \(\mathcal B_{\rm ARS}\).
	
	\subsection{Binary Interpolation}
	\label{proofsubsec:central-descent}
	
	\begin{lemma}[Binary interpolation]
		\label{lem:binary-interpolation}
		Let \(\Gamma\) be a finite algebraic signature set whose supported doubled
		charges include both signs, let
		\(0\ne B\in\Gamma\) be ordered binary, and put \(A=BX\).
		\begin{enumerate}
			\item If \(B=u\otimes v\) is singular, where \(u,v\ne0\), then
			factor closure gives \(u,v\in\Gamma\).
			\item If \(A\) is nonsingular and has infinite projective order, there is
			a nonzero rank-one binary \(H\in\Gamma\).
		\end{enumerate}
	\end{lemma}
	
	\begin{proof}
		This is \cite[Lemma~3.1]{Liu2026DEQ2Holant}.
	\end{proof}
	
	\subsection{Single-occurrence Reductions and Nonzero Unaries}
	\label{subsec:single-occurrence-carriers}
	
	The single-occurrence argument used below was originally developed for an
	expanded version of \cite{MengWangXiaZheng2025}. Since that version is not
	publicly available, we include the argument here for completeness.
	
	The dichotomy for a set containing a nonzero unary will apply below;
	both point masses are needed only in the constructions that explicitly use them.
	Write
	\[
	\Delta=\DeltaZero\otimes\DeltaOne,
	\]
	with the displayed order of its two ports.  For a signature \(g\), the
	superscripts \(g^{=1}\) and \(g^{\leq1}\) mean that an oracle instance contains,
	respectively, exactly one or at most one \(g\)-vertex.  They do not denote
	tensor powers.
	
	\begin{lemma}[Reduction with at most one \(\Delta\)-vertex]
		\label{lem:universal-directed-pair}
		For every fixed finite algebraic signature set \(\Gamma\),
		\[
		\KHolant(\Gamma,\Delta^{\leq1})
		\leT
		\KHolant(\Gamma).
		\]
		No EO, parity, charge, mountain, or tractability hypothesis is required.
	\end{lemma}
	
	\begin{proof}
		This is \cite[Lemma~3.2]{Liu2026DEQ2Holant}.
	\end{proof}
	
	The next elementary lemma isolates the only additional ingredient needed to
	turn one-occurrence access into unrestricted access to a new signature.  It
	applies verbatim to ordinary Holant and to K-Holant; write \(\mathsf H\) for
	either setting, with its native wires.
	
	\begin{lemma}[Single-occurrence replication]
		\label{lem:single-occurrence-replication}
		Let \(\Gamma\) be a fixed finite algebraic signature set, and let \(g\) and
		\(h\) be nonzero signatures of arities \(r\) and \(s\), respectively.  All
		connections below use the native binary edge: \(\Eq_2\) for ordinary Holant
		and \(X\) for K-Holant.  Suppose that
		\begin{enumerate}
			\item \(\mathsf H(\Gamma,g^{=1})\leT\mathsf H(\Gamma)\);
			\item there is a fixed gadget \(C\) over \(\Gamma\) whose dangling edges are
			the three disjoint ordered blocks
			\[
			A=(a_1,\ldots,a_r),\qquad
			P=(p_1,\ldots,p_r),\qquad
			Y=(y_1,\ldots,y_s).
			\]
			Connecting one \(g\)-vertex coordinatewise to \(A\) realizes
			\[
			\rho\,g(P)h(Y)
			\]
			on \(P,Y\), for a known \(\rho\in\AlgNums^\times\).  Moreover, for
			consecutive copies \(C_i,C_{i+1}\), the block \(P_i\) can be connected
			coordinatewise to \(A_{i+1}\), with no other port identifications.
			Thus a chain of \(k\) copies, beginning with one \(g\)-vertex,
			contains exactly that one \(g\)-vertex and realizes
			\[
			\rho^k g(P_k)\prod_{i=1}^k h(Y_i),
			\]
			where the blocks \(Y_1,\ldots,Y_k\) are pairwise disjoint; and
			\item there is a fixed instance \(J\) over \(\Gamma\cup\{g\}\), containing
			exactly one \(g\)-vertex, whose exactly computable partition value
			\(Z(J)\) is nonzero.
		\end{enumerate}
		Then
		\[
		\mathsf H(\Gamma,h)\leT\mathsf H(\Gamma).
		\]
	\end{lemma}
	
	\begin{proof}
		This is \cite[Lemma~3.3]{Liu2026DEQ2Holant}.  An input with no \(h\)-vertex
		is already an instance over \(\Gamma\).  Let \(\Omega\) be an input with
		\(k\ge1\) occurrences of \(h\).  Delete those vertices, retain their incident
		edges as \(k\) ordered boundary blocks, and attach the blocks
		\(Y_1,\ldots,Y_k\) of the displayed chain coordinatewise in the corresponding
		orders.  Let \(J^\circ\) be the network obtained by deleting the
		unique \(g\)-vertex of \(J\), and join its boundary coordinatewise to
		\(P_k\) in the port order inherited from that removed vertex.  Denote the
		resulting instance by \(\Omega'\).  It contains exactly one
		\(g\)-vertex, and the two tensor identities give
		\[
		Z(\Omega')=\rho^k Z(J)Z(\Omega).
		\]
		Divide by the known nonzero scalar \(\rho^kZ(J)\) and apply assumption~1.
	\end{proof}
	
	\begin{lemma}[Dichotomy with a nonzero unary]
		\label{lem:common-unary-exit}
		Let \(\cF\) and \(\Gamma\) be finite algebraic signature sets, and let
		\(u\ne0\) be unary.  If \(u\in\cF\), then
		\(\Holant(\cF)\) is \(\SharpP\)-hard unless \(\Tract(\cF)\) holds.  Likewise,
		if \(u\in\Gamma\), then
		\(\KHolant(\Gamma)\) is \(\SharpP\)-hard unless \(\Tract(K\Gamma)\) holds.
	\end{lemma}
	
	\begin{proof}
		This is \cite[Lemma~3.4]{Liu2026DEQ2Holant}.
	\end{proof}
	
	\begin{corollary}[Higher equality supplies unaries]
		\label{cor:higher-equality-unary}
		Let \(d\ge4\), and let \(\Gamma\) be a finite algebraic signature set in
		either the ordinary-Holant or K-Holant setting.  If \(\Eq_d\in\Gamma\), then
		\(\DeltaZero,\DeltaOne\in\Gamma\).  Hence
		\cref{lem:common-unary-exit} applies.
	\end{corollary}
	
	\begin{proof}
		This is \cite[Corollary~3.6]{Liu2026DEQ2Holant}.
	\end{proof}
	
	\subsection{Hardness from an Eulerian-orientation Signature}
	\label{subsec:hc-relative-eo}
	
	For the results that invoke \textnormal{(REO1)--(REO2)}, fix a finite
	preprocessed algebraic set \(\Gamma\) of even-arity signatures in
	\(K\)-coordinates satisfying:
	\begin{enumerate}[label=\textnormal{(REO\arabic*)},leftmargin=*,labelindent=0pt]
		\item Among the values \(\chi_{\arity(f)}(x)\), with \(f\in\Gamma\)
		and \(x\in\supp(f)\), at least one is positive and at least one
		is negative;
		\item \(\neg\Tract(K\Gamma)\).
	\end{enumerate}
	All gadget connections in this subsection use the native edge \(X\).
	Arrow-reversal symmetry (ARS) means
	\(h(x)=\overline{h(\bar x)}\) for every input \(x\).
	An EO signature \(h\) has \emph{projective arrow-reversal symmetry}
	(projective ARS) if, for some \(c\in\AlgNums^\times\),
	\[
	(c^{-1}h)(x)=\overline{(c^{-1}h)(\bar x)}
	\qquad\text{for every \(x\in\{0,1\}^{\arity(h)}\)}.
	\]
	Here \(\bar x\) is bitwise complement and the overline on the coefficient is
	complex conjugation.  An EO--ARS binary is both EO and ARS without rescaling.
	Let \(\mathcal B_{\rm ARS}\) consist of the nonzero scalar multiples of
	products on arbitrary perfect matchings, with an orientation chosen for each
	matched pair, of binary EO--ARS signatures
	\begin{equation}
		(0,a,\overline a,0),\qquad a\in\AlgNums^\times.
		\label{eq:hc-rel-B}
	\end{equation}
	We write
	\(\neq_4=\delta_{0011}+\delta_{1100}\).
	Nonzero scalars and output-port permutations are suppressed according to the
	equivalence convention.
	
	\begin{lemma}[Substitution of fixed gadgets in a larger signature set]
		\label{lem:ambient-eo-lifting}
		Let \(\Lambda\subseteq\Gamma\) be finite algebraic signature sets in
		K-coordinates.  Suppose that a finite source derivation starting from
		\(\Lambda\) adjoins a finite set \(S\) using only fixed direct \(X\)-gadgets.
		Then the same derivation is valid over \(\Gamma\), and
		\[
		\KHolant(\Gamma\cup S)\leT\KHolant(\Gamma).
		\]
		If, at some stage, a source subproblem \(\KHolant(\Lambda')\), with
		\(\Lambda'\subseteq\Lambda\cup S\), is \(\SharpP\)-hard, then
		\(\KHolant(\Gamma)\) is \(\SharpP\)-hard.
	\end{lemma}
	
	\begin{proof}
		Each direct gadget is substituted independently at every occurrence of its
		output signature, leaving vertices labelled from \(\Gamma\setminus\Lambda\)
		unchanged.  Induction over the finite source derivation proves the displayed
		reduction.  For the final claim, restriction and that reduction give
		\[
		\KHolant(\Lambda')
		\leT\KHolant(\Gamma\cup S)
		\leT\KHolant(\Gamma).
		\]
		A problem-level Turing reduction quoted from a source is not, by itself, a
		direct-gadget derivation: its local signature-producing operations must first
		be identified.  Tensor-factor extraction is not included in this
		statement; whenever it is used below, it is supplied separately by
		\cref{thm:external-factor}.
	\end{proof}
	
	\begin{theorem}[Hardness from an EO signature outside \(\mathcal B_{\rm ARS}\)]
		\label{thm:hc-relative-eo-terminal}
		Assume \textnormal{(REO1)--(REO2)}.  Let \(h\in\Gamma\) be a nonzero
		positive-arity EO signature.  Suppose there are no
		\(\lambda\in\AlgNums^\times\), perfect matching \(P_h\) of its ports
		together with an orientation \((i,j)\) of each matched pair, and coefficients
		\(a_{ij}\in\AlgNums^\times\) such that
		\[
		h=\lambda
		\bigotimes_{(i,j)\in P_h}(0,a_{ij},\overline{a_{ij}},0)_{ij}.
		\]
		Then \(\KHolant(\Gamma)\) is \(\SharpP\)-hard.
	\end{theorem}
	
	\begin{lemma}[Complementary quaternaries supply a unary]
		\label{lem:hc-rel-complementary-unary}
		Assume \textnormal{(REO1)--(REO2)}, and suppose \(D\in\Gamma\), where
		\[
		D=p\delta_s+q\delta_{\bar s},
		\qquad
		s\in\{0,1\}^4,\quad \wt(s)=2,\quad p,q\in\AlgNums^\times,
		\]
		then \(\Gamma\) contains a nonzero unary.
	\end{lemma}
	
	\begin{proof}
		First reduce \(D\) to the unweighted complementary quaternary.  After a
		port permutation take \(s=0101\), and connect the same last two port
		positions of two copies through native \(X\)-edges.  The resulting ordered
		quaternary is
		\[
		\begin{aligned}
			Q(x_1,x_2,y_1,y_2)
			&=\sum_{u,v\in\Ftwo}
			D(x_1,x_2,u,v)D(y_1,y_2,1-u,1-v)\\
			&=pq\bigl(\delta_{0110}+\delta_{1001}\bigr).
		\end{aligned}
		\]
		Thus, in output order \((x_1,y_2,x_2,y_1)\), the output is the nonzero
		scalar multiple \(pq\,\neq_4\); hence \(\neq_4\in\Gamma\).
		
		Put \(E=\neq_4\).  Joining the fourth ports of
		\(E(a_1,a_2,a_3,u)\) and \(E(b_1,b_2,b_3,v)\) through \(X(u,v)\)
		realizes, in output order \((a_1,a_2,b_3,a_3,b_1,b_2)\), the six-ary
		generalized disequality
		\[
		\delta_{000111}+\delta_{111000}.
		\]
		Attach the \(\DeltaZero\)-port of one \(\Delta\)-vertex through
		\(X\) to \(a_1\), and its \(\DeltaOne\)-port through \(X\) to \(a_3\).
		This selects the word \(111000\); in remaining-port order
		\((b_1,a_2,b_2,b_3)\), the output is
		\(\delta_{0101}=\Delta\otimes\Delta\).  In the notation of
		\cref{lem:single-occurrence-replication}, delete the unique
		\(\Delta\)-vertex to obtain a gadget with input block \(A=(a_1,a_3)\)
		and output blocks \(P=(b_1,a_2)\), \(Y=(b_2,b_3)\).
		Connecting a \(\Delta\)-vertex to \(A\) realizes \(\Delta(P)\Delta(Y)\).
		Connect \(P\) of each copy to \(A\) of the next, in the displayed orders,
		with no other port identifications.  A chain of \(k\) copies therefore
		uses exactly one \(\Delta\)-vertex and realizes
		\(\Delta(P_k)\prod_{i=1}^k\Delta(Y_i)\) on pairwise disjoint output blocks.
		Closing a single \(\Delta\)-vertex by
		\(X\) has value one.  The reductions in
		\cref{lem:universal-directed-pair,lem:single-occurrence-replication}
		therefore give unrestricted access to \(\Delta\), so the equivalence
		convention places \(\Delta\in\Gamma\).
		
		Now use \(\Delta=\DeltaZero\otimes\DeltaOne\).
		By \textnormal{(REO1)}, the common-one-sided exception in
		\cref{thm:external-factor} is impossible.  Factor closure gives
		\(\DeltaZero,\DeltaOne\in\Gamma\).
	\end{proof}
	
	\begin{lemma}[EO alternative]
		\label{lem:hc-rel-mwx}
		Assume \textnormal{(REO1)--(REO2)}, and let \(0\ne h\in\Gamma\) be a
		positive-arity EO signature.  Then either there is a
		\(c\in\AlgNums^\times\) such that
		\[
		(c^{-1}h)(x)=\overline{(c^{-1}h)(\bar x)}
		\qquad(x\in\{0,1\}^{\arity(h)}),
		\]
		or \(\KHolant(\Gamma)\) is \(\SharpP\)-hard.
	\end{lemma}
	
	\begin{proof}
		We first settle the scalar field.  Suppose an algebraic signature \(h_0\) is
		ARS after multiplication by an arbitrary nonzero complex scalar.  Then, for
		some \(\rho\in\mathbb C^\times\),
		\[
		h_0(x)=\rho\,\overline{h_0(\bar x)}.
		\]
		Choose \(x\) with \(h_0(x)\ne0\).  The same identity implies
		\(h_0(\bar x)\ne0\), and hence
		\[
		\rho=\frac{h_0(x)}{\overline{h_0(\bar x)}}\in\AlgNums^\times,
		\qquad \rho\overline\rho=1.
		\]
		Set \(c=1+\rho\) if \(\rho\ne-1\), and set \(c=i\) if
		\(\rho=-1\).  Then \(c\in\AlgNums^\times\) and
		\(c/\overline c=\rho\), so \(c^{-1}h_0\) is ARS.  Thus every source
		conclusion that ARS holds up to a nonzero scalar has the algebraic scalar
		required in the statement.
		
		We next handle the binary case.  Write \(h=(0,a,b,0)\).  If \(ab=0\), then \(h\) is a
		nonzero singular binary, so \cref{lem:binary-interpolation} and factor
		extraction give a nonzero unary.  By \textnormal{(REO2)} and
		\cref{lem:common-unary-exit}, this is hard.  Suppose \(ab\ne0\), and put
		\[
		A=hX=\operatorname{diag}(a,b).
		\]
		If \(A\) has infinite projective order, \cref{lem:binary-interpolation}
		first gives a nonzero rank-one binary; applying its singular case and factor
		extraction again gives a unary and hence hardness.  Otherwise \(a/b\) is a
		root of unity, so \(|a|=|b|\).  With
		\(\rho=a/\overline b\), one has
		\(h(x)=\rho\overline{h(\bar x)}\) for all \(x\), and the preceding scalar
		argument proves the required projective ARS conclusion.
		
		Assume now that \(\arity(h)\ge4\).  Let \(L\) be the number field generated
		over \(\mathbb Q\) by all entries of \(\Gamma\) and their complex conjugates.
		It is a finite extension of \(\mathbb Q\).  Every output of a fixed direct
		\(X\)-gadget over signatures already generated from \(h\) has entries in
		\(L\), and every nonzero ratio used to normalize such an output also lies in
		\(L\).
		
		Run the generating process of
		\cite[Definition~29]{MengWangXia2025EO} on the singleton source \(\{h\}\),
		with \(X\) as its initial binary.  Each self-loop and each binary path in that
		definition is a fixed direct \(X\)-gadget.  We therefore adjoin its output to
		\(\Gamma\) by \cref{lem:ambient-eo-lifting}; zero outputs are
		discarded.  If a generated nonzero binary is singular, the singular case of
		\cref{lem:binary-interpolation} gives unary factors.  If it is nonsingular
		but its transfer has infinite projective order, the other case of that lemma
		gives a rank-one binary and then unary factors.  Either event is hard by
		\textnormal{(REO2)} and \cref{lem:common-unary-exit}.  Notice that
		\textnormal{(REO1)} persists under these equivalent augmentations, so the
		factor-extraction exception never occurs.
		
		Otherwise every generated nonzero binary is nonsingular and,
		up to a nonzero scalar, has the form
		\[
		B_r=\begin{pmatrix}0&1\\ r&0\end{pmatrix},
		\qquad r\in L^\times,
		\qquad B_rX=\operatorname{diag}(1,r),
		\]
		where \(r\) is a root of unity.  Let \(R_i\) be the parameter set represented
		by the source family \(B_i(h)\) at stage \(i\).  The path clause in
		Definition~29 makes each nonempty \(R_i\) the subgroup generated by the
		root-of-unity parameters available at that stage; repeated self-loop descent
		\cite[Lemma~39]{MengWangXia2025EO} supplies nonemptiness.  To make the
		monotonicity explicit, write the source's stage operation as \(\Phi\).
		Its self-loop and path construction is monotone under inclusion of its
		input parameter set.  Finite order and the positive-power path clause give
		\(1\in R_1\), so \(R_0=\{1\}\subseteq R_1\).  Inductively,
		\(R_i=\Phi(R_{i-1})\subseteq\Phi(R_i)=R_{i+1}\).
		Equality at one stage forces equality thereafter by the same operation.
		Every \(R_i\) lies in the group of roots of unity of \(L\), which is finite.
		Consequently this increasing chain stabilizes at some finite stage \(i_0\).
		For a unique \(k\ge1\), writing \(\mu_k=\{\omega:\omega^k=1\}\), the stabilized
		source notation is therefore
		\[
		B(h)=B_{i_0}(h)=\{B_\omega:\omega\in\mu_k\}.
		\]
		Thus every gadget construction invoked below is finite, as
		required by \cref{lem:ambient-eo-lifting}.
		
		If \(k\ge3\), \cite[Lemmas~40--42]{MengWangXia2025EO} show that \(h\) is
		ARS up to a nonzero scalar.  The opening scalar argument completes this case.
		
		Suppose \(k=2\), and write \(Y=B_{-1}=(0,1,-1,0)\).  Both \(X\) and \(Y\)
		are available in the source derivation, and \(Y\) is realized by a direct
		gadget.  We use the proof of
		\cite[Lemma~47]{MengWangXia2025EO}, not merely its displayed problem-level
		reduction.  Its arity and support-distance inductions use only fixed
		\(X\)- or \(Y\)-self-loops.  At arity four, the proofs of
		\cite[Lemmas~45--46]{MengWangXia2025EO}, before invoking their general
		problem reduction, construct either an EO quaternary
		\(q\) for which the singleton problem is hard by
		\cite[Theorem~16]{MengWangXia2025EO}, or a nonzero complementary
		quaternary
		\[
		D=p\delta_s+q'\delta_{\bar s},
		\qquad pq'\ne0;
		\]
		the two-copy subcase realizes \(\neq_4\) itself.  All these signatures arise
		by the fixed self-loops and two-copy gadgets just identified, so
		\cref{lem:ambient-eo-lifting} adjoins them relative to \(\Gamma\).  The hard
		quaternary gives hardness by restriction.  The complementary quaternary gives a
		unary by \cref{lem:hc-rel-complementary-unary}, and then hardness by
		\cref{lem:common-unary-exit}.  In the remaining case, the proof of
		Lemma~47 gives ARS up to a nonzero scalar, to which the opening argument applies.
		
		Finally suppose \(k=1\).  Follow the proof of
		\cite[Lemma~48]{MengWangXia2025EO}, and let \(C(h)\) denote its set of
		nonzero quaternary signatures obtained by successive \(X\)-self-loops.
		Every member of \(C(h)\) is therefore available by
		\cref{lem:ambient-eo-lifting}.  Moreover, every binary generated from such a
		signature is also generated from \(h\).  Hence its generated binary set is
		contained in \(B(h)=\{X\}\), and Lemma~39 makes that set nonempty; it is
		therefore exactly \(\{X\}\).  If a member of \(C(h)\) lies outside
		\(\{X^{\otimes2},Y^{\otimes2}\}\), the proof of Lemma~45 reaches exactly
		one of the hard or complementary quaternaries described in the preceding
		paragraph.  If \(Y^{\otimes2}\in C(h)\), adjoin this signature and apply
		\cref{thm:external-factor}.  Assumption \textnormal{(REO1)} excludes its exceptional case, so
		\(Y\) is available; the proof of Lemma~47 then applies because
		\(B(h)\subseteq\{X,Y\}\) and both binaries are available.  This replaces,
		rather than imports, the problem-level tensor-power reduction in
		\cite[Lemma~38]{MengWangXia2025EO}.  The case \(C(h)=\varnothing\) is
		impossible for nonzero \(h\) by Lemma~39.  In the remaining case
		\(C(h)=\{X^{\otimes2}\}\), the reorganizing argument in Lemma~48, together
		with \cite[Lemma~43]{MengWangXia2025EO}, makes \(h\) a scalar multiple of a
		real domain-symmetric signature, hence projectively ARS.  Its real and
		imaginary parts are used only as vectors in that argument; they are
		not signatures adjoined to \(\Gamma\).  The opening scalar argument again
		makes the scalar algebraic.
	\end{proof}
	
	\begin{lemma}[ARS alternative]
		\label{lem:hc-rel-cfs}
		Assume \textnormal{(REO1)--(REO2)}, and let \(0\ne h\in\Gamma\) be a
		positive-arity EO signature for which there is a
		\(c\in\AlgNums^\times\) satisfying
		\[
		(c^{-1}h)(x)=\overline{(c^{-1}h)(\bar x)}
		\qquad(x\in\{0,1\}^{\arity(h)}).
		\]
		If there are no \(\lambda\in\AlgNums^\times\), perfect matching \(P_h\)
		together with an orientation \((i,j)\) of each matched pair, and
		\(a_{ij}\in\AlgNums^\times\) such that
		\[
		h=\lambda
		\bigotimes_{(i,j)\in P_h}(0,a_{ij},\overline{a_{ij}},0)_{ij},
		\]
		then \(\KHolant(\Gamma)\) is \(\SharpP\)-hard.
	\end{lemma}
	
	\begin{proof}
		Put \(h'=c^{-1}h\), which is algebraic, EO, and ARS.  By unique factorization
		and the ARS factorization lemma
		\cite[Lemmas~2.13--2.14]{CaiFuShao2020ARS}, choose a tensor-prime
		factorization of \(h'\), up to an overall scalar, into factors
		\(g_1,\ldots,g_m\) on disjoint port blocks, each satisfying ARS exactly.
		We now choose representatives without losing either
		algebraicity or ARS.  For each \(j\), choose
		\(\alpha_j\in\operatorname{supp}(g_j)\), let
		\(\alpha=\alpha_1\cdots\alpha_m\), and define on the port block of \(g_j\)
		\[
		q_j(z)=
		\frac{h'(\alpha_1,\ldots,\alpha_{j-1},z,
			\alpha_{j+1},\ldots,\alpha_m)}{h'(\alpha)}.
		\]
		Then \(q_j=g_j/g_j(\alpha_j)\).  Hence every \(q_j\) is algebraic, EO, tensor-prime,
		and projectively ARS, and
		\[
		h'=h'(\alpha)\bigotimes_{j=1}^m q_j.
		\]
		For each \(q_j\), the algebraic-scalar argument in the preceding lemma gives
		\(c_j\in\AlgNums^\times\) such that \(p_j=c_j^{-1}q_j\) satisfies ARS
		exactly.  Thus
		\[
		h'=h'(\alpha)\Bigl(\prod_{j=1}^m c_j\Bigr)
		\bigotimes_{j=1}^m p_j,
		\]
		with an algebraic overall scalar and algebraic EO--ARS tensor-prime factors \(p_j\).
		If every tensor-prime factor were binary, the prohibited displayed decomposition of
		\(h\) would hold.  Thus there is a nonbinary tensor-prime EO--ARS factor
		\(f\notin\mathcal B_{\rm ARS}\).  Iterated use of
		\cref{thm:external-factor} adjoins \(f\) to the current set; by
		\textnormal{(REO1)}, its common-one-sided exception is impossible.  Whenever
		the cited descent subsequently extracts another factor, we repeat the
		displayed ratio construction and rescale by a nonzero algebraic scalar to
		obtain an ARS representative.
		
		Apply the proof of
		\cite[Lemmas~3.2 and~4.1--4.7; Theorem~4.8]{CaiFuShao2020ARS} to the
		singleton EO--ARS source generated from \(f\), not to the possibly non-EO
		ambient language \(\Gamma\).  The signature-producing operations in that
		proof are local: the tensor connections in Lemma~3.2 and the mergings in Lemma~4.6 are
		fixed direct \(X\)-gadgets, while the selected-port attachments of the fixed
		binary \(iY=(0,i,-i,0)\) in Lemmas~4.5--4.6 are direct gadgets after that
		binary factor has been extracted.  These direct gadgets can be substituted over
		\(\Gamma\) by
		\cref{lem:ambient-eo-lifting}; every factor-extraction step is supplied
		separately by \cref{thm:external-factor}, whose exception is again excluded
		by \textnormal{(REO1)}.  Sections~3--4 of the cited proof use no polynomial
		interpolation.
		
		Concretely, Lemma~3.2 produces a direct EO quaternary whose singleton problem
		is hard, directly realizes \(\neq_4\), or enters its mutual-orthogonality
		case.  Lemmas~4.1--4.4 supply the divisibility propagation used in
		Lemma~4.6.  Lemmas~4.5--4.6 descend through signatures outside
		\(\mathcal B_{\rm ARS}\), and Lemma~4.7 rules out mutual orthogonality in the
		base arities \(4,6,8\).  Hence the descent adjoins either a hard quaternary
		or \(\neq_4\) by the reductions just described.  The first gives hardness by restriction;
		the second gives a unary by \cref{lem:hc-rel-complementary-unary} and then
		hardness by \cref{lem:common-unary-exit}.
		
		Finally, the coordinate restrictions called \emph{pinnings} in the proof of
		\cite[Lemma~2.14]{CaiFuShao2020ARS} are coordinate slices used to establish
		tensor-divisibility identities; they are not signatures obtained from direct
		pinning gadgets.  The induction above starts from \(h\in\Gamma\),
		or a factor adjoined by \cref{thm:external-factor}; a coordinate restriction of a
		larger tensor would not suffice, since connecting copies of the larger tensor
		could introduce additional terms.
	\end{proof}
	
	\begin{proof}[Proof of \cref{thm:hc-relative-eo-terminal}]
		Apply \cref{lem:hc-rel-mwx}.  Either \(\KHolant(\Gamma)\) is
		\(\SharpP\)-hard, or \(h\) is projectively ARS.  In the latter case,
		\(h\notin\mathcal B_{\rm ARS}\), so \cref{lem:hc-rel-cfs} gives hardness.
	\end{proof}
	
	\begin{remark}[Scope of the literature inputs]
		The ARS step uses \cite[Theorem~4.8]{CaiFuShao2020ARS}, not the Section~5
		theorem for a wholly EO--ARS signature set.  No theorem requiring the ambient
		set itself to be EO or EO--ARS is invoked: the source results are applied only
		to the indicated EO sublanguages, their direct gadgets are substituted using
		\cref{lem:ambient-eo-lifting}, and factor extraction is handled separately by
		\cref{thm:external-factor}.
	\end{remark}
	
	
	\Needspace{8\baselineskip}
	\subsection{Equalities and Generalized Equalities}
	\label{proofsubsec:equality-anchor}
	
	\begin{remark}[Simultaneous transformation of the edge and signature]
		\label{def:equality-anchor}
		The same matrix \(U\) in \cref{lem:equality-anchor} must transform both
		the native edge to a nonzero multiple of \(\Eq_2\) and the quaternary
		signature to a nonzero multiple of \(\Eq_4\).  Transforming the quaternary
		alone does not ensure that the transformed problem has equality edges.
	\end{remark}
	
	\begin{lemma}[Dichotomy from a transformed quaternary equality]
		\label{lem:equality-anchor}
		Let \(\cG\) be a finite algebraic signature set containing an ordered
		quaternary signature \(g\).  Suppose there are
		\[
		U\in\GL_2(\AlgNums),
		\qquad \kappa,\lambda\in\AlgNums^\times,
		\]
		such that
		\[
		(U^{-1})^{\mathsf T}XU^{-1}=\kappa\Eq_2,
		\qquad
		U^{\otimes4}g=\lambda\Eq_4.
		\]
		Then either \(\KHolant(\cG)\) is \(\SharpP\)-hard or
		\(\Tract(K\cG)\) holds.
	\end{lemma}
	
	\begin{proof}
		This is \cite[Lemma~3.8]{Liu2026DEQ2Holant}.
	\end{proof}
	
	\begin{lemma}[Dichotomy with a generalized equality]
		\label{lem:pure-ge}
		Let \(\cG\) be a finite algebraic signature set containing
		\[
		G_q=c_0\delta_{\zero^q}+c_1\delta_{\one^q},
		\qquad c_0,c_1\in\AlgNums^\times,\qquad q\ge4\text{ an integer}.
		\]
		Then either \(\KHolant(\cG)\) is \(\SharpP\)-hard or
		\(\Tract(K\cG)\) holds.
	\end{lemma}
	
	\begin{proof}
		This is \cite[Lemma~3.9]{Liu2026DEQ2Holant}.
	\end{proof}
	
	\begin{lemma}[Dichotomy from a two-point Walsh factor]
		\label{lem:walsh-ge-terminal}
		Let \(\Gamma\) be a finite algebraic signature set whose supported doubled
		charges include both signs.  Suppose a finite equivalent enlargement of
		\(\Gamma\) contains a nonzero signature with a tensor factor \(h\) of even
		arity \(2k\ge4\) with
		\[
		\widetilde h=(P^{-1})^{\otimes 2k}h
		=c_0\delta_a+c_1\delta_{\bar a},
		\qquad a\in\Ftwo^{2k},\quad c_0c_1\ne0.
		\]
		Then \(\KHolant(\Gamma)\) is \(\SharpP\)-hard or \(\Tract(K\Gamma)\)
		holds.  If \(c_1/c_0\) is a root of unity, a finite equivalent enlargement
		contains a quaternary \(g\) satisfying the two transformation conditions of
		\cref{lem:equality-anchor} with \(U=K\).  No reflection or prescribed group
		of binary transfers is required.
	\end{lemma}
	
	\begin{proof}
		Adjoin the specified signature and then its factor \(h\) using
		\cref{eq:finite-equivalent-enlargement,thm:external-factor}; the supported
		charges of both signs exclude the common-one-sided exception.  All subsequent gadgets are
		native \(K\)-gadgets.  In Walsh coordinates an edge is
		\(P^{\mathsf T}XP=2Z\).
		
		Choose two positions at which \(a\) has the same bit.  Interchanging the
		names of \(a,\bar a\), if necessary, makes that bit zero; this merely
		renames the two terms, and uses no bit-flip gadget.  Take two copies of
		\(h\), leave these two positions exposed in each copy, and join all other
		corresponding positions between the copies.  Write \(J\) for the joined
		positions.  The diagonal edges allow only the two equally oriented pairs
		of support words.  Since \(|J|=2k-2\) is even, their signs agree.  Thus the
		four-port gadget output \(q\) has
		\[
		\widetilde q=\lambda(\delta_{0000}+\rho\delta_{1111}),
		\quad \rho=(c_1/c_0)^2,
		\quad \lambda=2^{2k-2}(-1)^{\wt(a|_J)}c_0^2\ne0.
		\]
		Closing two ports of \(q\) gives a binary \(B\) with
		\[
		\widetilde B=2\lambda\diag(1,-\rho),
		\qquad P^{-1}(BX)P=4\lambda\diag(1,\rho).
		\]
		If \(\rho\) has infinite multiplicative order,
		\cref{lem:binary-interpolation} supplies a rank-one binary.  Factor
		extraction supplies a nonzero unary, and \cref{lem:common-unary-exit}
		gives the conclusion.
		
		Otherwise let \(\ell\) be the order of \(\rho\).  A chain of \(\ell-1\)
		copies of \(B\) has projective Walsh transfer \(\diag(1,\rho^{-1})\);
		when \(\ell=1\) use the empty chain.  Attaching this chain to one port of
		\(q\) gives a gadget output \(g\) with
		\(\widetilde g=c(\delta_{0000}+\delta_{1111})\), \(c\ne0\).
		Writing \(E_4=\mathbf1_{\rm even}\), we have
		\[
		g=2cE_4,\qquad K^{\otimes4}g=4c\Eq_4,
		\qquad (K^{-1})^{\mathsf T}XK^{-1}=I.
		\]
		These are the conditions of \cref{lem:equality-anchor} with \(U=K\), so
		that lemma applies.  The equivalent enlargements transfer
		hardness back to \(\Gamma\); tractability is inherited by its subset.
	\end{proof}
	
	\section{Obtaining a Mountain}
	\label{sec:factor-charge}
	
	For an even-arity signature set with supported half-charges of both signs,
	we can adjoin either both point masses or a mountain, as defined in
	\cref{def:mountain}, whose height is the gcd of the nonzero supported
	half-charges.
	
	\begin{lemma}[Producing the gcd-height signature or both point masses]
		\label{lem:charge-bridge}
		Let \(\cG\) be a finite preprocessed algebraic signature set of even-arity
		signatures.
		Assume that among the supported values
		\[
		\wt(x)-\frac{\arity(f)}2,
		\qquad f\in\cG,\quad f(x)\ne0,
		\]
		there is at least one positive value and at least one negative value.  Put
		\[
		d=\gcd\left\{
		\left|\wt(x)-\frac{\arity(f)}2\right|:
		f\in\cG,\ f(x)\ne0,\
		\wt(x)\ne\frac{\arity(f)}2
		\right\}.
		\]
		Then \(d\ge1\), every supported difference
		\(\wt(x)-\arity(f)/2\) lies in \(d\mathbb Z\), and at least one of the
		following reductions holds.
		\begin{enumerate}
			\item There is an arity-\(2d\) signature \(M\) such that
			\[
			\KHolant(\cG,M)\leT\KHolant(\cG),
			\]
			\[
			M(\zero^{2d})M(\one^{2d})\ne0,
			\qquad
			M(x)=0\quad\text{if }\wt(x)\notin\{0,d,2d\}.
			\]
			\item Both point masses are accessible:
			\[
			\KHolant(\cG,\DeltaZero,\DeltaOne)
			\leT\KHolant(\cG).
			\]
		\end{enumerate}
		In the second outcome, if \(\neg\Tract(K\cG)\), then
		\cref{lem:downward-inheritance,lem:common-unary-exit} and the displayed
		reduction imply that \(\KHolant(\cG)\) is \(\SharpP\)-hard.
	\end{lemma}
	
	\begin{proof}
		Let
		\[
		C=\{\chi_{\arity(f)}(x):
		f\in\cG,\ f(x)\ne0,\ \chi_{\arity(f)}(x)\ne0\}
		\]
		be the nonzero supported half-charges.  The gcd assertion is immediate.
		The next argument expresses \(d\) and \(-d\) as sums of elements of
		\(C\) with nonnegative multiplicities.  Tensor products add half-charges, but their factor
		multiplicities cannot be negative.  Because \(C\) contains both signs, take every pair
		\(a,-b\in C\) with \(a,b>0\).  If \(g=\gcd(a,b)\), assigning multiplicity
		\(b/g\) to \(a\) and \(a/g\) to \(-b\) gives total half-charge zero.  Summing
		these assignments over all positive--negative pairs produces a vector
		\(v\in\mathbb Z_{>0}^C\) satisfying
		\[
		\sum_{c\in C}v_c c=0.
		\]
		Bézout's identity, iterated over the finite set \(C\), supplies an integer
		vector \(u\in\mathbb Z^C\) with \(\sum_{c\in C}u_c c=d\)
		\cite[Section~1.3, Theorem~1.2]{Apostol1976}.  The entries of \(u\) need not be
		nonnegative, but every entry of \(v\) is positive.  Consequently, for all
		sufficiently large \(k\),
		\[
		n^+=u+kv,\qquad n^-=-u+kv
		\]
		are nonnegative integer vectors, while their total half-charges remain \(d\)
		and \(-d\), respectively.  Thus \(n^+\) and \(n^-\) are the required
		multiplicity lists.
		
		Choose a supported word of half-charge \(c\) for each \(c\in C\).  The two
		tensor products prescribed by \(n^+,n^-\) contain supported words of half-charges
		\(d,-d\), and all their supported half-charges remain divisible by \(d\).
		Applying the
		endpoint-preserving self-loop lemma
		\cite[full version, Lemma~40; conference version, Lemma~37]{MengWangXiaZheng2025}
		and the equivalence convention gives \(P,N\in\cG\) of arity \(2d\) with
		\(P(\one^{2d})N(\zero^{2d})\ne0\).  Self-loops preserve half-charge, so
		\[
		\supp(P)\cup\supp(N)
		\subseteq
		HW_{2d}^=\cup\{\zero^{2d},\one^{2d}\}.
		\]
		If either \(P\) or \(N\) has both endpoints, take that signature as the
		required \(M\).
		We may therefore assume
		\[
		P(\zero^{2d})=0,
		\qquad
		N(\one^{2d})=0;
		\]
		and their remaining entries are central.
		
		If both central parts are nonzero, permute chosen central words to
		\((\zero^d,\one^d)\) in \(P\) and \((\one^d,\zero^d)\) in \(N\), expose
		the first \(d\) ports of each signature, and join their last \(d\) ports
		coordinatewise.  The resulting arity-\(2d\) signature \(Q\) has unique
		endpoint extensions and
		\[
		Q(\zero^{2d})=P(\zero^d,\one^d)N(\zero^{2d}),\qquad
		Q(\one^{2d})=P(\one^{2d})N(\one^d,\zero^d),
		\]
		both nonzero.  If only \(P\) has a central part, write
		\(N=n\delta_{\zero^{2d}}\) for \(n\in\AlgNums^\times\), permute a supported central word of \(P\) to
		\((\zero^d,\one^d)\), and join the two \(d\)-port blocks of \(N\) to the
		last \(d\) ports of two copies of \(P\), leaving the first \(d\) ports of
		both copies exposed.  Again the endpoint extensions are unique, with
		values \(nP(\zero^d,\one^d)^2\) and \(nP(\one^{2d})^2\).  (The opposite
		one-central case is symmetric.)  Half-charge divisibility restricts either
		output to half-charges \(-d,0,d\); take that output as \(M\).
		
		If both central parts vanish, then
		\(P=p\DeltaOne^{\otimes2d}\) and
		\(N=n\DeltaZero^{\otimes2d}\) for
		\(p,n\in\AlgNums^\times\).  Since the hypothesis supplies both positive
		and negative supported differences, the common-one-sided exception to the
		tensor-factor theorem cannot occur.
		Factor closure therefore gives
		\(\DeltaZero,\DeltaOne\in\cG\), proving the unary outcome.  It is covered by
		\cref{thm:external-factor,thm:external-odd}.
	\end{proof}
	
	\begin{remark}
		The factor alternative is necessary.  For
		\(\{\operatorname{EXACT}^4_3,\operatorname{EXACT}^4_1\}\), where
		\(\operatorname{EXACT}^4_j(x)=\one[\wt(x)=j]\), conservation of integer
		half-charge places every fixed gadget on a single integer-half-charge layer.
		Thus no gadget produces a mountain, and \cref{lem:charge-bridge} gives
		both point masses by factor extraction instead.
	\end{remark}
	\section{Mountain Height One}
	\label{sec:binary-dynamics}
	
	A height-one mountain is an ordered binary signature whose two diagonal
	entries are nonzero.  We prove the dichotomy for any finite algebraic
	signature set containing one.  We use the equivalent-signature-set
	convention throughout: a signature supplied by a gadget, interpolation, or
	an applicable factor reduction is adjoined to the current set.  When direct
	realization matters, we specify the wiring, port order, and nonzero scalar.
	
	A singular binary, or a nonsingular binary whose transfer has infinite
	projective order, is handled by \cref{lem:binary-interpolation} and the
	odd-arity theorem.  Otherwise \cref{lem:bounded-torsion} makes the group of
	all directly realizable binary transfers finite, and
	\cref{thm:finite-group-split} gives a nonsingular weighted equality or an odd
	cyclic or odd dihedral group.  A weighted equality can be normalized to
	\(I=\Eq_2\) by one global diagonal transformation in \(GO(X)\), so
	\cref{thm:eq2-kholant-dichotomy} applies.  We treat the odd-dihedral case
	first, since its contraction arguments are also used in the odd-cyclic case.
	
	The distinguished binary remains the transported original mountain.
	A quaternary mountain satisfies \cref{cor:eight-vertex-complete-closure}.
	Generalized equalities and the signatures satisfying
	\cref{lem:equality-anchor} are handled by
	\cref{lem:pure-ge,lem:equality-anchor}; other tensor factors are adjoined
	using the factor-reduction equivalence.
	
	\begin{theorem}[Dichotomy with a height-one mountain]
		\label{cor:height-one-complete}
		Let \(\Gamma\) be a finite algebraic signature set containing a height-one
		mountain.  Then
		\[
		\Tract(K\Gamma)
		\qquad\text{or}\qquad
		\KHolant(\Gamma)\text{ is }\SharpP\text{-hard}.
		\]
	\end{theorem}
	
	The remainder proves the theorem using the current-set convention.  Only
	the global changes of basis and the changes in \(G(\Gamma)\) remain
	explicit.
	
	For an ordered binary signature \(B\), recall its transfer
	\[
	T_B=BX.
	\]
	Joining \(B\) and \(C\) in series through one \(X\)-edge realizes \(BXC\),
	and therefore
	\[
	T_{BXC}=T_BT_C.
	\]
	Reversing the ports gives the map defined in
	\cref{subsec:reduction-conventions}:
	\begin{equation}
		T^\#:=XT^{\mathsf T}X;
		\qquad
		T^\#=(\det T)ZT^{-1}Z\quad\text{if }T\text{ is nonsingular}.
		\label{eq:transfer-reversal}
	\end{equation}
	
	\begin{remark}[Relation to binary-signature classifications]
		The trichotomy above is the K-Holant form of the binary classification
		used by Cai--Lu--Xia \cite{CaiLuXia2009,CaiLuXia2018RealHolantc}:
		singular matrices, and nonsingular matrices of infinite or finite projective
		order.  The four ordered Bell
		signatures \(\mathcal B=\{B_{0,0},B_{0,1},B_{1,0},B_{1,1}\}\) represent
		the projective Bell classes of Shao--Cai \cite{ShaoCai2020}.  Their
		holographic matrix, denoted \(Z\) in that paper, is our \(K\); throughout
		this paper,
		\(Z=\left(\begin{smallmatrix}1&0\\0&-1\end{smallmatrix}\right)\), as in
		\cref{eq:transfer-reversal}.  Beyond
		the binary classification, we need the complete gadget transfer group
		and uniform control of its higher-arity contractions.
	\end{remark}
	
	\begin{lemma}[Bounded torsion]
		\label{lem:bounded-torsion}
		Fix a number field \(L\), and let \(G\) be the set of projective
		transfers of all directly realizable nonzero nonsingular ordered
		binary gadgets over a fixed signature set with entries in \(L\).
		If \(G\ne\varnothing\) and every element of \(G\) has finite
		projective order, then \(G\) is a finite subgroup of
		\(\PGL_2(L)\).
	\end{lemma}
	
	\begin{proof}
		This is \cite[Lemma~5.2]{Liu2026DEQ2Holant}.
	\end{proof}
	
	For a current signature set \(\cG\), suppose that every nonzero directly
	realizable ordered binary \(B\) is nonsingular and that \([BX]\) has finite
	order.  Series composition and \cref{lem:bounded-torsion} then give the finite
	group
	\begin{equation}
		G(\cG)
		=\{[BX]:B\ne0\text{ is an ordered binary directly realizable over }\cG\}.
		\label{eq:complete-safe-group}
	\end{equation}
	This definition includes all such binary gadgets, including those not yet
	adjoined to \(\cG\).  We use \(G(\cG)\) only under the stated
	nonsingularity and finite-order hypothesis; a singular binary has no
	projective transfer class.  The decision procedure in
	\cref{thm:predicate-algorithms} does not enumerate this group.
	
	\begin{theorem}[Weighted equality or an odd cyclic or dihedral group]
		\label{thm:finite-group-split}
		Let \(\cG\) be a finite algebraic signature set containing the height-one
		mountain
		\[
		B=\begin{pmatrix}a&b\\c&d\end{pmatrix},
		\qquad ad\ne0,
		\]
		and suppose every nonzero directly realizable binary over \(\cG\) is
		nonsingular and its transfer has finite projective order.  Then one of the following
		holds.
		\begin{enumerate}
			\item A nonsingular weighted equality
			\(E=\diag(p,q)\), with \(pq\ne0\), belongs to the current signature
			set.
			\item There exist an odd \(N\ge3\) and \(D\in GO(X)\) such that the legal
			global normalization \(\cG\mapsto D\cG\) puts
			\(G(D\cG)\) in exactly one of the forms
			\[
			\{[R_s]:s\in\mu_N\},
			\qquad
			\{[R_s],[ZR_s]:s\in\mu_N\},
			\qquad
			R_s=\frac12
			\begin{pmatrix}1+s&1-s\\1-s&1+s\end{pmatrix}.
			\]
			Every displayed projective class has a fixed directly realizable
			ordered-binary representative over \(D\cG\).
		\end{enumerate}
	\end{theorem}
	
	\begin{proof}
		Put \(G=G(\cG)\) and
		\[
		H=\angles{G,[Z]}.
		\]
		All gadget entries lie in the number field generated by the entries of
		\(\cG\), so \cref{lem:bounded-torsion} makes \(G\) finite.  If
		\([T]\in G\), port reversal and \cref{eq:transfer-reversal} give
		\([ZT^{-1}Z]\in G\).  Since inversion permutes \(G\), it follows that
		\([Z]G[Z]=G\).  Consequently
		\[
		G\triangleleft H,\qquad [H:G]\le2,
		\]
		and \(H=G\) or \(H=G\sqcup[Z]G\), so \(H\) is finite.
		
		For the distinguished mountain,
		\[
		T=BX=\begin{pmatrix}b&a\\d&c\end{pmatrix}
		\]
		has class \([T]\in G\).  Since \(a,d\ne0\), both off-diagonal entries of
		\(T\) are nonzero.  We first record a weighted equality.  Suppose \([S]\in G\)
		has an anti-diagonal representative
		\[
		S=\begin{pmatrix}0&p\\q&0\end{pmatrix},
		\qquad pq\ne0.
		\]
		By the definition of \(G\), there is a directly realizable ordered binary
		\(B_S\) and a scalar \(\rho\ne0\) such that \(B_SX=\rho S\).  Hence
		\[
		B_S=\rho SX=\rho\diag(p,q),
		\]
		which is the nonsingular weighted equality in outcome~1.
		
		We now use
		\cref{prop:finite-pgl2-classification,lem:finite-pgl2-facts}.  Their use
		here is only group-theoretic: the distinguished class is the fixed
		\([Z]\), and no conjugating matrix from the abstract classification is
		applied to the signature set.
		
		If \(H\) is cyclic, it is abelian, so \([T]\) commutes with \([Z]\).  The
		projective centralizer of \([Z]\) consists of the diagonal and
		anti-diagonal classes.  The two nonzero off-diagonal entries of \(T\)
		exclude the diagonal case; thus \([T]\in G\) is anti-diagonal and gives the
		weighted equality.
		
		Suppose \(H\cong A_4\).  Since \(A_4\) has no index-two subgroup, the
		index bound forces \(G=H\).  The centralizer of \([Z]\) in \(H\) is a
		\(V_4\), so it contains an involution \(u\ne1,[Z]\).  The class \(u\)
		belongs to \(G\) and commutes with \([Z]\).  Because \([Z]\) is the only
		nonidentity diagonal involution, \(u\) is anti-diagonal.  The same argument
		applies when \(H\cong A_5\): simplicity gives \(G=H\), and the
		centralizer of an involution is again \(V_4\).
		
		Suppose \(H\cong S_4\).  If \(G=H\) and \([Z]\) corresponds to a
		transposition, choose the disjoint transposition; if \([Z]\) corresponds
		to a double transposition, choose either of the other two double
		transpositions.  In both cases the chosen involution lies in \(G\), commutes
		with \([Z]\), and differs from it, so it is anti-diagonal.  If \(G\ne H\),
		then the unique index-two subgroup gives \(G=A_4\).  Here \([Z]\notin G\)
		is an involution outside \(A_4\), hence a transposition.  If \(z'\) is the
		disjoint transposition, then \(u=[Z]z'\) is a double transposition in
		\(A_4=G\); it commutes with \([Z]\), differs from it, and is therefore
		anti-diagonal.
		
		It remains to consider
		\[
		H=D_{2m}=\angles{r,z:r^m=z^2=1,\ zrz=r^{-1}},
		\]
		where the distinguished reflection is denoted by \(z=[Z]\) when \([Z]\)
		is a reflection.  If \(m\) is even and \([Z]=r^{m/2}\) is the central
		half-turn, then \([T]\in G\) commutes with \([Z]\); its two nonzero
		off-diagonal entries again make it anti-diagonal.  Suppose instead that
		\([Z]\) is a reflection and \(m\) is even.  The two involutions
		\[
		r^{m/2},\qquad [Z]r^{m/2}
		\]
		commute with \([Z]\) and differ from it.  If \(G=H\), both belong to
		\(G\).  If \([H:G]=2\), then \([Z]\notin G\), and the two displayed
		elements lie in opposite cosets of \(G\); hence exactly one belongs to
		\(G\).  In either event \(G\) contains an involution different from
		\([Z]\) that commutes with \([Z]\), and this involution is anti-diagonal.
		
		The sole unresolved case is therefore
		\[
		H\cong D_{2N},\qquad N\ge3\text{ odd},\qquad [Z]\text{ a reflection},
		\]
		with \(G=C_N\) or \(G=D_{2N}\).  Indeed, for odd \(N\) the rotation
		subgroup is the unique index-two subgroup; thus a proper \(G\) is the
		rotation subgroup, while otherwise \(G=H\).
		
		Choose a rotation class \([R]\) of exact order \(N\), and write
		\[
		R=\begin{pmatrix}u&v\\w&x\end{pmatrix}.
		\]
		The dihedral relation gives \([ZRZ]=[R^{-1}]\), while
		\[
		ZRZ=\begin{pmatrix}u&-v\\-w&x\end{pmatrix},
		\qquad
		R^{-1}\doteq\begin{pmatrix}x&-v\\-w&u\end{pmatrix}.
		\]
		The entries \(v,w\) cannot both vanish: otherwise \([R]=[R^{-1}]\), so
		\([R]^2=1\), contrary to the odd order \(N\ge3\).  Since at least one of
		\(v,w\) is nonzero, comparison in the preceding projective equality forces
		the projective scalar to be \(1\), and hence \(u=x\).  Exactly one of
		\(v,w\) cannot vanish.  In that case nonsingularity gives \(u=x\ne0\), and
		a projective rescaling makes \(R\) a nonidentity unipotent matrix; its
		powers are pairwise distinct in characteristic zero.  Therefore
		\[
		R=\begin{pmatrix}\alpha&\beta\\ \gamma&\alpha\end{pmatrix},
		\qquad \beta\gamma\ne0.
		\]
		
		Choose \(q\in\AlgNums^\times\) with \(q^2=\gamma/\beta\), and put
		\[
		D=\diag(q,1).
		\]
		Then
		\[
		D^{\mathsf T}XD=qX,
		\]
		so \(D\in GO(X)\); moreover \(DZD^{-1}=Z\).  With
		\(\delta=q\beta=q^{-1}\gamma\),
		\[
		DRD^{-1}
		=\begin{pmatrix}\alpha&\delta\\ \delta&\alpha\end{pmatrix}.
		\]
		The matrix \(P\) is used only for the following analytic diagonalization:
		\[
		P^{-1}(DRD^{-1})P
		=\diag(\alpha+\delta,\alpha-\delta).
		\]
		Both eigenvalues are nonzero because \(R\) is nonsingular.  Put
		\[
		s_0=\frac{\alpha-\delta}{\alpha+\delta}.
		\]
		The exact projective order of \([R]\) implies that \(s_0\) is a primitive
		\(N\)-th root of unity, and
		\[
		DRD^{-1}
		=(\alpha+\delta)R_{s_0},
		\qquad
		R_s=\frac12\begin{pmatrix}1+s&1-s\\1-s&1+s\end{pmatrix}.
		\]
		Since \(R_sR_t=R_{st}\), the normalized rotation classes are precisely
		\[
		\{[R_s]:s\in\mu_N\}.
		\]
		The diagonal matrix \(D\) fixes \([Z]\).  Thus, when \(G=C_N\), these are
		all normalized classes, and, when \(G=D_{2N}\), the remaining classes are
		precisely
		\[
		\{[ZR_s]:s\in\mu_N\}.
		\]
		
		Finally, the normalization gives equality of complete groups, not merely
		an inclusion.  If \(B_0\) is a directly realizable ordered binary over
		\(\cG\), legal global transport gives a directly realizable boundary tensor
		\(B_0^D\) over \(D\cG\) with
		\[
		B_0^D\doteq DB_0D^{\mathsf T},
		\qquad
		(DB_0D^{\mathsf T})X=qD(B_0X)D^{-1}.
		\]
		Hence \(T_{B_0^D}\doteq D T_{B_0}D^{-1}\).  Applying the same argument
		with \(D^{-1}\in GO(X)\) gives the converse, so
		\[
		G(D\cG)=[D]G(\cG)[D]^{-1}.
		\]
		Every displayed class therefore has a fixed directly realizable
		ordered-binary representative over \(D\cG\), with the nonzero scalar from
		the global transport recorded and restored.  This is a single legal
		\(GO(X)\)-normalization; neither \(P\) nor an arbitrary conjugator from the
		finite-subgroup classification is used as a gadget or a global
		normalization.
	\end{proof}
	
	\subsection{Odd cyclic and dihedral transfer groups}
	\label{sec:odd-residue}
	
	\begin{definition}[Equivalent updates with strict progress]
		\label{def:height-one-outer-successor}
		Let \((A,\Gamma)\) satisfy \cref{lem:normalization-ledger} for the initial
		set \(\Gamma_0\) and its fixed number field \(L_0\).  An update to
		\((A^+,\Gamma^+)\) must satisfy the same invariant and use only the
		augmentations and global transformations allowed there.  In particular,
		\(\Gamma^+\) is finite, preprocessed, algebraic, and factor-closed, and
		\[
		\begin{gathered}
			A^+\Gamma_0\subseteq\Gamma^+,
			\qquad (A^+)^{\mathsf T}XA^+=\mu_{A^+}X\quad(\mu_{A^+}\ne0),\\
			\KHolant(\Gamma^+)\equivT\KHolant(\Gamma).
		\end{gathered}
		\]
		The transformed height-one mountain and every old tensor-prime signature
		remain members up to the same invertible change of basis and their recorded
		nonzero scalars.  Every nonzero directly realizable ordered binary over
		both sets is nonsingular with finite projective transfer order.  Put
		\[
		H=[A]^{-1}G(\Gamma)[A],
		\qquad H^+=[A^+]^{-1}G(\Gamma^+)[A^+].
		\]
		Both groups are subgroups of \(\PGL_2(L_0)\).  When both values of \(\nu\)
		are defined, the update is required to satisfy
		\[
		H\subseteq H^+,
		\qquad \nu(\Gamma^+)\le\nu(\Gamma),
		\qquad
		H\ne H^+\ \text{or}\ \nu(\Gamma^+)\ne\nu(\Gamma).
		\]
		Thus at least one comparison is strict.  If no nonbinary tensor-prime
		signature remains, the arity-at-most-two alternative gives tractability
		and no further update is made.
	\end{definition}
	
	\begin{lemma}[Termination of the strict updates]
		\label{lem:height-one-progress-well-founded}
		Fix the number field \(L_0\) in
		\cref{lem:normalization-ledger}.  There is no infinite sequence
		\[
		(H_0,\nu_0),(H_1,\nu_1),\ldots
		\]
		such that every \(H_i\) is a finite subgroup of \(\PGL_2(L_0)\), every
		\(\nu_i\) is an integer at least three, and every update satisfies
		\[
		H_i\subseteq H_{i+1},
		\qquad
		\nu_{i+1}\le\nu_i,
		\]
		with at least one of these two inequalities strict.  In the height-one
		construction these quantities are
		\[
		H_i=[A_i]^{-1}G(\Gamma_i)[A_i],
		\qquad
		\nu_i=\min\{\arity(f):f\in\Gamma_i
		\text{ is tensor-prime and }\arity(f)\ge3\}.
		\]
		If the second set is empty, the arity-at-most-two case is tractable and no
		further strict update is made.  This assertion concerns only the
		transitions represented by the displayed pairs; operations preserving both
		\(H_i\) and \(\nu_i\) are handled by the finite recursive extension of
		\cref{lem:height-one-card-saturation}.
	\end{lemma}
	
	\begin{proof}
		Write \(\delta=[L_0:\mathbb Q]\).  If
		\([T]\in\PGL_2(L_0)\) has projective order \(m\), the ratio of the two
		eigenvalues of \(T\) is a primitive \(m\)-th root of unity \(\zeta_m\).
		The eigenvalues lie in an extension of \(L_0\) of degree at most two, whence
		\[
		\varphi(m)=[\mathbb Q(\zeta_m):\mathbb Q]
		\leq [L_0(\zeta_m):\mathbb Q]\leq2\delta.
		\]
		There are therefore only finitely many possible projective element orders;
		let \(M_0\) be their maximum.  By
		\cref{prop:finite-pgl2-classification}, every finite subgroup of
		\(\PGL_2(L_0)\) has order at most
		\[
		C_0=\max\{2M_0,60\}.
		\]
		
		For a sequence in the statement, define
		\[
		\Phi_i=(C_0-|H_i|)+(\nu_i-3)\in\mathbb Z_{\ge0}.
		\]
		The inclusions and inequalities make \(\Phi_i\) nonincreasing, and the strict
		condition decreases it by at least one at every update.  Hence the sequence
		has at most \(\Phi_0\) updates.  In the application,
		\cref{lem:normalization-ledger} places every pulled-back group in
		\(\PGL_2(L_0)\), equivalent enrichment makes the groups nested, and
		\cref{def:height-one-outer-successor} makes \(\nu\) nonincreasing.  Factor
		closure between updates is finite because every nontrivial tensor split replaces
		a signature by positive-arity factors of strictly smaller arity; it is not
		part of the displayed measure.
	\end{proof}
	
	For either group with odd \(N\), the proof will establish that the single matrix
	\(P\) fixed in \cref{eq:Walsh-coordinates} satisfies
	\begin{equation}
		(P^{-1})^{\otimes\arity(f)}f\in\cP
		\quad(f\in\Lambda),
		\qquad
		P^{\mathsf T}XP=2Z\in\cP.
		\label{eq:common-Walsh-product-presentation}
	\end{equation}
	These are the simultaneous signature and edge conditions for a common
	\(K\)-presentation in \(\cP\); the same matrix must handle every member
	of \(\Lambda\).
	
	\begin{theorem}[The odd-dihedral case]
		\label{thm:odd-dihedral}
		Let \((A,\cG)\) be a stage of the sequence in
		\cref{lem:normalization-ledger} at which \(\cG\) is finite, preprocessed,
		algebraic, and factor-closed and contains the transported height-one
		mountain.  Suppose that, after the normalization of
		\cref{thm:finite-group-split}, every nonzero directly realizable binary is
		nonsingular with finite projective transfer order, and
		\(G(\cG)\) is
		\[
		D_{2N}=\{[R_s],[ZR_s]:s\in\mu_N\}
		\]
		for an odd \(N\ge3\).
		Then at least one of the following holds.
		\begin{enumerate}
			\item \(\KHolant(\cG)\) is \(\SharpP\)-hard or \(\Tract(K\cG)\) holds.
			\item There is a later stage \((A^+,\cG^+)\) of the same sequence at which
			\(\cG^+\) is finite, preprocessed, algebraic, and factor-closed and
			still contains the correspondingly transported height-one mountain,
			such that
			\[
			\KHolant(\cG^+)\equivT\KHolant(\cG),
			\]
			every nonzero directly realizable ordered binary over \(\cG^+\) is
			nonsingular with finite projective transfer order, and either
			\[
			[A]^{-1}G(\cG)[A]
			\subsetneq
			[A^+]^{-1}G(\cG^+)[A^+]
			\]
			or, with both values defined, \(\nu(\cG^+)<\nu(\cG)\).
			In either case the full update conditions of
			\cref{def:height-one-outer-successor} hold.
			\item There is an equivalent finite preprocessed algebraic factor-closed
			set \(\cG^{\mathrm{sat}}\supseteq\cG\), in the same
			coordinates and with the same matrix \(A\), for which the fixed
			Walsh matrix \(P\) satisfies
			\[
			(P^{-1})^{\otimes\arity(f)}f\in\cP
			\quad(f\in\cG^{\mathrm{sat}}),
			\qquad
			P^{\mathsf T}XP=2Z\in\cP.
			\]
			Hence \(K\cG\) is \(\cP\)-transformable.
		\end{enumerate}
	\end{theorem}
	
	Fix a signature set \(\Gamma\) for which \(G(\Gamma)\) is defined, a
	subgroup \(H\le G(\Gamma)\), and one directly realizable ordered binary
	\(B_h\) with \([B_hX]=h\) for each \(h\in H\).  We will refer to the
	following contraction condition for \(f\in\Gamma\):
	\begin{enumerate}[label=\textnormal{(C)},leftmargin=*]
		\item\label{cond:height-one-contractions}
		For every ordered pair of distinct ports of \(f\), every chosen \(B_h\),
		and both kernel orientations \(XB_hX\) and \(XB_h^{\mathsf T}X\), the
		contraction is zero, or all its positive-arity tensor-prime factors are
		nonsingular binaries \(B\) with \([BX]\in H\).
	\end{enumerate}
	Any nullary factors are absorbed into the recorded scalar.  A zero
	contraction is allowed and is not assigned a projective transfer class.
	The conditions are in the current \(K\)-coordinates; after a Walsh change
	of coordinates the group membership is conjugated by \(P\) as well.
	Applications below take \(H=G(\Gamma)\), with the complete group from
	\cref{eq:complete-safe-group}.  Before using
	condition~\ref{cond:height-one-contractions}, any factor that gives
	hardness, tractability, or an update in
	\cref{def:height-one-outer-successor} is handled separately.  This condition
	does not assert that arbitrary gadget outputs have been adjoined.
	
	\begin{lemma}[Finite recursive contraction and factor extraction]
		\label{lem:height-one-card-saturation}
		Let \((A,\Gamma)\) be a stage of a sequence satisfying
		\cref{lem:normalization-ledger}, over its fixed number field \(L_0\).
		Suppose that \(\Gamma\) is finite,
		preprocessed, algebraic, and factor-closed, contains the transported height-one mountain,
		and every nonzero directly realizable ordered binary over \(\Gamma\) is
		nonsingular and has finite projective transfer order.  Let
		\(H=G(\Gamma)\), and fix one directly realizable ordered-binary
		representative of each class in this finite group.
		
		Start with every nonbinary tensor-prime factor currently in \(\Gamma\).  For
		each retained factor \(f\), process every choice of an ordered pair of
		distinct ports, a fixed representative of a class in \(H\), and either
		kernel orientation.  Adjoin the resulting nonzero two-port contraction and, through
		\cref{lem:factor-saturation}, all of its positive-arity tensor-prime factors;
		apply the same procedure to each newly retained nonbinary factor.  This
		finite recursive procedure reaches one of the following outcomes:
		\begin{enumerate}
			\item \(\KHolant(\Gamma)\) is \(\SharpP\)-hard or
			\(\Tract(K\Gamma)\) holds;
			\item there are a matrix \(A^+\) and a finite preprocessed algebraic
			factor-closed set \(\Gamma^+\), continuing the same coordinate-transport
			sequence and still containing the correspondingly transported
			height-one mountain, with
			\(\KHolant(\Gamma^+)\equivT\KHolant(\Gamma)\), such that
			every nonzero directly realizable ordered binary over \(\Gamma^+\) is
			nonsingular with finite projective transfer order, and either
			\[
			[A]^{-1}G(\Gamma)[A]
			\subsetneq
			[A^+]^{-1}G(\Gamma^+)[A^+]
			\]
			or, with both values defined,
			\(\nu(\Gamma^+)<\nu(\Gamma)\).
			The update satisfies \cref{def:height-one-outer-successor}, including
			the non-strict comparisons for both quantities;
			\item there is a finite preprocessed algebraic factor-closed set
			\(\Gamma^{\mathrm{sat}}\supseteq\Gamma\), with
			\[
			\KHolant(\Gamma^{\mathrm{sat}})\equivT\KHolant(\Gamma),
			\qquad G(\Gamma^{\mathrm{sat}})=H,
			\]
			for which the coordinate matrix remains \(A\).  Every indicated
			two-port contraction has been processed, every
			nonbinary tensor-prime factor supplied by factor closure is retained in
			\(\Gamma^{\mathrm{sat}}\), and every binary factor it produces is
			nonsingular with transfer class in \(H\).  A nonbinary factor produced
			from a parent \(f\) has arity at most \(\arity(f)-2\).
		\end{enumerate}
		The third outcome may preserve \(\nu\); it is exhaustive closure with the
		same displayed group and coordinate matrix.  It does not require a strict
		change in either quantity.
	\end{lemma}
	
	\begin{proof}
		Begin with the finitely many nonbinary tensor-prime factors in
		\(\Gamma\).  For each retained factor, process its finite set of
		ordered-port, representative, and orientation choices.  For each choice,
		adjoin the resulting contraction and all positive-arity tensor-prime factors
		under \cref{eq:finite-equivalent-enlargement,eq:factor-saturation}.
		The transported mountain rules out the common-one-sided factor exception;
		if a factor theorem instead proves hardness or tractability, this is
		outcome~1.
		
		An odd or unary factor is handled by the odd-arity or unary theorem.
		A singular binary, or a nonsingular binary whose transfer has infinite
		projective order, is handled by \cref{lem:binary-interpolation} and those
		same theorems.  After each augmentation, consider all directly realizable
		nonzero ordered binaries.  If any is singular or its transfer has infinite
		projective order, the preceding argument gives outcome~1.
		Otherwise the group is defined, and an element outside \(H\) strictly
		enlarges the pulled-back group,
		and a new nonbinary tensor-prime factor of arity below \(\nu(\Gamma)\)
		strictly decreases \(\nu\).  These give outcome~2, with both comparisons
		preserved as in \cref{def:height-one-outer-successor}.
		If neither occurs, retain all the binary factors and repeat the procedure
		on each newly retained nonbinary tensor-prime factor.
		
		There are finitely many initial factors and finitely many choices at each
		step.  Every newly produced nonbinary factor has arity at most the arity of
		its parent minus two.  The recursive depth is therefore bounded, so only
		finitely many factors and contraction choices occur.  If neither of the
		first two outcomes occurs, processing all these choices gives outcome~3.
		Every augmentation retains the transformed mountain and satisfies the
		fixed-field and scalar conditions of \cref{lem:normalization-ledger}.
	\end{proof}
	
	\subsection{Contractions for the odd cyclic and dihedral groups}
	\label{subsec:residual-phase-toolkit}
	
	The next two lemmas are shared by the odd cyclic and odd dihedral cases.
	For a signature \(f\in\cG\), recall that a tilde denotes only its
	Walsh coordinates
	\[
	\widetilde f=(P^{-1})^{\otimes\arity(f)}f.
	\]
	The matrix \(P\) is never inserted into a gadget: every contraction below is a
	native \(K\)-gadget, and every pencil identity is only an analysis of those
	contractions; no displayed linear combination is asserted to be realizable.
	
	\subsubsection{Contraction kernels and the four-coefficient relation}
	\label{subsec:odd-dih-cards}
	
	Fix odd \(N\ge3\).  With \(R_s\) as in
	\cref{thm:finite-group-split}, the actual rotation and reflection classes are
	\begin{equation}
		[R_s],\qquad[ZR_s],
		\qquad s\in\mu_N.
		\label{eq:odd-dih-native-transfers}
	\end{equation}
	
	They have Walsh forms
	\[
	P^{-1}R_sP=D_s:=\diag(1,s),
	\qquad P^{-1}ZP=X,
	\qquad F_s:=XD_s.
	\]
	
	Fix a direct binary-chain representative of every projective class in
	\cref{eq:odd-dih-native-transfers}.  Contracting two ports then gives, up to
	the recorded nonzero representative scalar, the Walsh kernels
	
	\begin{equation}
		\mathsf K_s^+=
		\begin{pmatrix}1&0\\0&-s\end{pmatrix},
		\qquad
		\mathsf K_s^-=
		\begin{pmatrix}0&1\\-s&0\end{pmatrix},
		\qquad s\in\mu_N.
		\label{eq:odd-dih-kernels}
	\end{equation}
	
	The plus kernels occur in both groups; the minus kernels occur only in the
	dihedral group.  Under port reversal, \(D_s\) is fixed, whereas \(F_s\) is
	sent projectively to \(F_{s^{-1}}\).  Attaching binary representatives to the ports adds no kernels
	because
	
	\[
	\tau_Z(T):=ZT^{\mathsf T}Z
	\]
	
	preserves \(\{D_s:s\in\mu_N\}\) and
	\(\{D_s,F_s:s\in\mu_N\}\).  Hence \cref{eq:odd-dih-kernels} lists every
	oriented kernel needed below.
	
	Although this subsection otherwise has \(N\) odd, the next pencil fact is
	valid for every integer \(N\ge3\) and will also be used in the even-order
	cyclic case.  For such an \(N\) and
	\(a,b,c,d\in\AlgNums\), let \(\mathcal R_N(a,b,c,d)\) mean that, for some
	\(\alpha\in\mu_N\), one of the following three systems holds:
	
	\begin{align}
		c&=-\alpha a,&d&=-\alpha b,
		&&\text{(proportional)},
		\label{eq:odd-dih-R-proportional}\\
		b&=c=0,&d&=\alpha a,
		&&\text{(diagonal monomial)},
		\label{eq:odd-dih-R-diagonal}\\
		a&=d=0,&c&=\alpha b,
		&&\text{(anti-diagonal monomial)}.
		\label{eq:odd-dih-R-antidiagonal}
	\end{align}
	
	The three systems are not required to be disjoint; in particular, their
	union contains the all-zero quadruple.
	
	\begin{lemma}[Shared cyclic phase pencil, including zeros]
		\label{lem:odd-dih-RN}
		For every integer \(N\ge3\), the relation
		\(\mathcal R_N(a,b,c,d)\) is equivalent to the following
		condition: for every \(s\in\mu_N\),
		\[
		\diag(a-sb,c-sd)Z
		\quad\text{is zero or is projectively }D_\eta
		\text{ for some }\eta\in\mu_N.
		\]
	\end{lemma}
	
	\begin{proof}
		This is \cite[Lemma~7.1]{Liu2026DEQ2Holant}.
	\end{proof}
	
	\subsubsection{Support parity}
	\label{subsec:odd-dih-parity}
	
	\begin{lemma}[Walsh parity from contractions]
		\label{lem:odd-dih-all-arity-parity}
		Let \(N\ge3\) be odd, let \(H=C_N\) or \(D_{2N}\), and let \(f\ne0\) have
		positive even arity.  At arity two, suppose the Walsh transfer of \(f\)
		belongs to \(H\).  At arity at least four, suppose that every positive-arity
		prime factor of every nonzero contraction of \(f\) using an ordered pair of distinct
		ports, a chosen representative of \(H\), and either kernel orientation is supported on
		one Walsh parity class; in the cyclic case suppose that class is even.  Then
		\(\widetilde f\) is supported on one parity class, which is even in the
		cyclic case.
	\end{lemma}
	
	\begin{proof}
		At arity two, a Walsh binary is, up to a nonzero scalar, a diagonal or
		anti-diagonal transfer times \(Z\), and hence has one parity.  Only the
		diagonal alternative occurs in the cyclic group.  At larger arity, the
		tensor product of the prime factors of a nonzero contraction has one parity; in the
		cyclic case it is even.
		
		At arity four, each \(\mathsf K_s^+\)-contraction is zero, diagonal, or
		anti-diagonal, and every entry is affine in \(s\).  Among three phases,
		either two contractions vanish, annihilating the pencil, or two nonzero contractions have
		the same support parity, which forces the opposite entries to vanish
		identically.  In the only remaining pattern---one zero, one diagonal, and
		one anti-diagonal contraction---each opposite pair acquires two roots, with the
		same conclusion.  Hence every equal-bit slice has one parity.
		Two supported words of opposite parity have a coordinate pair equal in both
		unless their four ordered coordinate pairs exhaust \(\Ftwo^2\); in the
		latter case both words have weight two and therefore already have the same
		parity, a contradiction.  In the former case deleting the repeated pair
		contradicts slice parity.  In the cyclic case, odd Walsh parity would
		realize a reflection and is therefore excluded.
		
		For arity at least six, fix distinct ports \(i,j\).  The actual plus contraction is
		\[
		\widetilde h_s(y)=\lambda_s\bigl(\widetilde f(00,y)
		-s\widetilde f(11,y)\bigr),
		\qquad \lambda_s\ne0.
		\]
		A zero contraction makes all others proportional.  Otherwise two phases have the
		same contraction parity.  Every coordinate in the opposite parity then
		vanishes at two distinct phases, so its affine coefficient vanishes
		identically; the two nonzero phases also recover both endpoint slices with
		the common parity.
		
		In the cyclic case, delete two equal coordinates of any supported word and
		choose a phase other than the at most one value cancelling its residual
		coefficient; the resulting contraction is even.  In the dihedral case, two
		opposite-parity supported words at arity at least six have two coordinates
		with the same ordered pair.  At most one phase cancels each residual
		coefficient, so \(|\mu_N|\ge3\) leaves a phase preserving both words.
		Deleting the chosen pair preserves their opposite parities, contradicting
		the common slice parity.
	\end{proof}
	
	For every \(f\) covered by \cref{lem:odd-dih-all-arity-parity}, Walsh
	duality gives the equivalent complement relation in the native
	\(K\)-coordinates,
	\[
	f(\bar x)=f(x)
	\quad\text{or}\quad
	f(\bar x)=-f(x),
	\]
	
	with only the plus sign in the cyclic subcase.
	
	\subsection{The odd-dihedral case}
	\label{proofsec:odd-dihedral}
	
	\subsubsection{Support and coefficients on even blocks}
	\label{subsec:odd-dih-phase-product-class}
	
	For a finite port set \(V\) and \(C\subseteq V\), let
	\(\one_C\in\Ftwo^V\) denote the incidence vector of \(C\).  For a partition
	\(\pi\) of \(V\) into nonempty even blocks, define the subspace
	\[
	L_\pi=\Span_{\Ftwo}\{\one_C:C\in\pi\}.
	\]
	Its cosets are precisely the sets obtained by toggling each block as one
	unit.
	
	The class \(\mathcal P_N(V)\) consists of the zero tensor and
	all tensors \(g\) for which there exist such a partition \(\pi\), an offset
	\(a\in\Ftwo^V\), a scalar \(c\in\AlgNums\setminus\{0\}\), and block ratios
	\(r_C\in(-1)^{|C|/2}\mu_N\) such that
	\begin{align*}
		\supp(g)&=a+L_\pi,
		\\
		g\!\left(a+\sum_{C\in\pi}t_C\one_C\right)
		&=c\prod_{C\in\pi}r_C^{t_C}
		\qquad\bigl((t_C)_{C\in\pi}\in\Ftwo^\pi\bigr).
	\end{align*}
	For \(C\subseteq V\), write \(a|_C\) for the restriction of \(a\) to \(C\),
	and \(\overline{a|_C}\) for its bitwise complement.
	Equivalently, with the tensor factors placed on their indicated port sets,
	\[
	g=c\bigotimes_{C\in\pi}
	\bigl(\delta_{a|_C}+r_C\delta_{\overline{a|_C}}\bigr).
	\]
	Thus, in the variables \((t_C)\), the coefficient function is the product
	\(c\prod_C r_C^{t_C}\).  A factor on one block has the displayed two-point
	form and the stated ratio restriction.  For nonzero \(g\), its support
	determines \(\pi\): the inclusion-minimal nonzero vectors of \(L_\pi\) are
	precisely the block indicators.  We omit \(V\) when clear.  The class
	\(\mathcal P_N\) differs from the standard product class \(\cP\) of
	\cref{subsec:standard-classes}; its even-block and coefficient-ratio
	restrictions remain part of its definition.
	
	\begin{lemma}[Closure of \(\mathcal P_N\) under contractions]
		\label{lem:phase-product-closure}
		The class \(\mathcal P_N\) is closed under nonzero scalar multiplication,
		port permutation, tensor product on disjoint port sets, and, up to the zero
		tensor, contraction of two ports against any kernel
		\(\mathsf K_s^+\), \(\mathsf K_s^-\), or its transpose from
		\cref{eq:odd-dih-kernels}, where \(s\in\mu_N\).  Consequently, if
		every positive-arity prime factor of a contraction has Walsh tensor in
		\(\mathcal P_N\), then the contraction itself has Walsh tensor in
		\(\mathcal P_N\).
	\end{lemma}
	
	\begin{proof}
		For a tensor product, take the disjoint union of the two even-block
		partitions and concatenate their offsets and block ratios.  Scalars and port
		permutations merely change the leading coefficient or relabel that data.
		
		Consider a two-port contraction of a nonzero tensor in \(\mathcal P_N\).  If the
		ports lie in the same block \(C\), the kernel with the wrong xor gives zero.
		For the supported xor, delete the two ports.  When ports remain in \(C\), its
		new block ratio is
		\[
		-s r_C\quad\hbox{or}\quad-r_C/s,
		\]
		and therefore belongs to
		\((-1)^{(|C|-2)/2}\mu_N\).  If \(|C|=2\), the block instead contributes a
		scalar, possibly zero, which is absorbed into the leading coefficient.
		
		If the ports lie in distinct blocks \(C,D\), the kernel affinely identifies
		their two block variables.  After complementing one residual offset when
		needed, replace those blocks by
		\[
		(C\setminus\{i\})\mathbin{\dot\cup}(D\setminus\{j\}).
		\]
		Its ratio is one of
		\[
		-s^{\pm1}r_Cr_D^{\pm1}
		\in(-1)^{(|C|+|D|-2)/2}\mu_N.
		\]
		All other blocks are unchanged.  Transposing \(\mathsf K_s^-\) only supplies
		a nonzero scalar and replaces \(s\) by \(s^{-1}\), so it is covered as well.
		Thus every nonzero contraction is again in \(\mathcal P_N\).
		The last assertion follows from prime factorization; nullary factors are
		absorbed into the scalar.
	\end{proof}
	
	\subsubsection{Four-port support and coefficients}
	\label{subsec:odd-dih-q4}
	
	Fix \(N\ge3\) and a quaternary signature \(q\) in the current set, and put
	\(\widetilde q=(P^{-1})^{\otimes4}q\).  When \(\widetilde q\) is supported
	on the even Hamming-weight positions in Walsh coordinates, write
	
	\begin{equation}
		p=\widetilde q_{0000},
		\qquad t=\widetilde q_{1111},
		\qquad x_{ab}=\widetilde q_{\one_{\{a,b\}}}.
		\label{eq:odd-dih-q4-coordinates}
	\end{equation}
	
	For every ordering \((i,j,k,l)\) of \([4]\), the equal and mixed phase-contraction
	matrices are, respectively,
	\begin{align*}
		H_{ij}^+(s)&=
		\begin{pmatrix}p-sx_{ij}&0\\0&x_{kl}-st\end{pmatrix},\\
		H_{ij}^-(s)&=
		\begin{pmatrix}
			0&x_{jl}-sx_{il}\\x_{jk}-sx_{ik}&0
		\end{pmatrix}.
	\end{align*}
	Thus condition~\ref{cond:height-one-contractions} for the cyclic group
	and \cref{lem:odd-dih-RN} give \(\mathcal R_N(p,x_{ij},x_{kl},t)\).
	For the dihedral group, the mixed contraction also gives
	\(\mathcal R_N(x_{jk},x_{ik},x_{jl},x_{il})\).
	
	When \(\widetilde q\) is supported on the odd Hamming-weight positions in Walsh coordinates, write
	\[
	u_i=\widetilde q(e_i),
	\qquad v_i=\widetilde q(\one^4\oplus e_i).
	\]
	
	For every such port ordering, condition~\ref{cond:height-one-contractions}
	for the dihedral group gives
	\[
	\mathcal R_N(u_k,v_l,u_l,v_k),
	\qquad
	\mathcal R_N(u_j,u_i,v_i,v_j).
	\]
	
	\begin{lemma}[Shared cyclic quaternary family]
		\label{lem:odd-cyc-q4-family}
		Let \(N\ge3\), let \(\Lambda\) be a finite algebraic signature set in
		\(K\)-coordinates, and, for every \(s\in\mu_N\), fix an actual ordered
		binary representative \(C_s\) over \(\Lambda\), in a specified orientation,
		whose Walsh matrix is
		\[
		\widetilde C_s=\kappa_sD_sZ,
		\qquad \kappa_s\in\AlgNums^\times.
		\]
		Let \(q\in\Lambda\) be quaternary, and
		suppose \(\widetilde q=(P^{-1})^{\otimes4}q\) is supported on the even Hamming-weight
		positions in Walsh coordinates.  Define \(p,t,x_{ab}\) by
		\cref{eq:odd-dih-q4-coordinates}.  Assume that, for every complementary pair,
		both orders of its two residual ports, and every \(s\in\mu_N\), contraction
		with \(C_s\) gives either the zero contraction or a nonsingular binary contraction whose
		Walsh transfer class belongs to
		\[
		C_N=\{[D_\eta]:\eta\in\mu_N\}.
		\]
		Organize the eight
		coordinates into four complementary pairs.  Put \(B_0=(p,t)\) and
		\(\mathfrak I=\{\{1,2\},\{1,3\},\{1,4\}\}\).
		For \(I=\{a,b\}\in\mathfrak I\), set \(x_I=x_{ab}\),
		put \(I^c=[4]\setminus I\), and define
		the other three complementary pairs \(B_I=(x_I,x_{I^c})\),
		\(I\in\mathfrak I\).  Here \(B_I\ne(0,0)\) initially means only that
		at least one coefficient is nonzero; the conclusion below makes both
		coefficients nonzero.  Then exactly one of the following holds:
		\begin{enumerate}
			\item If \(p,t\ne0\), then \(\lambda=t/p\in\mu_N\), and each pair \(B_I\)
			is zero or
			\begin{equation}
				B_I=(-p\lambda\alpha_I^{-1},-p\alpha_I),
				\qquad \alpha_I\in\mu_N.
				\label{eq:odd-cyc-q4-nonzero-endpoints}
			\end{equation}
			\item If \(p=t=0\), each pair \(B_I\) is zero or
			\[
			B_I=\rho_I(1,\alpha_I),
			\qquad \rho_I\ne0,
			\quad \alpha_I\in\mu_N.
			\]
		\end{enumerate}
		In particular, exactly one of \(p,t\) cannot be nonzero.
		Conversely, every even Walsh table in these alternatives satisfies the
		stated contraction condition under every phase contraction.
	\end{lemma}
	
	\begin{proof}
		This is \cite[Lemma~7.2]{Liu2026DEQ2Holant}.
	\end{proof}
	
	\begin{proposition}[Dihedral quaternary classification]
		\label{prop:odd-dih-q4}
		Let \(N\ge3\) be odd, and let \(q\in\Gamma\) be a nonzero quaternary
		satisfying condition~\ref{cond:height-one-contractions} with
		\(H=G(\Gamma)=D_{2N}\) in the normalized coordinates above.  Put
		\(\widetilde q=(P^{-1})^{\otimes4}q\).  Then
		\(\widetilde q\in\mathcal P_N([4])\).  More precisely,
		\(\widetilde q\) has one Walsh parity and exactly one or two complementary
		pairs whose coefficient pairs are nonzero.  After choosing an orientation, there is a scalar
		\(\rho\in\AlgNums\setminus\{0\}\) such that, for one nonzero complementary
		pair \(\{x,\bar x\}\), its nonzero coefficients in the order
		\((\widetilde q_x,\widetilde q_{\bar x})\) are
		
		\begin{equation}
			\rho(1,\delta),
			\qquad \delta\in\mu_N,
			\label{eq:odd-dih-q4-one-block}
		\end{equation}
		
		or, for two nonzero complementary pairs \(\{x,\bar x\}\) and
		\(\{y,\bar y\}\), its nonzero coefficients in the order
		\((\widetilde q_x,\widetilde q_{\bar x},
		\widetilde q_y,\widetilde q_{\bar y})\) are
		
		\begin{equation}
			\rho(1,\delta,-\gamma,-\delta\gamma^{-1}),
			\qquad \delta,\gamma\in\mu_N.
			\label{eq:odd-dih-q4-two-block}
		\end{equation}
		
		Conversely, every Walsh table in
		\cref{eq:odd-dih-q4-one-block,eq:odd-dih-q4-two-block} satisfies the
		stated contraction condition under all phase and reflection contractions.
	\end{proposition}
	
	\begin{proof}
		Equal kernels give a \(\mu_N\)-ratio on every nonzero complementary
		pair.  If \(p,t\ne0\), \cref{eq:odd-cyc-q4-nonzero-endpoints} puts all nonzero central
		coefficients in \(-p\mu_N\).  For two nonzero central coefficient pairs, the
		\(2\)-by-\(2\) coefficient matrix of the reflected two-port contraction has all
		four entries nonzero, and its only possible proportional alternative
		in \cref{lem:odd-dih-RN} requires a ratio in \(-\mu_N\), not \(\mu_N\).
		The cosets are disjoint for odd \(N\).
		
		With zero endpoints, three nonzero central coefficient pairs would require three
		pairwise ratios in \(-\mu_N\), although the product of two lies in
		\(\mu_N\).  Thus at most two complementary pairs have nonzero coefficients.  Applying the same argument
		to both orientations of the reflected two-port kernels
		\(\mathsf K_s^-\) gives the identical conclusion on the odd Hamming-weight positions.
		
		Substitution gives \cref{eq:odd-dih-q4-one-block,eq:odd-dih-q4-two-block}
		and proves their converse, including zeros.  One block uses
		\(\pi=\{[4]\}\); two blocks form a \(2+2\) partition coset with ratios
		\(-\gamma,-\delta\gamma^{-1}\in-\mu_N\), whose product is \(\delta\).
		Thus both parities lie in \(\mathcal P_N([4])\), with odd parity absorbed
		by the coset offset.
		
		A one-block tensor is a weighted GE\@.  Up to flips and permutation, a
		two-block support is \(\{0000,1111,1100,0011\}\), and
		\cref{eq:odd-dih-q4-two-block} gives
		\[
		\widetilde q_x\widetilde q_{\bar x}
		=\widetilde q_y\widetilde q_{\bar y},
		\]
		
		for the two complementary pairs \(\{x,\bar x\}\) and
		\(\{y,\bar y\}\) in the support.  Hence it is a product of two
		nonsingular weighted equality/disequality
		binaries.  Thus \(\widetilde q\in\cP\), as required by
		\cref{eq:common-Walsh-product-presentation}.
	\end{proof}
	
	\subsubsection{Products of binaries and the six-port case}
	\label{subsec:odd-dih-q6}
	
	Write \(\mathcal B_N(V)\subseteq\mathcal P_N(V)\) for zero together with
	the tensors whose partition consists entirely of pairs.  Every nonzero
	member is a product of nonsingular Walsh equality/disequality binaries,
	with ratio in \(-\mu_N\) on each pair.  For the resulting perfect matching
	\(M\), its support is \(a+L_M\), is complement-closed, and has
	cardinality \(2^{|V|/2}\).
	The contraction formulas in \cref{lem:phase-product-closure} show that
	\(\mathcal B_N\) is itself closed under all residual-group contractions:
	joining different pairs makes one pair, and closing a pair gives a scalar.
	
	\begin{theorem}[Six-port factorization from contractions]
		\label{thm:odd-dih-q6-base}
		Let \(N\ge3\) be odd.  If every nonzero oriented phase or reflection contraction
		of an arity-six signature \(f\) has Walsh tensor in \(\mathcal B_N([4])\),
		then \(f=0\) or \(f\) is a tensor product of three nonsingular binaries.
	\end{theorem}
	
	\begin{proof}
		Suppose \(f\ne0\), and write \(F=\widetilde f\).  By
		\cref{lem:odd-dih-all-arity-parity}, \(F\) has one parity.  Fix a port pair
		and fix \(x_i\oplus x_j=\xi\), with the two corresponding slices \(A,B\), so the normalized actual
		contractions are \(q_s=A-sB\).  Put \(U=\supp(A)\cup\supp(B)\).
		Each nonzero coordinate remains nonzero at at least \(N-1\) phases, while every
		nonzero contraction has four support points.  Thus
		\[
		(N-1)|U|\le4N.
		\]
		Moreover \(U=\bigcup_s\supp(q_s)\) is complement-closed.  For \(N\ge5\)
		it is therefore empty or consists of at most two complement pairs.  For
		\(N=3\), the same conclusion is exactly
		\cref{thm:cert-odd-dih-q6-interface}\textnormal{(i)}.  Part~\textnormal{(ii)}
		of that theorem now applies for every odd \(N\): the
		support of \(F\) is a partition coset of type \((6)\), \((4+2)\), or
		\((2+2+2)\), or belongs to the orbit of
		\begin{equation}
			S_\star=\{000000,000011,001100,110011,111100,111111\}.
			\label{eq:odd-dih-q6-nonproduct-support}
		\end{equation}
		Here and below, bit relabelings describe support coordinates only.  They
		permute the available xor kernels projectively, with possible phase
		inversion, so all contraction arguments cover the full indicated orbits.
		
		A \((6)\) support has a supported-xor contraction with exactly two
		nonzero entries, each with a unique lift.  This contradicts the four-point
		support required by \(\mathcal B_N([4])\).  On a \((4+2)\) support, use
		block variables and denote the four nonzero coefficients by \(c_{uv}\).
		Closing the binary block gives, after normalizing its orientation,
		\[
		(c_{00}-s c_{01})\delta_{0000}
		+(c_{10}-s c_{11})\delta_{1111}.
		\]
		At most two phases are forbidden by cancellation; since \(N\ge3\), some
		contraction again has exactly two nonzero entries.  No tensor factorization of the
		original four coefficients was assumed.  On \(S_\star\), contract ports
		\(3,5\) with \(\mathsf K_s^+\).  Only the two endpoints have equal bits
		there, giving
		\[
		F_{000000}\delta_{0000}-sF_{111111}\delta_{1111},
		\]
		again with two nonzero entries and unique lifts.
		
		Only the \((2+2+2)\) support remains.  Its eight coefficients
		\(c_{uvw}\) are nonzero.  Closing any one of its binary blocks leaves a
		\(2\)-by-\(2\) table whose determinant vanishes at every \(s\in\mu_N\):
		a nonzero contraction has the same four-point matching support, and a zero contraction
		also has zero determinant.  For example,
		\[
		\det\bigl(c_{uv0}-s c_{uv1}\bigr)_{u,v\in\Ftwo}=0.
		\]
		This polynomial has degree at most two and at least three roots.  Its
		constant and quadratic coefficients show that the two endpoint faces have
		rank one.  Doing this for all three axes makes every axis face minor zero.
		Since all eight entries are nonzero, the ratios along each axis are
		independent of the other two bits, so \(c_{uvw}=c r_1^u r_2^v r_3^w\).
		This is the required three-binary factorization, which pulls back through
		the analytic Walsh change of coordinates.
	\end{proof}
	
	\subsubsection{Matching pencils}
	\label{subsec:odd-dih-support-circuits}
	
	\begin{lemma}[Three matching products on a line]
		\label{lem:odd-dih-three-coset-circuit}
		Let three pairwise nonproportional tensor products of nonsingular diagonal
		or anti-diagonal binaries on the same \(m\) ports satisfy a linear relation
		with all coefficients nonzero.  Either their supports coincide and
		the three tensors lie in one
		Segre ruling, or their matchings differ only on four ports, where the three
		different pairings occur.  In the latter case all binary factors outside
		those four ports are common up to nonzero scalars.
	\end{lemma}
	
	\begin{proof}
		Let the supports be \(S_0,S_1,S_2\), matching cosets of dimension \(m/2\).
		No point can belong to just one support.  If two supports coincide, the
		third is contained in them and has the same cardinality, so all coincide.
		On their common block coordinates, every matrix flattening minor has degree at
		most two along the line.  Three rank-one points make these minors vanish
		identically.  A line of simple tensors varies in at most one factor: if
		two factors of two of its points were nonproportional, flattening across
		one of them would give rank two at a generic third point.  This is one
		Segre ruling.
		
		Otherwise each \(S_i\) is the union of its intersections with the other
		two supports.  Two proper affine subspaces cover an affine space only when
		they are disjoint parallel hyperplanes: each has at most half its points,
		so both must attain that bound.  Thus the pairwise intersections are
		hyperplanes and the triple intersection is empty.  The union of two
		perfect matchings is a disjoint union of alternating even cycles, counting
		a common edge as a two-cycle.  If there are \(k=m/2\) edges and \(b\)
		components, their common direction code has dimension \(b\).  Here
		\(b=k-1\), so exactly one component is a four-cycle and all other edges
		coincide.  The offsets on those common edges agree because the supports
		intersect.  After removing their free binary coordinates, \(S_2\) is the
		four-point symmetric difference of the other two local supports; its
		matching is the third pairing on the same four ports.  Finally, on each
		pairwise intersection only two tensors contribute to the dependence.
		Fixing one four-port assignment with nonzero coefficient leaves the full external matching
		cube, on which their products are proportional.  Uniqueness of the factors
		on disjoint pairs gives the asserted common factors.
	\end{proof}
	
	\begin{proposition}[Phase-independent matching and common factors]
		\label{prop:odd-dih-support-circuit-rigidity}
		Let \(N\ge3\) be odd, and let \(f\) have even arity \(n\ge8\).  Suppose
		every nonzero result of one or more successive two-port contractions using
		the kernels in
		\cref{eq:odd-dih-kernels} has Walsh tensor in \(\mathcal B_N\).  Then each
		two-port pencil can be written as \(q_s=H\otimes h_s\), where \(H\)
		is a fixed product of nonsingular binaries on disjoint pairs and \(h_s\)
		is a binary tensor on one fixed remaining port pair.  Zero members and proportional
		pencils are included.
	\end{proposition}
	
	\begin{proof}
		Write a pencil as \(q_s=A-sB\).  A zero member makes all other members
		proportional, which settles that case.  If the pencil is not projectively
		constant, any three distinct phases give a linear dependence with all three
		coefficients nonzero,
		so \cref{lem:odd-dih-three-coset-circuit} applies to these members of
		\(\mathcal B_N\).  In its common-support
		case, the quadratic flattening minors vanish identically, giving the fixed
		matching ruling for the whole pencil.
		
		In the second case of \cref{lem:odd-dih-three-coset-circuit}, the matchings
		differ on four ports and all other binary factors are common.  Close every such common pair, on the original gadget,
		with its supported-xor kernel at phase \(1\).  For a binary ratio
		\(r\in-\mu_N\), this gives a nonzero multiple of \(1-r\), nonzero because
		\(N\) is odd.  These operations multiply the three contractions by the same
		nonzero scalar.  At least one pair is closed since \(n\ge8\).  Before
		closing the original pencil ports, the resulting six-port tensor is a
		nonzero result of contractions of \(f\), hence is itself a binary matching
		product.  Its contractions cannot have three different matchings: closing ports
		of a fixed matching deletes an edge or merges two edges, independently of
		the phase.  This contradiction proves that the matching is independent of phase.
	\end{proof}
	
	\subsubsection{A common binary factor suffices}
	\label{subsec:odd-dih-cross-sector}
	
	\begin{lemma}[A binary factor from contractions]
		\label{lem:odd-dih-common-binary-factor}
		Under the hypotheses of \cref{prop:odd-dih-support-circuit-rigidity}, every
		nonzero \(f\) is tensor-decomposable.
	\end{lemma}
	
	\begin{proof}
		Put \(F=\widetilde f\), fix ports \(i,j\), and for each
		\(\xi\in\{0,1\}\) consider the restriction to \(x_i\oplus x_j=\xi\), viewed as an
		\(n\)-ary tensor that is zero outside that constraint.
		Its two slices are \(A=F_{ij}^{0,\xi}\), \(B=F_{ij}^{1,1-\xi}\).
		Exactly one of \(A,B\) cannot be nonzero.  Indeed, if \(A\ne0\) and \(B=0\), take a
		supported word of \(A\), close any two external ports with its supported
		xor, and choose a phase avoiding cancellation of the projected word.  The
		resulting nonzero contraction has complement-closed support.  Its
		complementary word would require an assignment in \(B\), a contradiction.
		The case \(A=0\), \(B\ne0\) is identical.
		
		For each nonzero restriction, \cref{prop:odd-dih-support-circuit-rigidity} writes
		the pencil as \(H\otimes h_s\), where \(H\) is a fixed product of external
		binaries and
		\[
		h_s=(a-sb)\delta_e+(c-sd)\delta_{\bar e}
		\]
		on one external pair.  Projectively constant pencils are included by
		choosing any pair.  Every nonzero \(h_s\) has ratio in \(-\mu_N\), so
		\cref{lem:odd-dih-RN} applies.  Restoring the ports \(i,j\), its
		proportional alternative gives a binary block \(\{i,j\}\); either
		monomial alternative gives one four-port GE block containing \(i,j\).
		All other blocks are nonsingular binaries.  Both slices \(A,B\) are
		nonzero, so the restored block has both coefficients nonzero.  This is a
		factorization in Walsh coordinates of the restriction; it does not assert
		that either this restriction or either slice is realizable by a gadget.
		
		If only one of the two restrictions is nonzero, this factorization of \(F\) already has an
		external binary factor, since \(n\ge8\).  Otherwise denote the restrictions for \(\xi=0,1\) by \(U,V\).  We use the following observation repeatedly.  Close
		an external pair with an xor supported by both restrictions, and choose a phase
		at which both resulting contractions are nonzero.  Such a phase exists: a contraction of
		either displayed block product with a supported xor has unique lifts, except when an entire
		binary block is closed, in which case only one phase can cancel its
		scalar.  The nonzero contraction \(G\) is a binary matching product with
		both \(ij\) xors present.  Its two \(ij\)-xor restrictions have:
		\begin{enumerate}[label=\textnormal{(O\arabic*)}]
			\item the same four-port block containing \(i,j\);
			\item the same external matching and the same offset on each external pair;
			\item proportional full binary factors on each external pair, including
			equality of their ordered coefficient ratios.
		\end{enumerate}
		Indeed, \(i,j\) belong to different matching edges in \(G\).  Either xor
		restriction joins those edges into one four-port block and leaves every
		other binary factor unchanged.  In particular, this observation identifies
		coefficients as well as supports.
		
		First, neither \(U\) nor \(V\) can have \(\{i,j\}\) as a binary block.
		If \(U\) did, close ports in two distinct external binary blocks of
		\(U\).  Both xors occur in \(U\), so one is also supported by \(V\).
		The observation would require a four-port block through \(i,j\), whereas
		this contraction leaves their binary block unchanged.
		
		Write the four-port blocks of \(U,V\) as
		\(K=\{i,j,a,b\}\), \(L=\{i,j,c,d\}\).  Close \(a,b\) with its fixed
		xor in \(U\).  That restriction now has \(\{i,j\}\) as a binary block, so the
		observation implies that \(V\) does not support this xor.  Therefore
		\(a,b\) lie in the same block of \(V\), with the opposite xor.  Either
		\(L=K\), or \(\{a,b\}\) is an external binary block of \(V\).  By
		symmetry, in the latter case \(\{c,d\}\) is an external binary block of
		\(U\), and the two pairs are disjoint.  Choose a port \(e\) outside
		\(\{i,j,a,b,c,d\}\), possible since \(n\ge8\), and close \(a,e\).
		Both restrictions have unique lifts for either xor.  The four-port block through
		\(i,j\) in the contracted \(U\) contains \(b\), whereas that in the
		contracted \(V\) remains \(\{i,j,c,d\}\).  This contradicts the
		observation.  Thus \(K=L\).
		
		Let \(M_U,M_V\) be the external matchings on \(E=[n]\setminus K\).
		For \(e\in E\), closing \(a,e\) makes the four-port block through
		\(i,j\) equal to \(\{i,j,b,M_U(e)\}\) in one restriction and
		\(\{i,j,b,M_V(e)\}\) in the other.  The observation gives
		\(M_U(e)=M_V(e)\), so the external matchings coincide.  There are at least
		two external pairs.  For any one of them, \(R\), choose \(e\) on another
		pair and close \(a,e\).  The pair \(R\) remains external.  The observation
		now identifies its full binary factor in the two restrictions, hence in
		\(U,V\), since the contraction avoids \(R\).  Therefore \(U+V=F\) has
		this common factor.  Invertible Walsh coordinates preserve tensor
		factorization, proving the lemma.
	\end{proof}
	
	\begin{proof}[Proof of \cref{thm:odd-dihedral}]
		Apply \cref{lem:height-one-card-saturation}.  Its first two outcomes are
		the first two conclusions of the theorem.  In outcome~3, work in the
		equivalent set \(\Gamma=\cG^{\mathrm{sat}}\), with the same matrix and
		group.  If a nonbinary prime remains, choose one, \(f\), of minimum arity
		\(n\).  Odd and unary factors have already been handled by their respective
		complexity theorems.  Every positive-arity prime factor of a two-port contraction of \(f\)
		has smaller arity, and hence is binary by minimality.  Its transfer is in
		\(D_{2N}\), so every nonzero two-port contraction is in \(\mathcal B_N\) in Walsh
		coordinates.  Closure of \(\mathcal B_N\) gives the same assertion for
		all further contractions involving at least one port pair.
		
		If \(n=4\), \cref{prop:odd-dih-q4} gives either a four-port generalized
		equality, to which \cref{lem:walsh-ge-terminal} applies, or a product of
		two binaries, contradicting primeness.  If \(n=6\),
		\cref{thm:odd-dih-q6-base} contradicts primeness.  If \(n\ge8\),
		\cref{lem:odd-dih-common-binary-factor} gives the same contradiction.
		Thus either hardness or tractability has been proved, or every tensor-prime
		in \(\Gamma\) is
		binary.  In the latter case one Walsh basis puts all signatures in \(\cP\),
		and \(P^{\mathsf T}XP=2Z\in\cP\), giving outcome~3 and the stated common \(K\)-presentation.
	\end{proof}
	
	\subsection{The odd-cyclic case}
	
	\begin{theorem}[The odd-cyclic case]
		\label{thm:odd-cyclic}
		Let \((A,\cG)\) be a stage of the sequence in
		\cref{lem:normalization-ledger} at which \(\cG\) is finite, preprocessed,
		algebraic, and factor-closed and contains the transported height-one
		mountain.  Let \(N\ge3\) be odd, and suppose every nonzero directly
		realizable binary is nonsingular with finite projective transfer order, and
		\[
		P^{-1}G(\cG)P
		=C_N=\{[\diag(1,s)]:s\in\mu_N\}.
		\]
		Then at least one of the following holds.
		\begin{enumerate}
			\item \(\KHolant(\cG)\) is \(\SharpP\)-hard or \(\Tract(K\cG)\) holds.
			\item There is a later stage \((A^+,\cG^+)\) of the same sequence at which
			\(\cG^+\) is finite, preprocessed, algebraic, and factor-closed and
			still contains the correspondingly transported height-one mountain,
			such that
			\[
			\KHolant(\cG^+)\equivT\KHolant(\cG),
			\]
			every nonzero directly realizable ordered binary over \(\cG^+\) is
			nonsingular with finite projective transfer order, and either
			\[
			[A]^{-1}G(\cG)[A]
			\subsetneq
			[A^+]^{-1}G(\cG^+)[A^+]
			\]
			or, with both values defined, \(\nu(\cG^+)<\nu(\cG)\).
			In either case the full update conditions of
			\cref{def:height-one-outer-successor} hold.
			\item There is an equivalent finite preprocessed algebraic factor-closed
			set \(\cG^{\mathrm{sat}}\supseteq\cG\), in the same
			coordinates and with the same matrix \(A\), for which the fixed
			Walsh matrix \(P\) satisfies
			\[
			(P^{-1})^{\otimes\arity(f)}f\in\cP
			\quad(f\in\cG^{\mathrm{sat}}),
			\qquad
			P^{\mathsf T}XP=2Z\in\cP.
			\]
			Hence \(K\cG\) is \(\cP\)-transformable.
		\end{enumerate}
	\end{theorem}
	
	\subsubsection{Phase-independent matching}
	\label{subsec:odd-cyc-pencil}
	
	Throughout this odd-cyclic proof, \(\widetilde f\) denotes the fixed
	Walsh coordinates, and
	\(P^{-1}G(\Gamma)P=C_N=\{[D_s]:s\in\mu_N\}\), with one direct representative fixed for
	each class and \(N\ge3\) odd.
	
	For the conditional analysis below, work in outcome~3 of
	\cref{lem:height-one-card-saturation}, retaining its finite, preprocessed,
	algebraic, factor-closed set \(\Gamma\), the matrix \(A\), and the
	transported height-one mountain.  Odd-arity tensor-prime signatures have
	already been handled by \cref{thm:external-odd}.  Choose a nonbinary
	tensor-prime \(f\in\Gamma\) of minimum arity \(n\), which is even.
	Attaching the chosen binary representative to distinct ports \(i,j\) and
	closing through the native edges gives, up to a known nonzero
	representative scalar, the Walsh kernel
	\[
	\mathsf K_s=\diag(1,-s),
	\qquad s\in\mu_N.
	\]
	After dividing out that recorded scalar, the contraction has Walsh coefficients
	\begin{equation}
		C_{ij,s}(y)=A(y)-sB(y),
		\qquad
		A(y)=\widetilde f(00,y),
		\quad B(y)=\widetilde f(11,y).
		\label{eq:odd-cyc-card-pencil}
	\end{equation}
	By \cref{lem:height-one-card-saturation}\textnormal{(3)}, each indicated
	nonzero contraction and all its positive-arity tensor-prime factors belong
	to \(\Gamma\).  A nonbinary factor would have arity at most \(n-2\),
	contradicting minimality.  Every binary factor is nonsingular and its Walsh
	transfer class belongs to \(C_N\).  Thus
	condition~\ref{cond:height-one-contractions} holds for \(f\); each nonzero
	contraction is a product of such binaries on a perfect matching.  A transfer \(D_t\) has
	Walsh tensor
	\[
	D_tZ=\diag(1,-t).
	\]
	For a perfect matching \(M\) of the residual ports, put
	\begin{equation}
		E_M=\{y\in\Ftwo^{[n]\setminus\{i,j\}}:
		y_a=y_b\text{ for every }\{a,b\}\in M\}.
		\label{eq:odd-cyc-matching-code}
	\end{equation}
	Every nonzero contraction has support \(E_M\), with a nowhere-zero rank-one table
	on the matching variables.  By
	\cref{lem:odd-dih-all-arity-parity}, \(\widetilde f\) has even Walsh
	parity.
	
	\begin{lemma}[Phase-independent matching]
		\label{lem:odd-cyc-one-matching}
		Let \(N\ge3\) be odd, and let \(f\in\Gamma\) be a nonbinary
		tensor-prime signature of minimum arity under the assumptions of this
		subsection, including condition~\ref{cond:height-one-contractions} with
		\(P^{-1}G(\Gamma)P=C_N\).  Define \(C_{ij,s}\) and \(E_M\)
		by \cref{eq:odd-cyc-card-pencil,eq:odd-cyc-matching-code}.  For each pair
		\(i,j\), either every contraction vanishes, or \(C_{ij,s}=0\) for at most one
		phase \(s\in\mu_N\) and there is a fixed perfect matching \(M_{ij}\)
		such that
		\begin{equation*}
			C_{ij,s}\ne0
			\quad\Longrightarrow\quad
			\supp(C_{ij,s})=E_{M_{ij}}
			\qquad(s\in\mu_N).
		\end{equation*}
		This includes the boundary case \(N=3\).
	\end{lemma}
	
	\begin{proof}
		If \(A-sB\) and \(A-tB\) vanish at distinct phases, then \(A=B=0\),
		so the first alternative holds.  Otherwise at most one phase gives a zero
		contraction.  Any two distinct nonzero contractions satisfy
		\begin{equation}
			\supp(A)\cup\supp(B)
			=\supp(A-sB)\cup\supp(A-tB),
			\label{eq:odd-cyc-two-card-union}
		\end{equation}
		because a scalar affine polynomial cannot vanish at two phases unless both
		coefficients vanish.
		
		If no contraction vanishes, choose three phases.  By
		\cref{eq:odd-cyc-two-card-union}, the third matching code lies in the union
		of the first two.  A vector space contained in \(U\cup V\) lies in one of
		them (otherwise sum vectors from its two exclusive parts), and all matching
		codes have the same dimension.  Thus the third code equals one of the first
		two.  Two phases with one code and a third with another would make a word in
		the latter difference vanish at two phases, contradicting the affine-pencil
		identity.  Hence all codes coincide.
		
		If \(A-s_0B=0\), then every other contraction is \((s_0-s)B\), and the conclusion
		is immediate.
	\end{proof}
	
	\subsubsection{Pairs with every phase contraction zero}
	\label{subsec:odd-cyc-zero-graph}
	
	Define
	\begin{equation*}
		Z_f^{\mathrm{cyc}}
		=
		\{\{i,j\}:C_{ij,s}=0\text{ for every }s\in\mu_N\}.
	\end{equation*}
	Two phases in \cref{eq:odd-cyc-card-pencil} show that an edge is exactly
	the support equation
	\begin{equation*}
		\widetilde f(x)\ne0
		\quad\Longrightarrow\quad
		x_i+x_j=1.
	\end{equation*}
	It follows that every connected component of
	\(Z_f^{\mathrm{cyc}}\) is complete bipartite.  An odd path forces unequal
	endpoint bits and hence supplies the corresponding edge; an even path
	forces equal bits and cannot also be a zero edge.
	
	\begin{lemma}[Every port pair has a nonzero phase contraction above arity four]
		\label{lem:odd-cyc-zero-graph-empty}
		Under the assumptions of \cref{lem:odd-cyc-one-matching}, if
		\(\arity(f)=n\ge6\), then
		\[
		Z_f^{\mathrm{cyc}}=\varnothing.
		\]
	\end{lemma}
	
	\begin{proof}
		Each color class has size at most two: if one contained \(p,q,r\) and the
		other \(b\), a nonzero contraction on the nonzero pair \(p,q\) would have an
		all-zero residual coefficient while retaining the forced disequality
		\(r\ne b\), impossible for a nonsingular matching product.
		
		An isolated \(\{1\}\sqcup\{2\}\) is impossible.  For every \(k\ge3\),
		the all-zero word of a nonzero \((1,k)\)-contraction has the unique lift
		\(x_1=x_k=1\), whence
		\begin{equation*}
			\widetilde f(e_1+e_k)\ne0
			\qquad(k=3,\ldots,n).
		\end{equation*}
		A nonzero \((2,3)\)-contraction then has each \(e_1+e_k\), \(k\ge4\), by the
		unique lift \(x_2=x_3=0\).  Its support matching would contain every
		\(\{1,k\}\), impossible.
		
		Thus only \(2+1\) and \(2+2\) remain.  Label the size-two color
		\(\{1,3\}\) and the other \(\{2\}\) or \(\{2,4\}\).  In a \((1,3)\)
		contraction the forced equations give the unique lift
		\[
		x_1=x_3=1-x_2.
		\]
		Put
		\[
		h(x_2,\mathbf y)=\widetilde f(1-x_2,x_2,1-x_2,\mathbf y),
		\]
		where \(\mathbf y\) denotes the remaining coordinates; in the \(2+2\)
		case, also impose the support equation \(x_4=x_2\).
		Then
		\begin{equation*}
			C_{13,s}(x_2,\mathbf y)=u_s(x_2)h(x_2,\mathbf y),
			\qquad
			u_s(0)=-s,\quad u_s(1)=1.
		\end{equation*}
		Every contraction is nonzero and equals a nowhere-zero unary \(u_s\) times \(h\).
		Dividing analytically by \(u_s\) only modifies the matching factor incident
		with port \(2\), so \(h\) factors over that matching.  Reversing the unique
		lift gives a four-port GE factor on \(\{1,2,3,4\}\) in the \(2+2\) case,
		or on \(\{1,2,3,k\}\) in the \(2+1\) case, tensor the remaining binaries.
		For \(n\ge6\) this contradicts tensor-primeness.
	\end{proof}
	
	\subsubsection{Weight-two coefficients and the arity bound}
	\label{subsec:odd-cyc-weight-two}
	
	\begin{lemma}[Rank-one diagonal restrictions]
		\label{lem:diagonal-slice-integration}
		Let \(r\ge3\), and let \(h:\Ftwo^r\to\mathbb C^\times\) have full support and
		\(h(0)=1\).  Suppose that, for every \(i\ne j\), the restriction
		\(h|_{t_i=t_j}\), viewed as a function of the common bit and the other
		\(r-2\) bits, has tensor rank one.  If \(r\ge4\), then
		\[
		h(t)=\prod_{i=1}^r\rho_i^{t_i}
		\qquad(\rho_i\ne0).
		\]
		If \(r=3\), there are \(\rho_1,\rho_2,\rho_3,\kappa\ne0\) such that
		\begin{equation}
			h(e_i)=\rho_i,\qquad
			h(e_i+e_j)=\kappa\rho_i\rho_j,\qquad
			h(111)=\kappa\rho_1\rho_2\rho_3.
			\label{eq:diagonal-slice-q3-table}
		\end{equation}
	\end{lemma}
	
	\begin{proof}
		This is \cite[Lemma~7.8]{Liu2026DEQ2Holant}.
	\end{proof}
	
	\begin{theorem}[The minimum nonbinary arity is four in the cyclic case]
		\label{thm:odd-cyc-higher-collapse}
		Let \(\Gamma\) be a finite preprocessed algebraic factor-closed signature
		set containing the transported height-one mountain.  Suppose every nonzero
		directly realizable binary is nonsingular and has finite projective transfer
		order, and, in the fixed
		Walsh coordinates,
		\[
		P^{-1}G(\Gamma)P
		=C_N=\{[\diag(1,s)]:s\in\mu_N\},
		\qquad N\ge3\text{ odd}.
		\]
		Let \(f\in\Gamma\) be an even-arity nonbinary tensor-prime signature of
		minimum arity among the nonbinary tensor-prime signatures in \(\Gamma\).
		For every ordered pair of distinct ports of \(f\), every fixed direct
		representative of a class in \(G(\Gamma)\), and either kernel orientation,
		assume that the resulting two-port contraction is zero or all of its positive-arity
		tensor-prime factors are nonsingular binaries \(B\) satisfying
		\[
		[P^{-1}(BX)P]\in C_N.
		\]
		Then \(\arity(f)\le4\), and hence \(\arity(f)=4\).
	\end{theorem}
	
	\begin{proof}
		The contraction hypothesis is condition~\ref{cond:height-one-contractions}
		for the stated group, so the two preceding lemmas apply.
		Assume that \(f\) has arity \(n\ge6\).  By
		\cref{lem:odd-cyc-zero-graph-empty}, every pair has a nonzero phase contraction.
		For \(S\subseteq[n]\), write \(a_S=\widetilde f(\one_S)\).
		We abbreviate \(a_{i_1\cdots i_k}=a_{\{i_1,\ldots,i_k\}}\).
		
		\emph{Step 1: the constant coefficient is nonzero.}
		If \(a_\varnothing=0\), the all-zero coefficient
		\(C_{ij,s}(0)=-sa_{ij}\) of a nonzero matching contraction forces
		\(a_{ij}\ne0\) for every pair.  Hence all contractions are nonzero.  For disjoint
		pairs \(ij,pq\),
		\[
		C_{ij,s}(e_p+e_q)=a_{pq}-sa_{ijpq}
		\]
		is nonzero at some phase.  The common matching from
		\cref{lem:odd-cyc-one-matching} would therefore contain every edge on its
		at least four residual ports, impossible.  Normalize henceforth
		\(a_\varnothing=1\), and put
		\[
		W_f=\{\{p,q\}:a_{pq}\ne0\}.
		\]
		
		\emph{Step 2: a nonempty \(W_f\) factors \(f\).}
		For deleted \(i,j\) and a disjoint edge \(pq\in W_f\), the coefficients
		\(1-sa_{ij}\) and \(a_{pq}-sa_{ijpq}\) exclude at most one phase each.
		Thus some contraction contains \(e_p+e_q\), and its common support matching
		contains \(pq\).  Deleting two ports outside any two adjacent
		\(W_f\)-edges would put both edges in one matching, so \(W_f\) is a
		matching.
		
		Moreover, every \(pq\in W_f\) is a global equality constraint.  Indeed, if
		some supported \(z\) had \(z_p\ne z_q\), choose equal coordinates \(i,j\)
		outside \(p,q\), possible by even parity and \(n\ge6\).  The residual word
		violates the \(pq\)-edge in every nonzero \((i,j)\)-contraction and therefore
		vanishes at all phases.  Two phases of
		\cref{eq:odd-cyc-card-pencil} then make both lifts vanish, including \(z\).
		Consequently
		\begin{equation}
			pq\in W_f\quad\Longrightarrow\quad
			\supp(\widetilde f)\subseteq\{z:z_p=z_q\}.
			\label{eq:odd-cyc-global-equality}
		\end{equation}
		
		Fix \(E_1=\{u,v\}\in W_f\) and \(R=[n]\setminus E_1\).  In an
		\((u,k)\)-contraction, the all-zero word has the unique lift
		\(x_u=x_k=x_v=0\); hence, for \(p,q\in R\setminus\{k\}\), its coefficient
		at \(e_p+e_q\) is exactly \(a_{pq}\).  Its matching therefore consists of
		the \(W_f\)-edges on \(R\setminus\{k\}\) and one edge incident with \(v\).
		Since that last edge cannot cover all \(n-2\ge4\) ports, the induced
		matching on \(R\) is nonempty.  Choose \(k\) on one of its edges, with
		partner \(\ell\).  In the \((u,k)\)-contraction, \(\ell\) has no incident
		\(W_f\)-edge left, so it must be the partner of \(v\).  Any vertex
		unmatched in \(W_f|_R\) would also have to be matched to \(v\), which is
		impossible.  Thus \(W_f=\{E_1,\ldots,E_r\}\) is perfect, where \(r=n/2\).
		
		Equation \eqref{eq:odd-cyc-global-equality} identifies \(\widetilde f\)
		with a coefficient function \(h(t_1,\ldots,t_r)\) on the common bits of the
		\(E_k\)'s.  Contracting one port from \(E_i\) and one from \(E_j\) imposes
		\(t_i=t_j\), and the remaining partners determine the deleted bits uniquely.
		The contraction is nonzero since its all-zero coefficient is one.  Its support is
		contained in the matching code obtained by merging \(E_i,E_j\).  Both this
		code and the contraction's matching support have \(2^{r-1}\) points, so they
		coincide.  Removing the nowhere-zero unary phase on the merged variable
		therefore gives a full-support rank-one diagonal slice of \(h\).  Every assignment in \(\Ftwo^r\) has two equal coordinates, so these
		restrictions cover \(\Ftwo^r\) and \(h\) has full support.  After \(h(0)=1\),
		\cref{lem:diagonal-slice-integration} factors \(h\) when \(r\ge4\).  If
		\(r=3\), use \cref{eq:diagonal-slice-q3-table}.  Let \(H_i(s)\) be the
		two-bit coefficient function obtained by closing \(E_i\) at phase \(s\).  For
		\(\{i,j,k\}=[3]\),
		\begin{equation*}
			\det H_i(s)=\rho_j\rho_k(\kappa-1)
			(1-\kappa s^2\rho_i^2).
		\end{equation*}
		If \(\kappa\ne1\), the contraction is nonzero and hence rank one for every phase,
		whereas \(s\mapsto s^2\) permutes the odd-order group \(\mu_N\); the last
		factor cannot vanish at two phases.  Thus \(\kappa=1\), and the six-port
		table also factors.  Both cases contradict tensor-primeness.
		
		\emph{Step 3: an empty \(W_f\) is impossible.}
		Fix \(i,j\).  Every contraction now has constant coefficient one and is nonzero.
		Choose two edges \(e,e'\) of its common residual matching.  Full matching
		support gives
		\(a_{\{i,j\}\cup e}a_{\{i,j\}\cup e'}\ne0\), while its two-bit rank-one
		identity is
		\begin{equation*}
			a_{e\cup e'}-s a_{\{i,j\}\cup e\cup e'}
			=s^2a_{\{i,j\}\cup e}a_{\{i,j\}\cup e'}
			\qquad(s\in\mu_N).
		\end{equation*}
		This nonzero quadratic cannot vanish at the \(N\ge3\) distinct phases.
		The contradiction completes the proof.
	\end{proof}
	
	\subsubsection{The quaternary cyclic case}
	\label{subsec:odd-cyc-quaternary}
	
	We now treat arity four without assuming that the quaternary is a mountain.
	The following lemma gives the output obtained by attaching binary
	representatives to the four ports.
	
	\begin{lemma}[Attaching cyclic binary representatives to four ports]
		\label{lem:odd-cyc-four-port-dressing}
		Let \(N\ge3\), let \(\Lambda\) be a finite algebraic signature set in
		\(K\)-coordinates, and, for every \(s\in\mu_N\), fix an actual ordered
		binary representative \(C_s\) over \(\Lambda\), in a specified orientation,
		whose Walsh matrix is
		\[
		\widetilde C_s=\kappa_sD_sZ,
		\qquad \kappa_s\in\AlgNums^\times.
		\]
		Let \(q\in\Lambda\) be quaternary.
		For \(r_1,r_2,r_3,r_4\in\mu_N\), attach \(C_{r_i}\) to port \(i\) of
		\(q\) through one native \(X\)-edge and leave the other port of each
		representative exposed.  After division by the recorded nonzero scalar
		\(2^4\prod_{i=1}^4\kappa_{r_i}\), the resulting actual output \(h\) satisfies
		\begin{equation*}
			\widetilde h_y
			=\widetilde q_y\prod_{i=1}^4r_i^{y_i}.
		\end{equation*}
		Consider the map
		\begin{equation*}
			(r_1,r_2,r_3,r_4)
			\longmapsto
			(R,r_{12},r_{13},r_{14}),
			\qquad
			R=\prod_{i=1}^4r_i,
			\quad r_I=\prod_{i\in I}r_i,
		\end{equation*}
		Write \((\mu_N)^2:=\{z^2:z\in\mu_N\}\).  A target
		\((R,u,v,w)\in\mu_N^4\) occurs exactly when
		\begin{equation}
			\frac{uvw}{R}\in(\mu_N)^2.
			\label{eq:cyclic-dressing-image}
		\end{equation}
		For odd \(N\), this map is bijective.
	\end{lemma}
	
	\begin{proof}
		This is \cite[Lemma~7.3]{Liu2026DEQ2Holant}.
	\end{proof}
	
	\begin{theorem}[The quaternary odd-cyclic case]
		\label{thm:odd-cyc-quaternary-closure}
		Let \(\Gamma\) be a finite preprocessed algebraic factor-closed signature
		set containing the transported height-one mountain.  Suppose every nonzero
		directly realizable ordered binary is nonsingular and has finite projective
		transfer order, and, in the fixed Walsh coordinates,
		\[
		P^{-1}G(\Gamma)P
		=C_N=\{[\diag(1,s)]:s\in\mu_N\},
		\qquad N\ge3\text{ odd}.
		\]
		Let \(0\ne q\in\Gamma\) be quaternary.  For every ordered pair of distinct
		ports of \(q\), every fixed direct representative of a class in
		\(G(\Gamma)\), and either kernel orientation, assume that the resulting
		two-port contraction is zero or all of its positive-arity tensor-prime factors are
		nonsingular binaries \(B\) with
		\[
		[P^{-1}(BX)P]\in C_N.
		\]
		Then at least one of the following holds.
		\begin{enumerate}
			\item \(\KHolant(\Gamma)\) is \(\SharpP\)-hard or
			\(\Tract(K\Gamma)\) holds.
			\item Up to a permutation of its ports,
			\[
			q\doteq a\otimes b,
			\]
			where \(a,b\) are nonzero positive-arity signatures of arity strictly
			below four, and
			\[
			\KHolant(\Gamma,a,b)\leT\KHolant(\Gamma).
			\]
		\end{enumerate}
	\end{theorem}
	
	\begin{proof}
		The stated contractions satisfy
		condition~\ref{cond:height-one-contractions} for the cyclic group.
		In the Walsh tensor \(\widetilde q\), put
		\[
		p=\widetilde q_{0000},\qquad
		t=\widetilde q_{1111},\qquad
		x_I=\widetilde q_{\one_I}\quad(|I|=2).
		\]
		For each such \(I\), put \(B_I=(x_I,x_{I^c})\), so
		\(B_{I^c}\) has the same two coefficients in reverse order.
		By \cref{lem:odd-dih-all-arity-parity}, \(\widetilde q\) is even; the
		shared family \cref{lem:odd-cyc-q4-family} excludes exactly one nonzero
		Walsh endpoint.  Suppose first that \(p,t\ne0\), and put \(\lambda=t/p\).
		For each \(B_I\ne(0,0)\), it gives
		\((x_I,x_{I^c})=(-p\lambda\alpha_I^{-1},-p\alpha_I)\).
		For each zero pair choose \(\alpha_I\in\mu_N\) arbitrarily, using one
		index from each complementary pair as in \cref{lem:odd-cyc-q4-family}.
		By \cref{lem:odd-cyc-four-port-dressing}, prescribe
		\(R=\lambda^{-1}\) and \(r_I=\alpha_I\lambda^{-1}\).  Let \(h\) be the
		normalized representative of the gadget output, obtained by division
		by the recorded nonzero gadget scalar.  Its Walsh coefficient array
		\(\widetilde h\) has endpoints
		\((p,p)\), value \((-p,-p)\) on each complementary pair with \(B_I\ne(0,0)\),
		and zero on every pair with \(B_I=(0,0)\).
		
		Here the eight even-parity Walsh positions form the endpoint pair and
		three central pairs \(\{\one_I,\one_{I^c}\}\), where \(|I|=2\).
		In the inverse Walsh formula of \cref{eq:Walsh-coordinates}, a pair
		\(\{y,\bar y\}\) with common coefficient \(b\) contributes
		\[
		b\bigl((-1)^{x\cdot y}+(-1)^{x\cdot\bar y}\bigr).
		\]
		This vanishes when \(\wt(x)\) is odd.  At either endpoint all signs are
		positive.  At \(x=\one_J\), \(|J|=2\), the endpoint pair and the
		pair \(\{\one_J,\one_{J^c}\}\) have positive signs, and the other two
		central pairs have negative signs.  Thus, if exactly \(k\) of the three pairs \(B_I\) are
		nonzero,
		\[
		\begin{aligned}
			h(\zero^4)=h(\one^4)&=2p(1-k),\\
			h(\one_J)&=
			\begin{cases}
				2p(k-1),&\text{if \(B_J\ne(0,0)\)},\\
				2p(k+1),&\text{if \(B_J=(0,0)\)}.
			\end{cases}
		\end{aligned}
		\]
		This computes the K-coordinate entries of the same normalized output; Fourier
		inversion is not an additional gadget operation.  Thus \(k=1\) gives a
		nonzero multiple of \(X\otimes X\), up to a port permutation, whose
		factorization pulls back through the invertible binary attachments just
		constructed.  For
		\(k=0\), the original Walsh tensor is already a higher-arity GE, so
		\cref{lem:walsh-ge-terminal} gives outcome~1.  For \(k=2,3\), the output in
		native coordinates is even and has two nonzero endpoints.  Adjoin this actual
		quaternary and apply \cref{cor:eight-vertex-complete-closure}.  Its
		conclusion transfers to \(\Gamma\) by the gadget
		reduction and downward inheritance, again giving outcome~1.
		
		It remains that \(p=t=0\).  If just one complementary pair has nonzero coefficients,
		\cref{lem:walsh-ge-terminal} gives outcome~1 directly; no reflection is
		needed.  Otherwise label the three complementary pairs
		\(\{\one_{I_\ell},\one_{I_\ell^c}\}\), \(\ell=1,2,3\).
		Write each nonzero pair \(B_I\) as
		\(\rho_I(1,\alpha_I)\), choose the unique \(\beta_I\in\mu_N\) with
		\(\beta_I^2=\alpha_I\); for a zero pair choose \(\beta_I\in\mu_N\)
		arbitrarily.  Apply \cref{lem:odd-cyc-four-port-dressing} with
		\(R=1,r_I=\beta_I\).  Again divide
		the actual output by its recorded nonzero gadget scalar and denote the
		normalized tensor by \(h\).  Its nonzero Walsh coefficient pairs
		become \((\sigma_I,\sigma_I)\), where
		\(\sigma_I=\rho_I\beta_I\); set \(\sigma_I=0\) when \(B_I=(0,0)\), and write
		\(\sigma_\ell=\sigma_{I_\ell}\).  These attachments again give an invertible operation on each port.  Put \(c=\sigma_1+\sigma_2+\sigma_3\).  The same sign
		calculation, now with zero Walsh endpoints, gives
		\[
		\begin{aligned}
			h(\zero^4)=h(\one^4)&=2c,\\
			h(\one_{I_\ell})=h(\one_{I_\ell^c})
			&=2(2\sigma_\ell-c)\qquad(\ell=1,2,3),
		\end{aligned}
		\]
		and all odd-weight entries vanish.  If \(c\ne0\), this is a
		complement-symmetric mountain; after division by \(2c\), its endpoints
		are one and its three central coefficient-pair values are
		\[
		r_\ell=2\sigma_\ell/c-1,
		\qquad 1+r_1+r_2+r_3=1+2c/c-3=0.
		\]
		Up to the recorded nonzero scalar, this mountain is the output of invertible binary
		attachments to \(q\), so it may be
		adjoined by \cref{eq:finite-equivalent-enlargement}.  It is
		endpoint-nondegenerate, and
		\cref{cor:eight-vertex-complete-closure} makes the enlarged problem either
		tractable or \(\SharpP\)-hard.  Downward inheritance in the first case and
		the gadget reduction in the second transfer that conclusion to
		\(\Gamma\), giving outcome~1.
		
		Assume \(c=0\).  Exactly one nonzero pair is impossible.  Two nonzero Walsh pairs have
		values \(\sigma,-\sigma\); after permuting ports, \(\widetilde h\) factors as
		\[
		\sigma
		\begin{pmatrix}0&1\\-1&0\end{pmatrix}_{14}\otimes
		\begin{pmatrix}0&1\\-1&0\end{pmatrix}_{23},
		\]
		so the factor-reduction equivalence applies.  Finally suppose that all three \(\sigma_i\)
		are nonzero.  Choose distinct \(i,j,k\) with
		\(\{i,j,k\}=\{1,2,3\}\) and \(\sigma_j^2+\sigma_k^2\ne0\), which exists since otherwise subtraction
		would give \(2\sigma_i^2=0\).  Put pair \(i\) at \(12\mid34\), and join
		ports \(3,4\) of one copy to ports \(1,2\) of another by native edges.  The
		actual output is
		\begin{equation*}
			\frac14\widetilde G_i=
			\begin{pmatrix}
				\sigma_i^2&0&0&0\\
				0&-(\sigma_j^2+\sigma_k^2)&-2\sigma_j\sigma_k&0\\
				0&-2\sigma_j\sigma_k&-(\sigma_j^2+\sigma_k^2)&0\\
				0&0&0&\sigma_i^2
			\end{pmatrix}.
		\end{equation*}
		If every phase contraction of \(G_i\) were zero or had transfer class in
		\(C_N\) in Walsh coordinates, its equal nonzero endpoints and
		\cref{lem:odd-cyc-q4-family} would give \(\lambda=1\); complement symmetry
		and odd \(N\) then force every nonzero central coefficient to be
		\(-\sigma_i^2\).  In particular,
		\[
		\sigma_j^2+\sigma_k^2=\sigma_i^2=(\sigma_j+\sigma_k)^2,
		\]
		because \(\sigma_i=-(\sigma_j+\sigma_k)\), contradicting
		\(\sigma_j\sigma_k\ne0\).  Hence an actual phase contraction of \(G_i\) is a
		nonzero binary which is singular or has Walsh transfer class outside
		\(C_N\).  A singular binary gives unary factors and the odd-arity theorem;
		a nonsingular binary of infinite projective transfer order is handled by
		\cref{lem:binary-interpolation}.  Otherwise its transfer has finite
		projective order, contradicting the definition of \(G(\Gamma)\).
		Every hardness or tractability conclusion above gives outcome~1 through the
		cited theorem and the equivalent-enlargement reduction.  In a
		factor case, the attached binaries act invertibly, so the displayed
		factorization pulls back to \(q\); the transported mountain excludes the
		common-one-sided exception, and the factor-reduction equivalence gives the
		reduction in outcome~2.
	\end{proof}
	
	\begin{proof}[Completion of \cref{thm:odd-cyclic}]
		Apply \cref{lem:height-one-card-saturation}.  Its first two outcomes are the
		first two conclusions of the theorem.  In outcome~3, the listed two-port
		contractions and all their positive-arity tensor-prime factors are retained
		in the equivalent set supplied by outcome~3.
		Other fixed gadget outputs and any required factors are adjoined separately
		by \cref{eq:finite-equivalent-enlargement,thm:external-factor}, as in the
		quaternary closure proof.  If a nonbinary
		prime remains, choose one of minimum arity.  The argument preceding
		\cref{lem:odd-cyc-one-matching} proves
		condition~\ref{cond:height-one-contractions} for this signature.
		\Cref{thm:odd-cyc-higher-collapse} bounds its arity by four, and
		\cref{thm:odd-cyc-quaternary-closure} then gives hardness, tractability, or
		a tensor factorization.  The last possibility contradicts the choice of a
		tensor-prime signature, so outcome~1 holds.
		
		In the remaining case every tensor-prime signature is binary.
		By \cref{lem:height-one-card-saturation}\textnormal{(3)}, each has transfer
		class in \(G(\Gamma)\) and Walsh tensor proportional to \(D_rZ\).
		Hence \cref{eq:common-Walsh-product-presentation} holds,
		which proves the third conclusion of \cref{thm:odd-cyclic}.
	\end{proof}
	
	\begin{theorem}[Hardness, tractability, or a strict update]
		\label{thm:height-one-transfer-interface}
		Let \(\Gamma_0\) be the initial signature set, and let
		\((A,\Gamma)\) be a stage of the sequence in
		\cref{lem:normalization-ledger}; in particular,
		\[
		A\Gamma_0\subseteq\Gamma,
		\qquad A^{\mathsf T}XA=\mu_AX\quad(\mu_A\ne0),
		\qquad \KHolant(\Gamma)\equivT\KHolant(\Gamma_0).
		\]
		Suppose \(\Gamma\) is finite, preprocessed, algebraic, and factor-closed,
		and contains the height-one mountain
		\[
		B=\begin{pmatrix}c_0&u\\v&c_1\end{pmatrix},
		\qquad c_0c_1\ne0.
		\]
		Then one of the following occurs:
		\begin{enumerate}
			\item \(\KHolant(\Gamma)\) is \(\SharpP\)-hard, or
			\(\Tract(K\Gamma)\) holds;
			\item there are \(A^+\in GO(X)\) and a finite preprocessed algebraic
			factor-closed set \(\Gamma^+\), continuing the same
			coordinate-transport sequence, such that
			\[
			A^+\Gamma_0\subseteq\Gamma^+,
			\qquad
			(A^+)^{\mathsf T}XA^+=\mu_{A^+}X\quad(\mu_{A^+}\ne0),
			\qquad
			\KHolant(\Gamma^+)\equivT\KHolant(\Gamma),
			\]
			and, for \(D=A^+A^{-1}\), the transported binary
			\(DBD^{\mathsf T}\) remains in \(\Gamma^+\) up to its known nonzero
			scalar.  In this alternative every nonzero directly realizable ordered
			binary over both \(\Gamma\) and \(\Gamma^+\) is nonsingular with
			finite projective transfer order, so
			\(G(\Gamma)\) and \(G(\Gamma^+)\) are defined.  Moreover, either
			\[
			[A]^{-1}G(\Gamma)[A]
			\subsetneq
			[A^+]^{-1}G(\Gamma^+)[A^+],
			\]
			or both values of \(\nu\) are defined and
			\[
			\nu(\Gamma^+)<\nu(\Gamma).
			\]
			Both quantities satisfy the non-strict comparisons and fixed-field
			conditions of \cref{def:height-one-outer-successor}.
		\end{enumerate}
	\end{theorem}
	
	\begin{proof}
		Maintain factor closure using the factor-reduction equivalence.  Each
		nontrivial factorization produces positive-arity factors of strictly smaller
		arity, so this process terminates.  An odd-arity member is handled by the
		odd-signature theorem; a rank-one binary supplies unaries and hence the same
		conclusion.  If no nonbinary tensor-prime signature remains, the
		arity-at-most-two alternative gives tractability.
		
		All equivalent augmentations and global transformations satisfy
		\cref{lem:normalization-ledger}.  Whenever the groups are defined, their
		pullbacks are nested.  The old tensor-prime signatures, up to one common
		invertible transformation, remain in the set, so \(\nu\) cannot increase.  A strict
		group inclusion or decrease of \(\nu\) gives outcome~2.  An augmentation
		preserving both quantities is retained with the same matrix and group;
		\cref{lem:height-one-card-saturation} proves separately that the required
		recursive contraction and factor extraction is finite.
		
		If \(B\) is singular, or if \(BX\) is nonsingular of infinite projective
		order, \cref{lem:binary-interpolation} supplies the rank-one binary used
		in the preceding paragraph.  The same argument applies to every other directly realizable binary that
		is singular or whose nonsingular transfer has infinite projective order.  We may therefore assume
		that every
		nonzero directly realizable binary is nonsingular and of finite projective transfer order.
		Then \cref{lem:bounded-torsion,thm:finite-group-split} give either
		a nonsingular weighted equality in the current set or, after one legal
		\(GO(X)\)-normalization, an odd cyclic or odd dihedral
		residual form.  In the equality case, write
		\(E=\diag(p,q)\in\Gamma\), where \(pq\ne0\), and choose
		a diagonal \(D\in GO(X)\) with \(DED^{\mathsf T}\doteq I\).  The resulting
		legal global normalization satisfies the hypotheses of
		\cref{thm:eq2-kholant-dichotomy}.  Its hard conclusion composes back to
		\(\Gamma\), while its tractable conclusion pulls back through the problem
		equivalence and \cref{lem:normalization-ledger}; hence this is outcome~1.
		
		Apply \cref{thm:odd-cyclic,thm:odd-dihedral} in the two residual forms.
		Their first conclusions give outcome~1, and their updates satisfy
		outcome~2 by \cref{def:height-one-outer-successor}.  Their remaining
		conclusions give \cref{eq:common-Walsh-product-presentation}, which
		\cref{lem:normalization-ledger} transports back to \(\Tract(K\Gamma)\).
		This again gives outcome~1 and completes the case analysis.
	\end{proof}
	
	\begin{proof}[Proof of \cref{cor:height-one-complete}]
		Maintain \((A,\Gamma)\) and the invariant of
		\cref{lem:normalization-ledger} over the fixed initial number field, and
		apply \cref{thm:height-one-transfer-interface}.  Its first outcome is the
		required hardness or tractability conclusion.  The second outcome gives
		an update satisfying \cref{def:height-one-outer-successor}; by
		\cref{lem:height-one-progress-well-founded}, only finitely many such
		updates are possible.  Between these updates, factor closure terminates by
		strict arity descent, and the recursive contractions and factor extraction
		terminate by \cref{lem:height-one-card-saturation}.  Repeated application
		therefore reaches the first outcome.  Finally,
		\cref{lem:normalization-ledger} transports the conclusion back to the
		initial signature set.
	\end{proof}
	
	\section{Mountain Height Two}
	\label{sec:height-two}
	
	A height-two mountain is a quaternary signature with
	\(M(0000)M(1111)\ne0\) that vanishes on odd Hamming weight.  It is an eight-vertex
	signature with those nonzero endpoints.  If it is tensor-prime,
	\cref{thm:quaternary-kholant-dichotomy} applies.  Otherwise the
	parity and endpoint conditions force a \(2+2\) product of nonsingular
	weighted equalities; the factor-reduction equivalence places such an equality
	in the current signature set.  A legal diagonal \(GO(X)\)-normalization then
	sends it to \(I\).  Then \cref{thm:eq2-kholant-dichotomy} closes that
	case.  Both arguments are included in
	\cref{cor:eight-vertex-complete-closure}.
	
	\begin{theorem}[Dichotomy with a height-two mountain]
		\label{cor:height-two-complete}
		Let \(\Gamma\) be a finite algebraic signature set containing a height-two
		mountain.  Then
		\[
		\Tract(K\Gamma)
		\qquad\text{or}\qquad
		\KHolant(\Gamma)\text{ is }\SharpP\text{-hard}.
		\]
	\end{theorem}
	
	\begin{proof}
		The mountain is an eight-vertex signature with two nonzero endpoints, so apply
		\cref{cor:eight-vertex-complete-closure}.  Hardness already
		reduces to \(\KHolant(\Gamma)\), while tractability is the
		predicate required here.
	\end{proof}
	
	\section{Mountain Height Three}
	\label{sec:height-at-least-three}
	
	\begin{theorem}[Height-three mountain dichotomy]
		\label{thm:height-three-complete}
		Let \(\Gamma\) be a finite algebraic set of even-arity signatures such that
		\(\chi_r(x)\in3\mathbb Z\) for every arity-\(r\) member \(f\in\Gamma\)
		and every \(x\in\supp(f)\).  Suppose that \(\Gamma\) contains a mountain
		\[
		M=a\delta_{\zero^6}+b\delta_{\one^6}+Q,
		\qquad ab\ne0,\qquad Q\text{ EO}.
		\]
		Then \(\Tract(K\Gamma)\) holds or \(\KHolant(\Gamma)\) is
		\(\SharpP\)-hard.
	\end{theorem}
	
	The aim is to adjoin a quaternary EO signature outside the binary
	matching-product class.  Then \cref{thm:hc-relative-eo-terminal} gives the
	conclusion for the entire signature set.  A two-copy construction handles a
	matching-product central layer; the other central layers are covered by
	the six-port loop and two-copy calculations below.
	
	\subsection{Matching Products and Two-Copy Constructions}
	
	Let \(V\) be an ordered set of even cardinality.  For distinct \(i,j\in V\),
	write \(\partial_{ij}f\) for the contraction with a native \(X\)-edge
	\[
	(\partial_{ij}f)(x_{V\setminus\{i,j\}})
	=f(0,1,x_{V\setminus\{i,j\}})
	+f(1,0,x_{V\setminus\{i,j\}}).
	\]
	We also write \(\partial_e\) when \(e=\{i,j\}\).
	Define the matching-product class
	\[
	\mathcal B_\times(V)=
	\left\{\lambda\bigotimes_{\{u,v\}\in P}
	\begin{pmatrix}0&p_{uv}\\q_{uv}&0\end{pmatrix}_{u,v}:
	\begin{array}{l}
		P\text{ is a perfect matching of }V,\\
		\lambda\prod_{\{u,v\}\in P}p_{uv}q_{uv}\ne0
	\end{array}\right\},
	\]
	and put \(\mathcal B_\times^0=\mathcal B_\times\cup\{0\}\).
	Each factor has a specified orientation; reversing its ports exchanges its
	two nonzero entries.  In particular,
	\(\mathcal B_{\rm ARS}\subseteq\mathcal B_\times\).
	An \emph{\(X\)-matching product} is a nonzero scalar multiple of
	\(\bigotimes_{e\in P}X_e\).  Such binary factors are tensor-prime, so
	uniqueness of tensor-prime factorization identifies both their port pairs
	and their projective values.  We abbreviate this uniqueness as UPF.
	
	
	\begin{lemma}[A matching-product central layer yields a hard quaternary]
		\label{lem:hc-transversal-cube-terminal}
		Let \(d\ge3\), and let \(\Gamma\) be a finite algebraic signature set
		containing
		\[
		M=a\delta_{\zero^{2d}}+b\delta_{\one^{2d}}+Q,
		\qquad ab\ne0,
		\]
		where, for a perfect matching
		\(P=\{\{u_1,v_1\},\ldots,\{u_d,v_d\}\}\) of the ordered ports,
		\[
		Q=\lambda\bigotimes_{k=1}^{d}
		\begin{pmatrix}0&p_k\\q_k&0\end{pmatrix}_{u_k,v_k},
		\qquad \lambda\prod_{k=1}^{d}p_kq_k\ne0.
		\]
		There is a fixed two-copy gadget over \(\{M\}\), using only native
		\(X\)-edges, whose ordered boundary signature \(h\) is EO and has the six
		weight-two entries
		\[
		\bigl(h_{0011},h_{0101},h_{0110},h_{1001},h_{1010},h_{1100}\bigr)
		=(ab,c_1,c_2,c_3,c_4,ab),
		\qquad c_1c_2c_3c_4\ne0.
		\tag{WT.18a}
		\]
		In particular,
		\[
		\KHolant(\Gamma,h)\leT\KHolant(\Gamma).
		\]
		The signature \(h\) is not a tensor product of two nonsingular
		anti-diagonal binaries.  Moreover, its supported triples include
		\[
		1100\oplus1010\oplus0110=0000,
		\qquad
		0011\oplus0101\oplus1001=1111.
		\]
		The first identity rules out the positive sign and the second rules out
		the negative sign in the support clause of \(\EOTract\).  Hence
		\(\neg\EOTract(\{h\})\), and \cref{thm:external-eo} makes
		\(\KHolant(\{h\})\) \(\SharpP\)-hard.  Since
		\[
		\KHolant(\{h\})
		\leT\KHolant(\Gamma,h)
		\leT\KHolant(\Gamma),
		\]
		the problem \(\KHolant(\Gamma)\) is \(\SharpP\)-hard.
	\end{lemma}
	
	\begin{proof}
		Write the matching of \(Q\) as
		\(P=\{\{w_0,w_1\},\ldots,\{w_{2d-2},w_{2d-1}\}\}\).  On two copies of
		the mountain, define a port bijection \(p\) by fixing
		\(w_0,w_{2d-1}\) and transposing
		\[
		(w_1,w_2),(w_3,w_4),\ldots,(w_{2d-3},w_{2d-2}).
		\]
		Leave \((w_0,w_1)\) external on the first copy and \((w_0,w_2)\)
		external on the second, in the external order
		\((w_0^{(1)},w_1^{(1)},w_0^{(2)},w_2^{(2)})\).  For every
		\(r\notin\{0,1\}\), join the
		first-copy port \(w_r\) to the second-copy port \(p(w_r)\) by one native
		\(X\)-edge.  With both copies central, the matching edges and wires form
		two alternating components: the edge between the first external pair and
		a path of length $4d-3$ between the second.  Both lengths are odd, so each
		weight-one/weight-one boundary pattern has a unique extension and a
		nonzero coefficient $c_i$.
		
		If both copies use endpoints, the native wires force the two constant words
		to be complementary.  This gives exactly the other two EO patterns, with
		unique extensions of value \(ab\).  A mixed endpoint--central extension is
		impossible: for $d\ge3$, some central matching edge has both ports wired to
		the constant copy and is forced equal.  Thus each nonzero coefficient in \emph{(WT.18a)} has exactly one nonzero extension.
		
		The six-word support excludes membership in \(\mathcal B_\times\).  Moreover
		\[
		1100\oplus1010\oplus0110=0000,
		\qquad
		0011\oplus0101\oplus1001=1111.
		\]
		These two supported triples rule out the two common signs in the support
		clause of \(\EOTract\).  The weighted EO dichotomy therefore makes the
		output hard.  By the equivalence convention, the output may be adjoined to
		the signature set containing the mountain.
	\end{proof}
	
	\subsection{Six-Port Loop and Two-Copy Contractions}
	\label{subsec:hc-q6-base}
	
	The finite calculation for both families of contractions is recorded in
	the appendix; the reduction uses only its stated conclusion.
	
	\begin{definition}[Loop and two-copy contractions of a six-port mountain]
		\label{def:hc-q6-joint-deck}
		Let \(U=[6]\), and let
		\[
		M=a\delta_{\zero^6}+b\delta_{\one^6}+Q,
		\qquad
		Q=\sum_{S\in\binom U3}q_S\delta_{\one_S},
		\qquad ab\ne0.
		\]
		For each unordered pair \(\{i,j\}\in\binom U2\), let
		\(L_{ij}=\partial_{ij}M\) be the four-port EO signature obtained by one
		native \(X\)-self-loop, with the residual ports in increasing order.  For
		\(R\in\binom{U\setminus\{i,j\}}2\),
		\begin{equation*}
			L_{ij}(\one_R)=q_{R\cup\{i\}}+q_{R\cup\{j\}}.
			\tag{Q6.3}\label{eq:hc-q6-loop-deck}
		\end{equation*}
		These are the fifteen one-copy contractions.
		
		For a two-copy contraction, choose unordered two-element external sets
		\(A,B\in\binom U2\), give them their canonical increasing orders
		\[
		A=(a_0,a_1),
		\qquad
		B=(b_0,b_1),
		\]
		and retain the first/second labels of the two copies of \(M\).
		Put \(I_A=U\setminus A\) and \(I_B=U\setminus B\), and choose a
		bijection \(\pi:I_A\to I_B\).  Leave \(A\) external on the first copy
		and \(B\) external on the second copy, in the order
		\((a_0,a_1,b_0,b_1)\), and join every \(r\in I_A\) to
		\(\pi(r)\in I_B\) by a native \(X\)-edge.  For
		\(x,y\in\{0,1\}^2\), put
		\[
		A_x=\{a_r:x_r=1\},
		\qquad
		B_y=\{b_r:y_r=1\}.
		\]
		The resulting four-port EO signature \(C_{A,B,\pi}\) has, when
		\(\wt(x)+\wt(y)=2\),
		\begin{equation*}
			C_{A,B,\pi}(x,y)
			=ab\,\one[(x,y)\in\{(00,11),(11,00)\}]
			+\sum_{\substack{T\subseteq I_A\\|T|=3-\wt(x)}}
			q_{A_x\cup T}\,
			q_{B_y\cup(I_B\setminus\pi(T))}.
			\tag{Q6.4}\label{eq:hc-q6-two-copy-deck}
		\end{equation*}
		There are \(15^2\cdot4!=5400\) canonically labelled two-copy contractions.
		The right-hand side of \eqref{eq:hc-q6-two-copy-deck} is the coefficient of
		the complete direct contraction; no individual summand is treated as a
		separately realizable signature.
	\end{definition}
	
	\begin{lemma}[Support forced by six-port contractions]
		\label{lem:hc-q6-joint-rigidity}
		Let \(M\), \(L_{ij}\), and \(C_{A,B,\pi}\) be as in
		\cref{def:hc-q6-joint-deck}.  Suppose that every one of the fifteen
		signatures \(L_{ij}\) and every one of the \(5400\) canonically labelled
		signatures \(C_{A,B,\pi}\) belongs to
		\(\mathcal B_\times^0([4])\) on its four ordered external ports.  Then
		\[
		\operatorname{supp}(Q)\subseteq\{S,U\setminus S\}
		\qquad\text{for some }S\in\binom U3.
		\tag{Q6.5}
		\]
		This holds for arbitrary complex values of the \(q_S\) and every
		\(ab\ne0\).
	\end{lemma}
	
	\begin{proof}
		With \(t=ab\), the signatures \(L_{ij}\) and \(C_{A,B,\pi}\) in
		\cref{eq:hc-q6-loop-deck,eq:hc-q6-two-copy-deck} are precisely the loop
		and labelled two-copy arrays covered by
		\cref{thm:cert-high-charge-q6-interface}.  Its conclusion is
		\emph{(Q6.5)}.
	\end{proof}
	
	\begin{theorem}[A non-matching quaternary output from a six-port mountain]
		\label{thm:hc-q6-complete}
		Let \(\Gamma\) be a finite algebraic signature set containing
		\[
		M=a\delta_{\zero^6}+b\delta_{\one^6}+Q,
		\qquad ab\ne0,
		\qquad Q\text{ EO}.
		\]
		There is a nonzero quaternary EO signature
		\(h\notin\mathcal B_\times\) such that
		\[
		\KHolant(\Gamma,h)\leT\KHolant(\Gamma).
		\]
		The signature \(h\) can be chosen as one of the fifteen one-copy contractions
		\(L_{ij}\), one of the \(5400\) labelled two-copy contractions
		\(C_{A,B,\pi}\) from \cref{def:hc-q6-joint-deck}, or the explicit
		two-copy output of \cref{lem:hc-transversal-cube-terminal}.
	\end{theorem}
	
	\begin{proof}
		We select the signature from one contraction \(L_{ij}\) or one contraction
		\(C_{A,B,\pi}\) when \(Q\notin\mathcal B_\times\), and from the
		explicit two-copy output of
		\cref{lem:hc-transversal-cube-terminal} otherwise.
		If $Q=0$, a contraction \(C_{A,B,\pi}\) has exactly the two endpoint--endpoint
		coefficients $ab$, hence support two rather than four; adjoin this contraction to
		\(\Gamma\).
		
		Suppose $Q\ne0$.  If a loop or
		two-copy contraction lies outside $\mathcal B_\times^0$, adjoin it to \(\Gamma\).
		Otherwise \cref{lem:hc-q6-joint-rigidity} gives
		$\operatorname{supp}(Q)\subseteq\{S,U\setminus S\}$.  Choose $i\in S$
		and $j\notin S$.  The two possible nonzero coefficients of
		$\partial_{ij}M$ have the unique extensions $S,U\setminus S$; hence this
		contraction is nonzero with support at most two, contradicting its
		membership in $\mathcal B_\times^0$.
		
		In particular, if
		\[
		Q=\lambda\bigotimes_{r=0}^2
		\begin{pmatrix}0&p_r\\q_r&0\end{pmatrix}_{w_{2r},w_{2r+1}},
		\qquad \lambda p_0q_0p_1q_1p_2q_2\ne0,
		\]
		the case $d=3$ of \cref{lem:hc-transversal-cube-terminal} gives the
		six-word output with one nonzero extension per supported word \emph{(WT.18a)}, which belongs to
		\(\Gamma\) by \cref{eq:finite-equivalent-enlargement}.
	\end{proof}
	
	
	\begin{proof}[Proof of \cref{thm:height-three-complete}]
		Apply zero/nullary preprocessing and assume \(\neg\Tract(K\Gamma)\).
		By \cref{thm:hc-q6-complete}, a finite equivalent enlargement contains a
		nonzero quaternary EO signature \(h\notin\mathcal B_\times\).
		It has charge zero, so the enlargement preserves the charge hypothesis;
		it retains the two-sided mountain and remains nontractable by
		\cref{lem:downward-inheritance}.
		Since \(\mathcal B_{\rm ARS}\subseteq\mathcal B_\times\),
		\cref{thm:hc-relative-eo-terminal} proves hardness, which
		transfers to the original \(\Gamma\).
	\end{proof}
	
	\Needspace{16\baselineskip}
	\section{Mountain Height at Least Four}
	\label{sec:height-at-least-four}
	
	\begin{theorem}[Mountain dichotomy at every height at least four]
		\label{thm:height-at-least-four}
		Let \(d\ge4\), and let \(\Gamma\) be a finite algebraic set of even-arity
		signatures such that
		\[
		\chi_r(x)\in d\mathbb Z
		\qquad(f\in\Gamma,\ \arity(f)=r,\ x\in\supp(f)).
		\]
		Suppose that \(\Gamma\) contains a mountain
		\[
		M=a\delta_{\zero^{2d}}+b\delta_{\one^{2d}}+Q,
		\qquad ab\ne0,\qquad Q\text{ EO}.
		\]
		Then \(\Tract(K\Gamma)\) holds or \(\KHolant(\Gamma)\) is
		\(\SharpP\)-hard.
	\end{theorem}
	
	The common step attaches extracted binary signatures to selected ports so
	that every self-loop contraction becomes an \(X\)-matching product.
	A single scalar then makes the central tensor satisfy ARS.  For \(d\ge5\),
	the matching-product reconstruction of Shao and Cai makes
	\cref{lem:hc-transversal-cube-terminal} applicable.  We prove this general
	case first.  At height four, a classification of the matching family gives
	the Reed--Muller central layer.  We then normalize the endpoint coefficients
	and obtain equality by a Turing reduction.
	
	\begin{remark}[Relation to the real-valued Holant argument]
		\label{rem:sc20-high-height}
		The port normalization below follows the contraction and binary-attachment
		strategy of \cite[full version, Lemma~8.1]{ShaoCai2020}.  Their transformed
		ambient language satisfies ARS globally.  Here the ambient language is
		complex-valued: charge conservation makes the self-loop outputs EO, and
		\cref{thm:hc-relative-eo-terminal} forces their binary ARS factorization
		unless hardness has already been proved.  ARS alone implies neither EO support nor this
		factorization.
		
		We therefore prove the relative normalization and phase recovery in
		\cref{lem:hc-port-normalization,lem:hc-global-phase}, before applying
		the algebraic reconstruction restated in
		\cref{lem:sc-x-product-reconstruction}.  The central layer is not adjoined
		as a signature.  Hardness comes from our actual mountain gadgets and, at
		height four, the relative Bell reduction; it does not require extending
		the global-ARS hardness arguments to the whole language.
	\end{remark}
	
	\subsection{Common Normalization and Phase Recovery}
	\label{subsec:hc-normalization}
	
	For the proof of \cref{thm:height-at-least-four}, apply zero/nullary
	preprocessing and assume \(\neg\Tract(K\Gamma)\).
	If \(Q=0\), \cref{lem:pure-ge} already gives hardness.  Hence it suffices to
	consider \(Q\ne0\).
	Every native self-loop on the complete mountain kills its endpoints:
	\[
	\partial_{ij}M=\partial_{ij}Q.
	\tag{HN.1}\label{eq:hc-central-card}
	\]
	More generally, charge conservation makes every native gadget output of
	positive arity less than \(2d\) EO.  By
	\cref{thm:hc-relative-eo-terminal}, any such nonzero output outside
	\(\mathcal B_{\rm ARS}\) gives hardness.  Otherwise its factors are
	nonsingular binaries
	\[
	B(z)=\begin{pmatrix}0&z\\\bar z&0\end{pmatrix},
	\qquad z\in\AlgNums^\times,
	\]
	up to nonzero scalars.
	These factors may be adjoined using \cref{thm:external-factor}: the
	retained mountain is two-sided, excluding the common-one-sided exception.
	We absorb each finite equivalent enlargement into \(\Gamma\) by
	\cref{eq:finite-equivalent-enlargement}.  The extracted binaries have
	charge zero, and \cref{lem:downward-inheritance} preserves nontractability.
	
	Attaching the first port of an ordered \(B(z)\) to a mountain port through
	\(X\) acts on that port by
	\[
	T_z=B(z)^{\mathsf T}X=\diag(\bar z,z),
	\qquad T_zB(z)=|z|^2X.
	\tag{HN.2}\label{eq:hc-port-mating}
	\]
	This invertible diagonal map preserves support and both nonzero endpoints.
	For consistently oriented factors,
	\[
	B(z)^{\mathsf T}XB(w)\doteq X
	\quad\Longleftrightarrow\quad
	w/z\in\mathbb R^\times
	\quad\Longleftrightarrow\quad B(w)\doteq B(z).
	\tag{HN.3}\label{eq:hc-binary-splice}
	\]
	Indeed the two off-diagonal entries are \(\bar z w\) and \(z\bar w\).
	In this case \(T_z^{\mathsf T}XT_w\) is also a nonzero scalar multiple of
	\(X\).  All these operations use the extracted ordered binaries on the
	complete mountain.  The conjugate entry in \(B(z)\) is already part of
	that binary; no conjugate mountain is required.
	
	\begin{lemma}[Vanishing pair contractions]
		\label{lem:wt-zero-loop}
		Let \(h\) have arity at least three.  If \(\partial_{ij}h=0\) for every
		pair \(i\ne j\), then \(h(x)=0\) whenever
		\(0<\wt(x)<\arity(h)\).
		In particular, an EO signature with all pair contractions zero is zero.
	\end{lemma}
	
	\begin{proof}
		This is \cite[Lemma~2.7]{CaiFuShao2020ARS}; its three-equation
		argument is purely linear and uses neither ARS nor nonvanishing contractions.
	\end{proof}
	
	\begin{lemma}[Propagation of \(X\)-matching products]
		\label{lem:hc-x-line-propagation}
		Let \(F\) be EO on an even port set \(V\), with \(|V|\ge8\), and suppose
		every nonzero \(\partial_eF\) lies in \(\mathcal B_{\rm ARS}\).
		Either of the following conditions implies that every \(\partial_eF\)
		is an \(X\)-matching product.
		\begin{enumerate}
			\item Every \(\partial_eF\) with \(e\cap\{1,2\}=\varnothing\) is an
			\(X\)-matching product.
			\item Both \(\partial_{12}F\) and \(\partial_{ij}F\) are
			\(X\)-matching products,
			\(\{1,2\}\cap\{i,j\}=\varnothing\), and the matching of
			\(\partial_{ij}F\) pairs \(1,2\) with distinct mates \(u,v\).
		\end{enumerate}
		The remaining pair contractions are allowed to be zero in the hypotheses.
	\end{lemma}
	
	\begin{proof}
		Contracting any pair in an \(X\)-matching product gives another nonzero
		\(X\)-matching product: it either removes a matching edge with scalar
		\(2\), or splices two edges with scalar \(1\).  Thus, whenever a deleted
		pair \(e\) is disjoint from a pair \(r\) for which
		\(\partial_rF\) is an \(X\)-matching product,
		\[
		0\ne\partial_e\partial_rF=\partial_r\partial_eF.
		\tag{HN.4}\label{eq:hc-nonzero-double}
		\]
		In particular \(\partial_eF\ne0\).  A factor whose ports avoid \(r\)
		survives in this double contraction, so UPF makes it projectively \(X\).
		
		For condition~1, only contractions with \(e\cap\{1,2\}\ne\varnothing\)
		remain.  A pair \(r\) disjoint from \(e\), with \(\partial_rF\) an
		\(X\)-matching product, first proves nonvanishing.
		For each binary factor on ports \(a,b\), choose a pair \(r\) avoiding
		\(\{1,2\}\cup e\cup\{a,b\}\).  This excludes at most five ports, so such
		an \(r\) exists and \(\partial_rF\) is an \(X\)-matching product.  The factor survives in
		\eqref{eq:hc-nonzero-double}, proving condition~1.
		
		For condition~2, put \(W=V\setminus\{1,2,i,j\}\), where
		\(u,v\in W\) and \(|W|\ge4\).  For
		\(e\in\binom W2\setminus\{\{u,v\}\}\), nonvanishing follows using the
		\(12\)-contraction.  The ports \(1,2\) cannot be paired in \(\partial_eF\):
		that pair would survive deletion of \(ij\), whereas deleting \(e\) from
		the \(X\)-matching product \(\partial_{ij}F\) pairs \(1,2\) only if \(e=\{u,v\}\).
		Consequently write
		\[
		\partial_eF=C(x_1,x_r)D(x_2,x_s)H.
		\]
		The \(12\)-double contraction makes \(H\) an \(X\)-product and the splice of
		\(C,D\) projectively \(X\).  By \eqref{eq:hc-binary-splice},
		\(C\doteq D\) in these orientations.
		Both factors cannot be affected by deletion of \(ij\), since this would
		splice \(1,2\) together, contradicting the other contraction order.
		A surviving factor is projectively \(X\) by comparison with the
		\(X\)-matching product \(\partial_{ij}F\).  Hence \(C,D\doteq X\), and the whole contraction is an
		\(X\)-matching product.
		
		For \(e=\{u,v\}\), choose two distinct ports \(r,s\in W\setminus\{u,v\}\).
		The contractions at the three disjoint pairs \(12,ij,rs\) are now
		\(X\)-matching products.
		The \(uv\)-contraction is nonzero by the \(12\)-comparison.  Each of its binary
		factors avoids at least one of these three pairs, and hence is
		projectively \(X\) by \eqref{eq:hc-nonzero-double}.
		Thus every \(\partial_eF\) with \(e\subseteq W\) is an
		\(X\)-matching product.
		
		Next let \(e=\{i,k\}\) or \(\{j,k\}\), with \(k\in W\).
		It is nonzero by the \(12\)-comparison.  For a factor on \(a,b\), use
		\(r=12\) if \(\{a,b\}\) avoids \(12\); otherwise choose
		\[
		r\in\binom{W\setminus(\{k\}\cup\{a,b\})}{2}.
		\]
		At most two ports of \(W\) are excluded, so the choice exists.
		Here \(r\) is disjoint from both \(e\) and the factor's ports, and
		\(\partial_rF\) is an \(X\)-matching product.  The factor remains in
		\eqref{eq:hc-nonzero-double}, so it is proportional to \(X\).
		Along with the already established \(X\)-matching product
		\(\partial_{ij}F\), this covers all deleted pairs avoiding
		\(12\).  Condition~1 completes the proof.
	\end{proof}
	
	\begin{lemma}[Relative port normalization at every arity at least eight]
		\label{lem:hc-port-normalization}
		Under the hypotheses of \cref{thm:height-at-least-four}, assume
		\(\neg\Tract(K\Gamma)\) and \(Q\ne0\).
		Then either \(\KHolant(\Gamma)\) is \(\SharpP\)-hard, or a finite equivalent
		enlargement of \(\Gamma\), still supported on half-charges in
		\(d\mathbb Z\), contains an actual mountain
		\[
		M^\circ=A\delta_{\zero^V}+B\delta_{\one^V}+F,
		\qquad |V|=2d,\quad AB\ne0,\quad 0\ne F\text{ EO},
		\]
		such that
		\[
		\partial_eF=\lambda_e\bigotimes_{t\in P_e}X_t,
		\qquad \lambda_e\in\AlgNums^\times
		\qquad\left(e\in\binom V2\right).
		\tag{HN.5}\label{eq:hc-normalized-deck}
		\]
		Here \(P_e\) is a perfect matching of \(V\setminus e\).
		No nonvanishing assumption is imposed on the original contractions of \(Q\).
	\end{lemma}
	
	\begin{proof}
		At every stage retain the complete mountain, adjoin the finitely many
		contractions or factors used, and apply \cref{thm:hc-relative-eo-terminal} again.
		The charge and nontractability invariants described above remain valid.
		Thus every current nonzero contraction is an ARS matching product, or
		hardness has already been proved.
		
		By \cref{lem:wt-zero-loop}, some contraction of \(Q\) is nonzero;
		relabel it as \(\partial_{12}Q\).
		Extract its \(d-1\) binary factors and extend one end of each, on the
		complete mountain, by a matching copy as in \eqref{eq:hc-port-mating}.
		The resulting mountain \(M_0\) has nonzero endpoints and central layer
		\(F_0\), with \(\partial_{12}F_0\) an \(X\)-matching product.
		For every pair \(e\) avoiding \(12\),
		\(\partial_{12}\partial_eF_0=\partial_e\partial_{12}F_0\ne0\).
		Hence all these contractions are nonzero.
		
		Suppose first that \(1,2\) form a matching edge in every such contraction.
		Write \(\partial_eF_0=C_e(x_1,x_2)H_e\).
		The \(12\)-comparison shows that the contraction of \(C_e\) is nonzero
		and \(H_e\) is an \(X\)-product.  If \(e,f\) are disjoint pairs avoiding
		\(12\), their double contraction is nonzero: in the \(e\)-order, deletion of
		\(f\) acts only on the \(X\)-product \(H_e\).
		The factor on \(12\) survives in both orders, so \(C_e\doteq C_f\).
		The graph of disjoint pairs on \(V\setminus\{1,2\}\) is connected.
		Indeed, intersecting distinct pairs exclude only three of its at least
		six ports, leaving a pair disjoint from both.
		Thus all \(C_e\) are projectively one ordered factor \(C\).
		Extract \(C\) from one actual contraction and extend port \(1\) of \(M_0\) by
		a matching copy.  Every contraction whose deleted pair avoids \(12\) is now an
		\(X\)-matching product.  After reapplying
		\cref{thm:hc-relative-eo-terminal}, condition~1 of
		\cref{lem:hc-x-line-propagation} makes every pair contraction an
		\(X\)-matching product.
		
		Otherwise choose a pair \(ij\) avoiding \(12\) whose contraction has distinct
		mates \(u,v\) for \(1,2\), and write
		\[
		\partial_{ij}F_0=C(x_1,x_u)D(x_2,x_v)H.
		\]
		The \(12\)-comparison makes \(H\) an \(X\)-product and gives
		\(C\doteq D\) by \eqref{eq:hc-binary-splice}.
		Extract both factors from this actual contraction and extend ports \(1,2\) of
		\(M_0\) by matching copies.  The new \(ij\)-contraction is an \(X\)-matching product.
		The effective \(12\)-kernel is \(T_C^{\mathsf T}XT_D\doteq X\), so the
		\(12\)-contraction remains an \(X\)-matching product as well.  The
		diagonal maps preserve the distinct mates of ports \(1,2\).  Reapply
		\cref{thm:hc-relative-eo-terminal} to the current contractions, then use condition~2 of \cref{lem:hc-x-line-propagation}.
		
		All extensions are invertible and preserve the central support and the
		nonzero endpoints.  Their number depends only on the fixed source
		signature.  Factor extraction and direct gadget substitution give the
		claimed finite equivalence.  Only contractions of the central layer are
		adjoined; the central layer itself is used only in the algebraic argument.
	\end{proof}
	
	\begin{lemma}[A common phase for EO contractions]
		\label{lem:hc-global-phase}
		Let \(F\) be algebraic and EO on an even port set \(V\), with \(|V|\ge6\), and suppose
		\eqref{eq:hc-normalized-deck} holds with every \(\lambda_e\ne0\).
		There is \(c\in\AlgNums^\times\) such that \(G=c^{-1}F\) is real-valued
		and invariant under bitwise complement, and every coefficient
		\(\lambda_e/c\) in \eqref{eq:hc-normalized-deck} is positive real.
		In particular \(G\) satisfies exact ARS.
	\end{lemma}
	
	\begin{proof}
		For disjoint pairs \(e,f\), the two contraction orders give
		\[
		\partial_f\partial_eF=m_{ef}\lambda_e X_{P'},
		\qquad
		\partial_e\partial_fF=m_{fe}\lambda_f X_{P''},
		\qquad m_{ef},m_{fe}\in\{1,2\}.
		\]
		Here \(X_P=\bigotimes_{t\in P}X_t\) has only \(0,1\) entries.
		Both tensors are nonzero, so commutation identifies their supports and
		gives
		\[
		\lambda_e/\lambda_f=m_{fe}/m_{ef}\in\{1,2,1/2\}.
		\]
		The graph of disjoint pairs on \(V\) is connected: two intersecting pairs
		have a common disjoint pair since \(|V|\ge6\).
		Fix one pair \(e_0\) and take \(c=\lambda_{e_0}\).
		Every contraction of \(G=c^{-1}F\) is therefore real and
		complement-invariant, with a positive coefficient.
		
		Let \(\overline G\) denote coefficientwise conjugation and
		\(G^{\rm rev}(x)=G(\bar x)\).
		Contraction commutes with both operations.  Hence every pair contraction
		of \(G-\overline G\) and of \(G-G^{\rm rev}\) vanishes.
		Both differences are EO, so \cref{lem:wt-zero-loop} makes them zero.
		This is an algebraic conclusion and does not assert that \(G\) is a
		realizable signature.
	\end{proof}
	
	We use precisely the following algebraic part of the real-valued
	classification.
	
	\begin{lemma}[Shao--Cai matching-product reconstruction
		{\cite[full version, Lemma~2.11]{ShaoCai2020}}]
		\label{lem:sc-x-product-reconstruction}
		Let \(G\) be an EO tensor of even arity \(n\ge8\) satisfying exact ARS.
		Suppose
		\[
		\partial_eG=\mu_e\bigotimes_{t\in P_e}X_t,
		\qquad \mu_e\in\mathbb R^\times
		\qquad\left(e\in\binom{[n]}2\right),
		\]
		where each \(P_e\) is a perfect matching of the residual ports.
		If \(n\ge10\), then \(G\) is an \(X\)-matching product.
		The same conclusion holds for \(n=8\) if \(P_A\cap P_B\ne\varnothing\)
		for some \(A\ne B\) with \(A\cap B\ne\varnothing\); in that case the
		\(X\)-factor on that shared edge is also a factor of \(G\).
	\end{lemma}
	
	\subsection{Height at Least Five}
	\label{subsec:height-at-least-five}
	
	\begin{proof}[Proof of \cref{thm:height-at-least-four} for \(d\ge5\)]
		As above, assume \(\neg\Tract(K\Gamma)\) and \(Q\ne0\).
		Apply \cref{lem:hc-port-normalization}.  Outside its hard outcome there
		is an actual mountain
		\(M^\circ=A\delta_{\zero^{2d}}+B\delta_{\one^{2d}}+F\)
		with \(AB\ne0\) and every pair contraction of \(F\) an
		\(X\)-matching product.
		By \cref{lem:hc-global-phase}, \(G=c^{-1}F\) is EO with exact ARS
		and nonzero real \(X\)-product contractions.  Since \(2d\ge10\),
		\cref{lem:sc-x-product-reconstruction} makes \(G\), and hence \(F\),
		an \(X\)-matching product.
		Apply \cref{lem:hc-transversal-cube-terminal} to the complete
		\(M^\circ\).  Its two-copy output is hard, and the finite equivalence
		transfers hardness to the original \(\Gamma\).
	\end{proof}
	
	\subsection{Height Four}
	\label{subsec:height-four}
	
	The eight-port exception is the central-layer Reed--Muller tensor
	\[
	R_8^{\rm cen}(x)=
	\one\left[x=(a+u\cdot t)_{t\in\Ftwo^3}
	\text{ for some }a\in\Ftwo,\ 0\ne u\in\Ftwo^3\right].
	\tag{WT.20}
	\]
	This tensor has fourteen supported words and is not a matching product, but each contraction
	is an \(X\)-product.  For deleted ports \(p,q\), the surviving pairs are
	\(\{t,t+p+q\}\): each choice of one endpoint from every remaining pair
	has a unique extension with
	\(u\cdot(p+q)=1\).  A matching edge can therefore belong to both \(P_A\) and \(P_B\),
	for distinct \(A,B\), only when \(A\cap B=\varnothing\), so the eight-port overlap clause of
	\cref{lem:sc-x-product-reconstruction} does not apply.
	
	For an eight-port set \(V\) and a bijection \(\iota:V\to\Ftwo^3\), write
	\[
	(\iota_*F)((x_t)_{t\in\Ftwo^3})
	=F((x_{\iota(v)})_{v\in V}).
	\]
	All canonical Reed--Muller statements below use this explicit port
	relabelling.
	
	\subsubsection{Translation Form of the Matching Family}
	
	\begin{lemma}[A shared edge forces a matching product]
		\label{lem:q8-scope-collision}
		Let \(F\) be an arity-eight EO tensor satisfying
		\eqref{eq:hc-normalized-deck}.
		If \(P_A\cap P_B\ne\varnothing\) for some \(A\ne B\) with
		\(A\cap B\ne\varnothing\), then \(F\) is an \(X\)-matching product.
		Consequently, any finite algebraic signature set containing an actual
		mountain
		\(M^\circ=A\delta_{\zero^8}+B\delta_{\one^8}+F\), with \(AB\ne0\),
		has a \(\SharpP\)-hard \(K\)-Holant problem.
	\end{lemma}
	
	\begin{proof}
		Apply \cref{lem:hc-global-phase} and then the eight-port clause of
		\cref{lem:sc-x-product-reconstruction}.
		This makes \(F\) an \(X\)-matching product.
		The complete mountain now satisfies
		\cref{lem:hc-transversal-cube-terminal}.
	\end{proof}
	
	
	\begin{lemma}[Translation form of the matching family]
		\label{lem:q8-collision-free-atlas}
		Let \(F\) be an arity-eight EO tensor on a port set \(V\), and suppose
		that for every \(u\ne v\) one has
		\[
		\partial_{uv}F
		=\lambda_{uv}\bigotimes_{e\in P_{uv}}X_e,
		\qquad \lambda_{uv}\ne0,
		\]
		where \(P_{uv}\) is a perfect matching of \(V\setminus\{u,v\}\).
		Suppose \(P_A\cap P_B=\varnothing\) whenever \(A\ne B\) and
		\(A\cap B\ne\varnothing\).  Then there is a bijection
		\(\iota:V\to\Ftwo^3\) such that, after relabelling each port \(v\) by
		\(\iota(v)\), for $p\ne q$ and $h=p+q$ one has
		\[
		P_{pq}
		=\bigl\{\{t,t+h\}:t\in \Ftwo^3\setminus\{p,q\}\bigr\},
		\tag{Q8.14}
		\]
		where duplicate unordered pairs are suppressed.
	\end{lemma}
	
	\begin{proof}
		We apply uniqueness of tensor-prime factorization and a finite matching
		classification.  For disjoint deleted pairs
		$e,r$, let $\kappa_r(P_e)$ be the two-edge matching left after contracting
		$r$ in the $X$-product on $P_e$: delete $r$ when $r\in P_e$, and
		otherwise splice the mates of the two ends of $r$ and retain the untouched
		edge.  Both operations are nonzero for $X$ factors.  Commutativity of the
		two loops and unique factorization give the exact constraint
		\[
		\kappa_r(P_e)=\kappa_e(P_r)
		\qquad(e\cap r=\varnothing).
		\tag{Q8.14a}
		\]
		
		The condition on \(P_A\cap P_B\) and \emph{(Q8.14a)} are exactly the
		hypotheses of \cref{thm:cert-q8-collision-free-interface}, which supplies
		the required bijection and \emph{(Q8.14)}.  Thus the identity obtained by commuting two direct contractions forces
		the matching family in \emph{(Q8.14)}; an arbitrary partition of the
		complete graph into perfect matchings would not suffice.
	\end{proof}
	
	\begin{lemma}[The matching family determines the central $RM(1,3)$ coefficients]
		\label{lem:q8-canonical-support}
		Let \(F\) be an arity-eight EO tensor on a port set \(V\) such that
		\[
		\partial_{uv}F
		=\lambda_{uv}\bigotimes_{e\in P_{uv}}X_e,
		\qquad \lambda_{uv}\ne0,
		\]
		for every $u\ne v$, where $P_{uv}$ is a perfect matching of
		$V\setminus\{u,v\}$.  Suppose \(P_A\cap P_B=\varnothing\) whenever
		\(A\ne B\) and \(A\cap B\ne\varnothing\).  Choose a bijection
		\(\iota:V\to\Ftwo^3\) supplied by
		\cref{lem:q8-collision-free-atlas}, put \(F^\iota=\iota_*F\), and
		write
		\[
		c_S=F^\iota(\one_S),
		\qquad S\in\binom{\Ftwo^3}4.
		\]
		Then
		\[
		\supp(F^\iota)
		=\left\{\one_S:
		S\in\binom{\Ftwo^3}4,\ 
		\bigoplus_{s\in S}s=0\right\},
		\tag{Q8.17}
		\]
		the set of the fourteen affine planes.  Moreover all fourteen nonzero
		coefficients in \emph{(Q8.17)} are equal.
	\end{lemma}
	
	\begin{proof}
		Relabel the ports through \(\iota\), continue to denote the resulting tensor
		by \(F^\iota=\iota_*F\), and write \(V=\Ftwo^3\).  For
		$S\in\binom V4$, $p\in S$, and $q\notin S$, put
		\[
		S'=S\mathbin\triangle\{p,q\},\qquad
		T=S\setminus\{p\},\qquad
		\sigma(S)=\bigoplus_{s\in S}s,\qquad h=p+q.
		\]
		The coefficient of the actual contraction at $\one_T$ is the complete
		two-extension sum
		\[
		(\partial_{pq}F^\iota)(\one_T)=c_S+c_{S'}.
		\]
		The word \(\one_T\) selects one endpoint from each edge of the
		matching \emph{(Q8.14)} if and only if
		\[
		\sigma(S)\in\{0,h\}.
		\tag{Q8.19}
		\]
		Indeed, in $V/\langle h\rangle$ the fibre $\{p,q\}$ contains one point of
		$S$, while the other occupancies are $1,1,1$ or a permutation of $2,1,0$.
		Their quotient xor is respectively zero or the sum of the two odd fibres,
		proving \emph{(Q8.19)}.
		Coefficients at these words equal $\lambda_{pq}$ and all others vanish.
		Thus the \(c_S\) and nonzero \(\lambda_{pq}\) satisfy the swap rule in
		\cref{thm:cert-q8-canonical-support-interface}.  That theorem eliminates
		all coefficients on four-sets that are not affine planes and identifies
		the fourteen plane coefficients
		with one \(c\ne0\), proving \emph{(Q8.17)}.
	\end{proof}
	
	Equivalently, the central-layer Reed--Muller tensor from \emph{(WT.20)} is
	\[
	R_8^{\rm cen}(x)
	=\one\left[\wt(x)=4,
	\bigoplus_{t:x_t=1}t=0\right].
	\]
	It consists of the fourteen nonconstant words of $RM(1,3)$.
	Consequently the full core defined in \cref{eq:RM-core} is
	\begin{equation*}
		\RMcore(x)
		=\delta_{\zero^8}(x)+\delta_{\one^8}(x)+R_8^{\rm cen}(x)
		=\one[x\in RM(1,3)].
	\end{equation*}
	Thus \(\RMcore\) is the full sixteen-word Reed--Muller signature.
	
	\subsubsection{Normalizing the Endpoint Coefficients}
	
	The preceding lemma has made all central coefficients equal.
	One information-set contraction detects the remaining endpoint invariant.
	
	\begin{lemma}[Normalizing the endpoint coefficients of \(\RMcore\)]
		\label{lem:q8-canonical-endpoint-repair}
		Let
		\[
		M^\circ=A\delta_{\zero^8}+B\delta_{\one^8}+cR_8^{\rm cen},
		\qquad ABc\ne0.
		\]
		Then either the two-copy native-\(X\) construction below yields a nonzero
		arity-six EO signature outside
		\(\mathcal B_\times\), or
		\[
		AB=c^2
		\tag{Q8.25}\label{eq:q8-canonical-endpoint}
		\]
		and, for any algebraic \(s\) with \(s^8=c/A\),
		\[
		D(s)^{\otimes8}M^\circ=c\RMcore,
		\qquad
		D(s)=\diag(s,s^{-1}),
		\qquad
		D(s)^{\mathsf T}XD(s)=X.
		\tag{Q8.26}\label{eq:q8-canonical-normalization}
		\]
		The first outcome uses two source vertices, four information-set edges, and
		one final \(X\)-closure.
	\end{lemma}
	
	\begin{proof}
		Index the ports by \(V=\Ftwo^3\) and take the complementary affine
		information sets
		\[
		I_0=\{000,100,010,001\},
		\qquad
		I_1=V\setminus I_0.
		\]
		Join corresponding \(I_0\)-ports of two copies of \(M^\circ\) by four
		native \(X\)-edges.  An affine word is determined by its values on either
		information set.  Hence the internal equations force the two
		\(RM(1,3)\) words to be complementary on all eight coordinates.  The
		remaining ports form four opposite pairs indexed by \(I_1\), and their
		coefficient function on the four independent pair-choice bits is
		\[
		H(y)=
		\begin{cases}
			AB,&y=0^4\text{ or }1^4,\\
			c^2,&\text{otherwise}.
		\end{cases}
		\]
		Close one of these four pairs by a native \(X\)-edge.  On the three
		remaining matching bits the two antipodal coefficients are \(AB+c^2\),
		whereas the other six coefficients are \(2c^2\).  If \(AB=-c^2\), the
		output is nonzero but its support omits two of the eight choices of one
		endpoint from each remaining pair.  Otherwise that
		support uniquely determines its three opposite pairs, and a rank-one
		coefficient table would have all \(2\times2\) minors vanish after fixing
		any one of the three pair-choice bits.  Choosing a minor through
		an antipodal word instead gives determinant
		\[
		\det\begin{pmatrix}
			AB+c^2&2c^2\\
			2c^2&2c^2
		\end{pmatrix}
		=2c^2(AB-c^2).
		\]
		Thus \(AB\ne c^2\) gives the first outcome, and the remaining case is
		\eqref{eq:q8-canonical-endpoint}.
		
		Finally, \(D(s)^{\otimes8}\) multiplies a weight-\(k\) coefficient by
		\(s^{8-2k}\).  It fixes the central layer, sends \(A\) to \(As^8=c\),
		and sends \(B\) to \(Bs^{-8}=AB/c=c\).  This proves
		\eqref{eq:q8-canonical-normalization}; direct multiplication gives the
		displayed \(GO(X)\) identity.
	\end{proof}
	\subsubsection{Obtaining Equality from the Reed--Muller Signature}
	
	We use the Bell contraction identity and the occurrence reduction of
	\cite[full version, Definition~8.7 and Lemma~8.9]{ShaoCai2020}.
	The proof below supplies the ambient-relative step for our complex-valued
	high-charge language; it does not assume global ARS.
	Recall that \(\RMcore\) is the full sixteen-word \(RM(1,3)\) signature,
	and put \(\mathcal B=\{I,X,Y,Z\}\), with the ordered matrices of
	\cref{eq:bell-matrix-dictionary}.
	
	\begin{lemma}[The Reed--Muller Bell reduction under high charge]
		\label{thm:q8-normalized-bell-lift}
		Let \(\Gamma\) be a finite algebraic set of even-arity signatures such that
		\[
		\chi_r(x)\in4\mathbb Z
		\qquad
		\bigl(f\in\Gamma,\ \arity(f)=r,\ x\in\supp(f)\bigr).
		\]
		Suppose \(\Gamma\) contains \(c\RMcore\) for a known \(c\ne0\).  Then
		\[
		\KHolant(\Gamma,\mathcal B)\leT\KHolant(\Gamma).
		\tag{Q8.31}\label{eq:q8-relative-bell-reduction}
		\]
		Consequently, if \(\neg\Tract(K\Gamma)\), then
		\(\KHolant(\Gamma)\) is \(\SharpP\)-hard.
	\end{lemma}
	
	\begin{proof}
		Throughout, boundary tensors refer to direct native-\(X\) gadgets over
		the indicated signature set, with ordered boundary ports.  The gadgets
		may be disconnected; components without boundary contribute scalars.
		
		\smallskip
		\noindent\emph{Occurrence reduction.}
		For each ordered \(b\in\mathcal B\), attaching \(b\) to two ports of
		\(c\RMcore\) leaves, up to a known nonzero scalar and a fixed ordering of
		the six remaining ports, \(b^{\otimes3}\).
		This is the identity used in
		\cite[full version, Lemma~8.9]{ShaoCai2020}.
		It remains valid in native \(X\)-coordinates, since
		\(XbX=\pm b\).  Thus any three \(b\)-vertices can be replaced by one
		\(b\)-vertex and one core vertex, reducing their number by two.
		Repeating this operation leaves at most two occurrences, with linear
		growth in the instance size and a tracked nonzero scalar.
		This is a tensor identity, so arbitrary complex labels in the remaining
		instance do not affect the reduction.
		
		For \(b=Y\) or \(I\), a self-loop on a \(b\)-vertex has value zero,
		while two such vertices sharing an edge collapse to a native \(X\)-wire,
		up to their recorded orientation signs.  If they share both edges, they
		form a component of known scalar value.  We may therefore assume that
		the remaining \(b\)-vertices are disjoint.  Removing \(k\le2\) of them
		leaves a direct gadget with \(2k\) ordered boundary ports.
		
		\smallskip
		\noindent\emph{Adjoining \(Y\).}
		The following identity holds on all sixteen four-bit assignments:
		\[
		Y_{12}\otimes Y_{34}
		=X_{14}\otimes X_{23}-X_{13}\otimes X_{24}.
		\tag{Q8.32}\label{eq:q8-y-square-closures}
		\]
		Hence two disjoint \(Y\)-occurrences can be evaluated by the difference
		of two native-\(X\) closures of the residual gadget.
		The effective kernel of a \(Y\)-vertex attached through two native edges
		is \(XYX=-Y\), so the two signs cancel here.
		
		It remains to evaluate one occurrence.  Write
		\(\langle A,B\rangle=\sum_x A(x)B(x)\) for the bilinear coefficient
		pairing.  If every binary boundary tensor \(H\) over \(\Gamma\) satisfies
		\(\langle H,Y\rangle=0\), every such instance has value zero.
		Otherwise fix one direct binary gadget \(G\) over \(\Gamma\) with
		\[
		\kappa=\langle G,Y\rangle\ne0.
		\]
		Its size and the nonzero algebraic number \(\kappa\) depend only on the
		fixed language \(\Gamma\), not on the input instance; its truth table is
		computed by finite exact summation.
		For the residual binary \(H\) of an input, apply
		\eqref{eq:q8-y-square-closures} to the disconnected gadget \(H\otimes G\).
		Two oracle closures compute
		\[
		\langle H\otimes G,Y\otimes Y\rangle
		=\kappa\langle H,Y\rangle.
		\]
		Dividing by the known \(\kappa\), and restoring the sign from \(XYX=-Y\),
		gives the required one-occurrence value.
		The case distinction is a statement about the fixed language: in the
		second case one fixed witness gadget is part of the reduction, and in the
		first case the zero rule applies to every input.
		Together with occurrence reduction, this proves
		\[
		\KHolant(\Gamma,Y)\leT\KHolant(\Gamma).
		\tag{Q8.33}
		\]
		
		\smallskip
		\noindent\emph{Adjoining \(I\).}
		Put \(\Lambda=\Gamma\cup\{Y\}\).  The added signature has charge zero,
		so all supported half-charges in \(\Lambda\) still lie in \(4\mathbb Z\).
		Charge conservation across native edges therefore makes every binary
		boundary tensor over \(\Lambda\) anti-diagonal and every quaternary
		boundary tensor EO.  No factorization or ARS assumption on these tensors
		is required.
		
		Use occurrence reduction for \(I\).  One disjoint occurrence evaluates
		to zero, since \(XIX=I\) is orthogonal to every anti-diagonal binary for
		the bilinear pairing.  For two disjoint occurrences, let \(H\) be the
		four-port residual tensor, with the removed vertices on \(12\) and \(34\).
		On its EO support,
		\[
		\left.(I_{12}\otimes I_{34})\right|_{HW_4^=}
		=
		\left.
		\frac{X_{13}\otimes X_{24}
			+X_{14}\otimes X_{23}
			-X_{12}\otimes X_{34}}{2}
		\right|_{HW_4^=}.
		\tag{Q8.34}\label{eq:q8-equality-eo-closures}
		\]
		Thus three native-\(X\) closures, with respective coefficients
		\(1/2,1/2,-1/2\), evaluate the desired contraction.  These are
		\(\KHolant(\Lambda)\)-oracle queries, valid for arbitrary complex
		coefficients of \(H\).  We have proved
		\[
		\KHolant(\Gamma,Y,I)\leT\KHolant(\Gamma,Y).
		\tag{Q8.35}
		\]
		Finally \(YX I=Z\) supplies the remaining Bell signature by a direct
		gadget, and an \(X\)-vertex merely subdivides a native wire.
		This proves \eqref{eq:q8-relative-bell-reduction}.
		
		If \(\neg\Tract(K\Gamma)\), restriction gives
		\(\neg\Tract(K(\Gamma\cup\{Y,I\}))\).
		The binary-equality \(K\)-Holant dichotomy
		\cref{thm:eq2-kholant-dichotomy}, namely
		\cite[Theorem~1.2 and Lemma~2.13]{Liu2026DEQ2Holant},
		makes the enlarged problem hard.  The two displayed reductions transfer
		this hardness to \(\KHolant(\Gamma)\).
	\end{proof}
	
	
	
	\begin{proof}[Proof of \cref{thm:height-at-least-four} for \(d=4\)]
		Assume \(\neg\Tract(K\Gamma)\) and \(Q\ne0\), as in
		\cref{subsec:hc-normalization}.
		Apply \cref{lem:hc-port-normalization} and suppose hardness has not
		already been proved.  If \(P_A\cap P_B\ne\varnothing\) for some
		\(A\ne B\) with \(A\cap B\ne\varnothing\), hardness follows from
		\cref{lem:q8-scope-collision}.  Otherwise
		\cref{lem:q8-collision-free-atlas,lem:q8-canonical-support} give a
		bijection \(\iota:V\to\Ftwo^3\) under which the actual mountain is
		\[
		M^\circ=A\delta_{\zero^8}+B\delta_{\one^8}+cR_8^{\rm cen},
		\qquad ABc\ne0.
		\]
		The first outcome of \cref{lem:q8-canonical-endpoint-repair} is an actual
		nonzero EO signature outside \(\mathcal B_\times\), so
		\cref{thm:hc-relative-eo-terminal} gives hardness.
		Otherwise \eqref{eq:q8-canonical-normalization} supplies
		\(D(s)=\diag(s,s^{-1})\in GO(X)\) with
		\[
		D(s)^{\otimes8}M^\circ=c\RMcore.
		\]
		Apply \(D(s)\) to the entire current signature set and again call it
		\(\Gamma\).  This whole-set equivalence preserves support and charge,
		and \cref{lem:GOX-tract-invariance} preserves nontractability.
		Now \(c\RMcore\in\Gamma\), so \cref{thm:q8-normalized-bell-lift}
		completes the proof.
	\end{proof}
	
	\section{Synthesis and Proof of the Dichotomy}
	
	\subsection{Combining the Height Cases}
	\label{sec:mountain-synthesis}
	
	The height analyses now give a dichotomy for the entire signature set;
	every tractable conclusion concerns one common presentation of that set.
	
	\begin{theorem}[Dichotomy for sets containing a mountain]
		\label{thm:mountain-interface}
		Let \(\cG\) be a finite algebraic set of even-arity signatures containing a
		mountain of height \(d\).  If \(d\ge3\), assume that
		\(\chi_r(x)\in d\mathbb Z\) for every \(f\in\cG\) of arity \(r\) and
		every \(x\in\supp(f)\), as for the gcd-height mountain produced by
		\cref{lem:charge-bridge}.  Then
		\[
		\Tract(K\cG)
		\qquad\text{or}\qquad
		\KHolant(\cG)\text{ is }\SharpP\text{-hard}.
		\]
	\end{theorem}
	
	\begin{proof}
		Apply \cref{cor:height-one-complete} for \(d=1\),
		\cref{cor:height-two-complete} for \(d=2\),
		\cref{thm:height-three-complete} for \(d=3\), and
		\cref{thm:height-at-least-four} for \(d\ge4\).
		These cases exhaust all heights.
	\end{proof}
	
	\subsection{Proof of the Main Theorem}
	\label{sec:assembly}
	
	Recall that \cref{thm:main} gives an \(\FP\)-versus-\(\SharpP\)-hard
	dichotomy characterized by the tractability predicate \(\Tract\) defined
	in \cref{sec:predicate}.  We now complete its proof by combining the
	preceding results to establish the corresponding dichotomy in
	\(K\)-coordinates.  For a signature set \(\cG\) in these coordinates,
	write \(\Tract_K(\cG):=\Tract(K\cG)\).
	
	\begin{theorem}[Structural classification in \(K\)-coordinates]
		\label{thm:K-structural}
		Let \(\cG\) be a finite algebraic signature set.  If
		\(\Tract_K(\cG)\) holds, then \(\KHolant(\cG)\in\FP\); otherwise
		\(\KHolant(\cG)\) is \(\SharpP\)-hard.  The applicable side is decidable
		from exact algebraic input.
	\end{theorem}
	
	\begin{proof}
		The tractable implication and its effective recognition follow from
		\cref{thm:predicate-algorithms}.  For the converse, assume that
		\(\KHolant(\cG)\) is not \(\SharpP\)-hard; we prove
		\(\Tract_K(\cG)\).
		Apply zero and nullary preprocessing and let \(\Lambda\) be the resulting
		positive-arity signature set.  We use the equivalence convention throughout;
		\cref{lem:normalization-ledger} transports both common presentations and
		complexity conclusions back to \(\cG\).
		
		Odd arity is closed first by \cref{thm:external-odd}.  In the all-even case,
		test common one-sided support before factor extraction; this is
		alternative~\ref{case:one-sided} by charge conservation and
		\cref{thm:external-eo}.  If every member is single-layer, the remaining
		two-sided case is alternative~\ref{case:single-weight} by
		\cref{thm:external-single-weight}.  Otherwise choose two signature--support-point pairs with
		\begin{equation*}
			\begin{aligned}
				f_\pm&\in\Lambda,
				\qquad x_\pm\in\supp(f_\pm),\\
				\chi_{\arity(f_+)}(x_+)&>0>
				\chi_{\arity(f_-)}(x_-).
			\end{aligned}
		\end{equation*}
		The two signs need not occur in the same signature.
		
		Apply \cref{lem:factor-saturation} and continue to write \(\Lambda\) for
		the factor-closed set.  Both selected signature--support-point pairs remain present.  An
		odd-arity factor returns to \cref{thm:external-odd}; if no higher-arity
		prime remains,
		alternative~\ref{case:arity-two} holds.
		
		We are left with a two-sided, multilayer all-even signature set containing a
		higher-arity prime.  Let \(d\) be the gcd of its nonzero supported
		half-charges, computed on the current set \(\Lambda\) after factor closure.
		Apply \cref{lem:charge-bridge} to \(\Lambda\); the two selected pairs
		verify its sign hypothesis.  In
		its pin outcome, take \(\DeltaZero,\DeltaOne\in\Lambda\) by the equivalence
		convention and apply the odd-arity theorem.  Otherwise that lemma places in
		\(\Lambda\) a mountain of the charge gcd \(d\), with
		\[
		\chi_{\arity(f)}(x)\in d\mathbb Z
		\qquad(f\in\Lambda, x\in\supp(f)).
		\]
		Thus \cref{thm:mountain-interface} applies.  Its hard outcome is excluded,
		so \(\Tract_K(\Lambda)\), and hence \(\Tract_K(\cG)\), holds.
		
		For completeness, if the proof gives a common \(K\)-presentation
		\(W\), the ordinary-coordinate basis is \(M=KW\).  Since
		\(K^{\mathsf T}K=X\),
		\begin{equation*}
			M^{-1}(K\cG)=W^{-1}\cG,
			\qquad M^{\mathsf T}M=W^{\mathsf T}XW,
		\end{equation*}
		so one basis handles both the edge and the entire signature set.
	\end{proof}
	
	\begin{proof}[Proof of \cref{thm:main}]
		Decidability and polynomial-time evaluation in every tractable case follow
		from \cref{thm:predicate-algorithms}.  For hardness, put
		\(\cG=K^{-1}\cF\).  By \cref{eq:K-equivalence},
		\begin{equation*}
			\Holant(\cF)\equivT\KHolant(\cG).
		\end{equation*}
		By the definition of \(\Tract_K\),
		\begin{equation*}
			\Tract_K(\cG)=\Tract(K\cG)=\Tract(\cF).
		\end{equation*}
		If \(\neg\Tract(\cF)\), then \(\neg\Tract_K(\cG)\), so
		\cref{thm:K-structural} implies that \(\KHolant(\cG)\) is
		\(\SharpP\)-hard.  The displayed problem equivalence transfers this
		hardness to \(\Holant(\cF)\).  All reductions are polynomial-time Turing
		reductions, as required.
	\end{proof}
	
	
	\section{Conclusion}
	\label{sec:conclusion}
	
	We have proved an \(\FP\)-versus-\(\SharpP\)-hard dichotomy for Boolean
	Holant on general graphs with arbitrary finite sets of algebraic
	complex-valued signatures, under polynomial-time Turing reductions.
	The classification uses the same seven conditions as the odd-arity
	dichotomy, and this criterion is decidable from exact algebraic input.
	
	For the all-even cases, charge conservation and tensor-factor extraction
	either lead to a known classification or produce a mountain signature.
	The analysis by height, together with the existing odd-arity,
	Eulerian-orientation, binary-equality, and quaternary dichotomies, yields
	the full result.  In each holographic tractable case, one basis works
	simultaneously for the edge and every signature in the set.
	
	
	
	\appendix
	\section{Exact Finite Certificates}
	\label{sec:exact-finite-certificates}
	
	This appendix gives the exact finite calculations used in the
	odd-dihedral order-three and high-charge six- and eight-port arguments.
	Each detailed proposition is followed by the consequence used in the main
	proof.  We retain the exhaustive counts, equations, verification details,
	and the distinction between the finite calculations and the reductions
	proved in the text.
	
	The accompanying programs reconstruct and check the finite objects from
	the definitions below. Reproduction instructions and the correspondence
	between manuscript statements and programs are given under Code Availability.
	The algebraic reductions and direct-gadget arguments remain in the text.
	
	Throughout this appendix, ranks and linear systems are computed exactly
	over the indicated characteristic-zero field.  A projective vector is
	normalized by making its first nonzero coordinate one, except when a
	primitive integer representative with first nonzero coordinate positive
	is explicitly requested.  For a nonzero rational projective vector, this
	integer representative is obtained by clearing denominators, dividing all
	coordinates by their greatest common divisor, and choosing the sign so
	that the first nonzero coordinate is positive.  No numerical tolerance or
	random sampling is used.
	
	As a supplementary check, the script
	\texttt{verify\_high\_charge\_q8\_normalized.py} verifies the
	rank of each relevant \(560\)-by-\(98\) integer matrix over
	\(\mathbb Q\): a nonzero vector in its kernel gives rank at most \(97\),
	while rank \(97\) modulo the prime \(1{,}000{,}000{,}007\) exhibits a
	nonzero integer minor and gives rank at least \(97\).  For its integer
	relation matrix, an explicit kernel and a nonzero \(13\)-by-\(13\) minor
	of absolute value \(16\) give rank \(13\); rank \(9\) modulo two then
	forces the nonzero Smith invariant factors to be nine \(1\)'s and four
	\(2\)'s.  These are exact rank and divisibility arguments over the
	integers, not heuristic finite-field specializations.
	
	\subsection{Odd-dihedral and high-charge certificates}
	\label{proofsec:finite-certificates}
	
	\subsubsection{The order-three odd-dihedral six-port certificate}
	\label{proofsubsec:cert-odd-dih-q6}
	
	Fix \(\omega\in\AlgNums\) with \(\omega^2+\omega+1=0\), and identify a
	residue \(r\in\mathbb F_3\) with its representative \(0,1,2\) whenever it
	appears as an exponent \(\omega^r\).  For
	\(\epsilon\in\Ftwo\), let
	\[
	\mathcal B_\epsilon
	=
	\bigl\{\{x,\bar x\}:x\in\Ftwo^4,\
	\wt(x)\equiv\epsilon\pmod2\bigr\}
	\]
	be the four complement blocks of parity \(\epsilon\).  Choose an orientation
	\((x_p,\bar x_p)\), \(0\le p\le3\), of these four blocks, and for
	\(x\in\Ftwo^4\) let \(e_x\in\mathbb C^{\Ftwo^4}\) be the coordinate vector
	given by \(e_x(y)=\one[y=x]\).  For
	\(r,t\in\mathbb F_3\), put
	\begin{equation}
		V_{p,r}=\omega^r e_{x_p}+\omega^{2r}e_{\bar x_p},
		\qquad
		E_{pq}(r,t)=V_{p,r}-V_{q,t}.
		\label{eq:cert-odd-dih-q4-atlas-vectors}
	\end{equation}
	Define the finite projective set
	\begin{equation*}
		\mathcal Q_3^\epsilon
		=
		\{[V_{p,r}]:0\le p\le3,\ r\in\mathbb F_3\}
		\,\cup\,
		\{[E_{pq}(r,t)]:0\le p<q\le3,\ r,t\in\mathbb F_3\}.
	\end{equation*}
	
	For a six-port tensor \(F\), an unordered pair \(i<j\), and
	\(a,b\in\Ftwo\), write \(F_{ij}^{ab}\) for the quaternary slice obtained by
	setting the two selected ports to \(a,b\), with the residual ports kept in
	increasing order.  For \(\xi\in\Ftwo\) and \(s\in\mu_3\), define
	\begin{equation}
		C_{ij,\xi,s}(F)
		=
		F_{ij}^{0,\xi}-sF_{ij}^{1,1-\xi}.
		\label{eq:cert-odd-dih-q6-card}
	\end{equation}
	For \(\eta\in\Ftwo\), if
	\(\supp(F)\subseteq\{x\in\Ftwo^6:\wt(x)\equiv\eta\pmod2\}\), then this
	contraction has residual parity \(\eta\oplus\xi\).
	
	For a support \(S\subseteq\Ftwo^6\), define
	\begin{equation*}
		U_{ij}^{\xi}(S)
		=
		\{x_{\widehat{\{i,j\}}}:x\in S,\ x_i\oplus x_j=\xi\}.
	\end{equation*}
	The \((i,j,\xi)\) \emph{xor sector} consists of the assignments satisfying
	\(x_i\oplus x_j=\xi\).
	We will use the condition that \(U_{ij}^{\xi}(S)\) is empty or is the
	union of one or two complete blocks from a single parity class
	\(\mathcal B_\epsilon\).  In the table below, the middle column lists
	the numbers of complement blocks supporting the three projective points;
	for example, \(1/1/2\) means that two points each use one block and the
	third uses two.  For three distinct blocks \(p,q,u\), the points
	\([E_{pq}(r,t)]\), \([E_{pu}(r,v)]\), and \([E_{qu}(t,v)]\) are collinear
	because \(E_{pq}(r,t)-E_{pu}(r,v)+E_{qu}(t,v)=0\).
	
	We call a pencil \(A-sB\) \emph{projectively nonconstant} when
	\(A,B\) are linearly independent.  Otherwise it is zero or has constant
	projective class wherever it is nonzero.
	Consider labelled pencils \(q_s=A-sB\), \(s\in\mu_3\), modulo one common
	nonzero scalar, such that all three \(q_s\) are nonzero and their three
	distinct projective classes are points of \(\mathcal Q_3^\epsilon\)
	assigned bijectively to the phase labels.  Only pencils satisfying these
	conditions are counted below.
	
	\begin{proposition}[Order-three odd-dihedral six-port classification]
		\label{prop:cert-odd-dih-q6-order-three}
		The following statements hold.
		
		\begin{enumerate}[label=\textnormal{(\roman*)}]
			\item Each set \(\mathcal Q_3^\epsilon\) has \(66\) points.  Its projective
			lines containing at least three of its points are exactly
			\begin{center}
				\small
				\begin{tabular}{ccl}
					\toprule
					number & supported blocks & geometry\\
					\midrule
					\(4\)   & \(1/1/1\) & one fixed complement block\\
					\(36\)  & \(2/2/2\) & one fixed Segre ruling\\
					\(54\)  & \(1/1/2\) & a secant between two blocks\\
					\(108\) & \(2/2/2\) & the displayed relation on three blocks\\
					\bottomrule
				\end{tabular}
			\end{center}
			Every such line contains exactly three points of \(\mathcal Q_3^\epsilon\).
			Hence there are \(202\) such lines and \(202\cdot6=1212\) labelled
			pencils satisfying the stated conditions in each parity.
			
			\item Let \(\eta\in\Ftwo\), let
			\(\supp(F)\subseteq\{x\in\Ftwo^6:\wt(x)\equiv\eta\pmod2\}\), and suppose that every contraction
			in \cref{eq:cert-odd-dih-q6-card} is zero or has projective class in the
			appropriate set \(\mathcal Q_3^{\eta\oplus\xi}\).  For no fixed
			\(i<j\) and \(\xi\) does the projectively nonconstant family
			\(s\mapsto C_{ij,\xi,s}(F)\) lie on a line from the third or fourth row of
			part~\textnormal{(i)}.  Consequently, for every such \(F\), each projectively nonconstant
			family \(s\mapsto C_{ij,\xi,s}(F)\) either is supported on one complement block or
			lies in a fixed Segre ruling.
			
			\item Exactly \(873\) supports \(S\subseteq\Ftwo^6\) make every
			\(U_{ij}^{\xi}(S)\) empty or a union of one or two complete complement
			blocks of a single parity.  They are the empty support; the \(392\)
			distinct cosets, counted without duplication,
			\[
			a+L_\pi,
			\qquad
			L_\pi=\Span_{\Ftwo}\{\one_C:C\in\pi\},
			\]
			where \(a\in\Ftwo^6\) and \(\pi\) is a partition of the six ports into even
			blocks; and the
			\(480\)-element bit-flip/port-permutation orbit of
			\begin{equation*}
				S_\star=
				\{000000,000011,001100,110011,111100,111111\}.
			\end{equation*}
			The distinct affine supports occur in the respective numbers
			\[
			32,\qquad240,\qquad120
			\]
			for partition types \((6)\), \((4+2)\), and \((2+2+2)\), where a partition
			type is the multiset of its block sizes.  Every nonempty
			support in the classification has one parity.
		\end{enumerate}
	\end{proposition}
	
	\begin{proof}
		The vectors in \cref{eq:cert-odd-dih-q4-atlas-vectors} give twelve
		one-block and fifty-four two-block projective points.  Normalize every
		nonzero vector at its first nonzero coordinate.  For every one of the
		\(\binom{66}{2}\) point pairs, take the exact reduced-row-echelon form of
		its two-dimensional span and group equal row spaces.  The groups containing
		at least three points of \(\mathcal Q_3^\epsilon\) all have size three and give the table in
		part~\textnormal{(i)}.
		
		The table also has a direct block interpretation.  The three points on one
		fixed block give four one-block lines.  On a fixed pair of blocks, the nine
		two-block points form a \(3\)-by-\(3\) Segre grid in the transformed
		phase coordinates \(u=r+t\), \(v=r-t\) in \(\mathbb F_3\).
		Fixing \(u\) or \(v\), not the original \(r\) or \(t\) in
		\cref{eq:cert-odd-dih-q4-atlas-vectors}, gives its two rulings and hence
		\(6\cdot2\cdot3=36\) lines.  Every \(E_{pq}(r,t)\) lies on the secant
		through \(V_{p,r}\) and \(V_{q,t}\), giving \(6\cdot3^2=54\) secants.
		For three distinct blocks,
		\[
		E_{pq}(r,t)-E_{pu}(r,v)+E_{qu}(t,v)=0
		\]
		gives \(\binom43 3^3=108\) lines on three blocks.  Conversely, in a dependent
		triple every block occurring in one vector must occur in another.  According
		as the triple uses one, two, or three blocks, it is one of the displayed
		types; four blocks would leave a block occurring only once.
		
		Assigning the phases \(1,\omega,\omega^2\) to the three points on one of these \(202\) lines gives a
		unique affine pencil up to one common nonzero scalar.  Indeed,
		after the unique dependence is scaled, solve
		\[
		q_1=A-B,\qquad q_\omega=A-\omega B
		\]
		and verify \(q_{\omega^2}=A-\omega^2B\).  All six phase assignments are
		distinct.  This proves the last assertion of part~\textnormal{(i)}.
		
		For part~\textnormal{(ii)}, after a bit flip to even parity, fix
		\(x_0=0000\), \(x_1=0011\), \(x_2=0101\), and \(x_3=0110\).  For
		\(t=(t_1,\ldots,t_4)\in\mathbb F_3^4\), define
		\((R_tq)(x)=\omega^{\sum_{j=1}^4t_jx_j}q(x)\).  Under \(R_t\),
		the four oriented block labels change by \(Ht\),
		where
		\[
		H=
		\begin{pmatrix}
			1&1&1&1\\
			1&1&-1&-1\\
			1&-1&1&-1\\
			1&-1&-1&1
		\end{pmatrix},
		\qquad
		\det H=-16\ne0\pmod3.
		\]
		The phase permutations are generated by \(s\mapsto\alpha s\) and
		\(s\mapsto\alpha/s\), with
		\(\alpha\in\mu_3\); these maps form \(D_6\cong S_3\) and realize all six
		permutations of \(\mu_3\).  Together with port permutations and bit flips,
		these actions reduce the two-block secant and three-block cases to the following
		representative pencils for a fixed restriction \(x_i\oplus x_j=\xi\).
		A semicolon separates
		the four oriented complement blocks, and \(0\) denotes the pair \((0,0)\):
		\begin{align*}
			A_{\rm sec}&=(1,1;\omega,\omega;0;0),&
			B_{\rm sec}&=(\omega^2,\omega^2;\omega,\omega;0;0),
			\\
			A_{\rm tri}&=(1,1;\omega,\omega;\omega^2,\omega^2;0),&
			B_{\rm tri}&=(\omega,\omega;1,1;\omega^2,\omega^2;0).
		\end{align*}
		Each displayed pair \((A_*,B_*)\) denotes the pencil
		\(q_s=A_*-sB_*\), \(s\in\mu_3\).
		Since \(H\) is invertible over \(\mathbb F_3\), the maps \(R_t\) shift the
		block labels independently.  Port permutations and bit flips act transitively
		on the block choices and the two total parities.  Together with the phase permutations, this proves that every orbit of
		the two-block secant and three-block cases under these operations is
		represented.
		
		Fix the selected pair \(i<j\) and sector \(\xi\).  It remains to exclude
		six-port tensors \(F\) satisfying the hypotheses of part~\textnormal{(ii)}
		and realizing either selected representative pencil.  Put
		\(\zeta=\xi\oplus1\).  Under the restriction \(x_i\oplus x_j=\zeta\),
		a projectively nonconstant pencil is one
		of the \(1212\) labelled pencils from
		part~\textnormal{(i)}.
		Its slices \(C=F_{ij}^{0,\zeta}\) and \(D=F_{ij}^{1,1-\zeta}\) are fixed
		up to a common scalar \(\lambda\in\AlgNums^\times\).  For each of the
		fourteen port pairs other than the selected pair, and for each xor sector and
		phase, substitution into \cref{eq:cert-odd-dih-q6-card} gives a quaternary
		contraction \(q_0+\lambda q_1\), where the computed vectors
		\(q_0,q_1\in\AlgNums^{\Ftwo^4}\) are independent of \(\lambda\).  For a
		fixed representative \(v\in\AlgNums^{\Ftwo^4}\setminus\{0\}\) of a point
		of the appropriate set \(\mathcal Q_3^\epsilon\), the set of \(\lambda\)
		for which this vector is zero or belongs to
		\(\Span\{v\}\) is all of the field, empty, or one
		point; exact two-column row reduction determines which case occurs.
		Intersecting the resulting finite unions over all remaining pairs, both xor
		sectors, and all three phases leaves no nonzero \(\lambda\) for either
		representative.
		
		For an opposite pencil with constant projective class, write \(C=cv\) and
		\(D=dv\), where \(v\in\AlgNums^{\Ftwo^4}\setminus\{0\}\) represents one
		of the \(66\) points of the appropriate \(\mathcal Q_3^\epsilon\) and \((c,d)\in\AlgNums^2\setminus\{(0,0)\}\).
		For each remaining contraction, substitution gives
		\(q_0+cq_1+dq_2\), where the computed
		\(q_0,q_1,q_2\in\AlgNums^{\Ftwo^4}\) are independent of \(c,d\).  Its
		preimage over all target points of \(\mathcal Q_3^\epsilon\) is a finite union of affine linear
		systems in \((c,d)\).  Exact RREF,
		followed by intersection over all contractions in the hypothesis of
		part~\textnormal{(ii)}, leaves no
		nonzero pair \((c,d)\).  The identically zero opposite sector is checked
		separately and also fails.  All matrices have entries in
		\(\mathbb Q(\omega)\); their ranks and consistency are unchanged after
		extension to \(\AlgNums\).  Reversing the two ports when \(\xi=0\) changes nothing, whereas for
		\(\xi=1\) reversal at \(s\) gives a nonzero scalar multiple of the contraction
		at \(s^{-1}\), so unordered port pairs exhaust the ordered conditions.
		
		For part~\textnormal{(iii)}, fix one port pair and one xor sector.  A union satisfying the stated condition is empty, one of the eight one-block unions across
		the two parities, or one of the twelve two-block unions.  Each residual word
		in a nonempty union occurs only in \(F_{ij}^{0,\xi}\), only in
		\(F_{ij}^{1,1-\xi}\), or in both.  Hence the exact finite superset of local candidates has
		\[
		1+8\cdot3^2+12\cdot3^4=1045
		\]
		members.  Combine the two disjoint xor sectors on the first port pair,
		giving \(1045^2\) candidates.  Retain a candidate precisely when, on
		each of the other fourteen pairs, each of the two unions is empty or is
		the union of one or two complete complement blocks of its own parity.  Encode a support by a
		\(64\)-bit integer whose bit \(k\) corresponds to the \(k\)-th word in
		lexicographic order on \(\{0,1\}^6\), starting at \(k=0\).  Exact
		intersection of these encoded support lists leaves \(873\) distinct supports.
		
		Independently generate every partition into nonempty even blocks and every
		coset of its block-indicator space.  This gives
		\[
		32+240+120=392
		\]
		distinct supports.  Generating all bit flips and port permutations of
		\(S_\star\) gives \(480\) further supports.  These two families and the
		empty support are disjoint, and their union is exactly the \(873\)-support
		intersection.  Since every block indicator has even weight and
		\(S_\star\) has even parity, every nonempty support has one parity.
		
		The exact calculation constructs the sets \(\mathcal Q_3^\epsilon\), their lines containing at least three points, labelled pencils,
		affine preimages, local support lists, partition cosets, and exceptional
		orbit from the displayed definitions.  It loads no precomputed data or
		stored accepted-case list and uses no random choice, floating-point comparison,
		finite-field specialization, or external solver.
	\end{proof}
	
	\begin{remark}[Scope of the odd-dihedral certificate]
		\label{rem:cert-odd-dih-q6-scope}
		\Cref{prop:cert-odd-dih-q6-order-three} is a normalized finite statement.
		The algebraic quaternary classification, the construction of the realizable
		kernels in \cref{eq:odd-dih-kernels}, and the reduction to
		\cref{eq:cert-odd-dih-q6-card} remain in the main proof.  The certificate does
		not prove the \(N\ge5\) support bound, the exclusions of contractions supported on two points, the
		coefficient factorization for support that selects one endpoint of each
		matching edge, the arbitrary-arity common-factor argument, or the reduction
		for the entire signature set.
	\end{remark}
	
	\Needspace{7\baselineskip}
	\begin{theorem}[Six-port support classification]
		\label{thm:cert-odd-dih-q6-interface}
		The order-three classification has the following consequences.
		\begin{enumerate}[label=\textnormal{(\roman*)}]
			\item Let \(F\) be a six-port tensor of one parity, and suppose that every
			contraction \(C_{ij,\xi,s}(F)\), for every port pair, xor, and
			\(s\in\mu_3\), is zero or belongs to \(\mathcal B_3([4])\).
			Then, for every fixed pair and xor, with
			\(A=F_{ij}^{0,\xi}\), \(B=F_{ij}^{1,1-\xi}\),
			\(\supp(A)\cup\supp(B)\) is empty or is the union of at most two
			complement pairs in its residual parity.
			\item Let \(F\) be a six-port tensor of one parity.  If the support union
			in part~\textnormal{(i)} is empty or consists of at most two
			complement pairs for every pair and xor, then \(\supp(F)\) is empty,
			a partition coset of type \((6)\), \((4+2)\), or \((2+2+2)\), or a
			bit-flip/port-permutation image of \(S_\star\) in
			\cref{eq:odd-dih-q6-nonproduct-support}.
		\end{enumerate}
	\end{theorem}
	
	\begin{proof}
		For part~\textnormal{(i)}, every nonzero contraction over all port pairs and xors
		has projective class in the corresponding set \(\mathcal Q_3^\epsilon\).  By
		\cref{prop:cert-odd-dih-q6-order-three}\textnormal{(ii)}, a projectively nonconstant
		pencil lies on a one-block line or one fixed Segre ruling.  Its support
		union consequently contains at most two complete complement pairs.
		Coordinatewise,
		\[
		\supp(A)\cup\supp(B)=\bigcup_{s\in\mu_3}\supp(A-sB),
		\]
		since a nonzero affine coefficient has at most one root.  This also handles
		zero and projectively constant pencils.  For part~\textnormal{(ii)}, the
		hypothesis says precisely that every local union
		\(U_{ij}^{\xi}(\supp(F))\) is empty or consists of one or two complete
		complement blocks of a single parity, as required in the enumeration of
		\cref{prop:cert-odd-dih-q6-order-three}\textnormal{(iii)}.  That enumeration
		gives the asserted alternatives.  Neither statement asserts coefficient
		factorization merely from the support shape; the required coefficient
		argument and the exclusions of contractions supported on two points are proved in
		\cref{thm:odd-dih-q6-base}.
	\end{proof}
	
	\subsubsection{High-charge six-port loop and two-copy arrays}
	\label{proofsubsec:cert-high-charge-q6}
	
	Let
	\[
	\mathcal W_4=(0011,0101,0110,1001,1010,1100).
	\]
	The arrays below use this order.  The Python loop generator lists its
	entries in reverse order; both
	\eqref{eq:cert-q6-zero-pairs} and \eqref{eq:cert-q6-pair-products} are
	invariant under \(z_k\mapsto z_{5-k}\).  The two-copy generator uses
	\(\mathcal W_4\) directly.
	Let \(U=[6]\), let \(t\in\mathbb C^\times\), and assign arbitrary complex
	coefficients \(q_S\) to the twenty triples \(S\in\binom U3\).  For each
	deleted pair \(i<j\), order the residual ports as
	\(v_0<v_1<v_2<v_3\); for \(w\in\mathcal W_4\), write
	\(w=(w_0,w_1,w_2,w_3)\) and put
	\(R_w=\{v_r:w_r=1\}\), and define the six EO entries
	\begin{equation}
		L_{ij}(w)=q_{R_w\cup\{i\}}+q_{R_w\cup\{j\}},
		\qquad w\in\mathcal W_4.
		\label{eq:cert-q6-loop-array}
	\end{equation}
	For unordered two-element external sets \(A,B\in\binom U2\), use the
	canonical increasing orders \(A=(a_0,a_1)\), \(B=(b_0,b_1)\) and the
	external order \((a_0,a_1,b_0,b_1)\).  The first and second copies retain
	their first/second copy labels.  Put \(I_A=U\setminus A\),
	\(I_B=U\setminus B\), and let
	\(\pi:I_A\to I_B\) be a bijection.  Define the six EO entries, indexed by
	the concatenated words \(xy\in\mathcal W_4\), where
	\(x=(x_0,x_1)\) and \(y=(y_0,y_1)\), by
	\begin{equation}
		C_{A,B,\pi}(x,y)
		=t\,\one[(x,y)\in\{(00,11),(11,00)\}]
		+\sum_{\substack{T\subseteq I_A\\|T|=3-\wt(x)}}
		q_{A_x\cup T}\,
		q_{B_y\cup(I_B\setminus\pi(T))},
		\label{eq:cert-q6-two-copy-array}
	\end{equation}
	where \(A_x=\{a_r:x_r=1\}\), and similarly for \(B_y\).
	
	For any quaternary EO array
	\(z=(z_0,z_1,z_2,z_3,z_4,z_5)\) indexed in the order \(\mathcal W_4\),
	put \(P_0=z_0z_5\), \(P_1=z_1z_4\), and \(P_2=z_2z_3\).
	
	\begin{proposition}[High-charge six-port coefficient classification]
		\label{prop:cert-high-charge-q6-joint-deck}
		Suppose every one of the fifteen arrays in
		\cref{eq:cert-q6-loop-array} and all
		\(15^2\cdot4!=5400\) labelled arrays in
		\cref{eq:cert-q6-two-copy-array} satisfy
		\begin{align}
			z_0=0&\Longleftrightarrow z_5=0,&
			z_1=0&\Longleftrightarrow z_4=0,&
			z_2=0&\Longleftrightarrow z_3=0,
			\tag{Q6.6}\label{eq:cert-q6-zero-pairs}\\
			P_0P_1P_2&=0,&
			P_iP_j(P_i^2-P_j^2)&=0\quad(0\le i<j\le2).
			\tag{Q6.7}\label{eq:cert-q6-pair-products}
		\end{align}
		Then
		\[
		\{R\in\tbinom U3:q_R\ne0\}\subseteq\{S,U\setminus S\}
		\qquad\text{for some }S\in\binom U3.
		\]
		This holds for arbitrary complex \(q_S\) and every \(t\ne0\).
	\end{proposition}
	
	\begin{proof}
		For each of the fifteen loop arrays,
		\cref{eq:cert-q6-zero-pairs,eq:cert-q6-pair-products} force at least one
		of its three complementary EO-entry pairs to vanish.  Selecting such a
		pair gives, for a suitable \(S\in\binom U3\) containing exactly one of
		\(i,j\), two signed equations
		\begin{equation*}
			q_S=-q_{S\mathbin\triangle\{i,j\}},
			\qquad
			q_{U\setminus S}=-q_{U\setminus(S\mathbin\triangle\{i,j\})}.
		\end{equation*}
		Thus every coefficient vector lies in at least one of exactly
		\(3^{15}=14\,348\,907\) branches defined by signed linear equations.  Signed
		union--find deduplicates them to \(1868\) linear subspaces; modulo the \(S_6\) port action
		and simultaneous complementation, they form exactly \(24\) orbits.
		For the C++ enumeration, number ports from \(0\) to \(5\) and order
		triple coordinates by increasing \(\sum_{k\in S}2^k\).  In each nonzero
		component \(C\) of the signed equations, write its coordinates as
		\(q_i=\varepsilon_i u_C\), where \(\varepsilon_i\in\{1,-1\}\),
		choose \(\varepsilon_{\min C}=1\), and use \(\min C\) to name
		the free parameter \(u_C\).
		After this sign normalization, choose the lexicographically least encoding
		under port permutations and simultaneous complementation, taking exactly
		one representative per orbit.
		
		For each family, substitute independent complex parameters into the
		displayed polynomials
		\cref{eq:cert-q6-loop-array,eq:cert-q6-two-copy-array} and impose only
		\cref{eq:cert-q6-zero-pairs,eq:cert-q6-pair-products}.  Split every complex
		parameter into real and imaginary parts.  A selected nonzero parameter may
		be normalized to \(1\), because all equations are homogeneous under
		\((q,t)\mapsto(cq,c^2t)\) for \(c\in\mathbb C^\times\).  Exact satisfiability checks for the resulting real polynomial systems
		find no solution with two distinct nonzero coefficient pairs
		\(\{q_S,q_{U\setminus S}\}\).
		
		The resulting exhaustive calculation is
		\[
		14{,}348{,}907\text{ branches}
		\longrightarrow 1868\text{ signed spaces}
		\longrightarrow 24\text{ orbits}.
		\]
		All generated expressions have integer coefficients, and all reductions
		performed before the solver use exact rational polynomial arithmetic.  The
		all-zero orbit is immediate.
		The exact real-polynomial queries for the remaining twenty-three nonzero
		orbit families were evaluated with Z3~4.14.0.  The answers are
		\texttt{UNSAT}, meaning that no satisfying assignment exists. The
		calculation uses 61 queries: 51 for projective charts of sparse orbits
		1--19 and 10 for the implications used in dense orbits 20--23.
		For sparse orbits, the script checks that parameter sets omitted by the
		chosen normalizations are supported on at most one complementary block and hence cannot meet the counterexample condition.  The dense
		branches explicitly reduce to a previously checked sparse family, the zero
		vector, or one complementary block.  No parameter sample replaces a
		quantified polynomial query.
		
		The six-port wrapper requires Z3~4.14.0 and records the compiler and
		solver versions and source-file hashes. Before any solver query, the
		Python stage reindexes the C++ output and checks that both stages give
		exactly the same orbit representatives. The programs regenerate the
		orbit representatives and exact real-polynomial queries on each run.
		The suite option \texttt{--q6-audit-dir} saves the generated inputs,
		parameter encoding, branch labels, and solver responses for inspection;
		the package README gives the reproduction commands.

		The runs neither request nor check solver proof objects.
		The \texttt{UNSAT} answers therefore rely on the correctness of Z3;
		exporting or rerunning the inputs does not remove this dependence.
		Subject to this dependence, the exhaustive calculation establishes
		the proposition over arbitrary complex coefficients.
	\end{proof}
	
	\begin{remark}[Scope of the six-port certificate]
		\label{rem:cert-high-charge-q6-scope}
		\Cref{prop:cert-high-charge-q6-joint-deck} starts from the coefficient
		arrays \cref{eq:cert-q6-loop-array,eq:cert-q6-two-copy-array}.  It does not
		prove that these arrays are direct loop or two-copy gadgets, does not
		make any summand of \cref{eq:cert-q6-two-copy-array} separately realizable,
		and does not establish a reduction for the entire signature set.  Those
		arguments remain in \cref{subsec:hc-q6-base,thm:hc-q6-complete}.
	\end{remark}
	
	\begin{theorem}[Consequence for loop and two-copy arrays]
		\label{thm:cert-high-charge-q6-interface}
		Let \(U=[6]\), let \(t\in\mathbb C^\times\), and let
		\((q_S)_{S\in\binom U3}\) be arbitrary complex coefficients.  Form the
		fifteen loop arrays in \cref{eq:cert-q6-loop-array} and all \(5400\)
		labelled two-copy arrays in \cref{eq:cert-q6-two-copy-array}.  If every one
		of these arrays belongs to \(\mathcal B_\times^0([4])\), then
		\[
		\{R\in\tbinom U3:q_R\ne0\}
		\subseteq\{S,U\setminus S\}
		\qquad\text{for some }S\in\binom U3.
		\]
	\end{theorem}
	
	\begin{proof}
		Every member of \(\mathcal B_\times^0([4])\) satisfies
		\cref{eq:cert-q6-zero-pairs,eq:cert-q6-pair-products}; hence
		\cref{prop:cert-high-charge-q6-joint-deck} applies and yields the assertion.
	\end{proof}
	
	\begin{remark}[Why the two-copy six-port arrays are necessary]
		\label{rem:cert-high-charge-q6-one-loop-insufficient}
		Order each triple as \(S=\{s_0<s_1<s_2\}\), and order the twenty triples by
		the lexicographic order of \((s_0,s_1,s_2)\).  In that order, the
		full-support coefficient vector
		\[
		(1,-1,-1,1,i,-i,-1,1,-i,-i,-i,-i,-1,1,-i,i,-1,1,1,-1)
		\]
		does not define a member of \(\mathcal B_\times([6])\), although all fifteen
		arrays \cref{eq:cert-q6-loop-array} lie in
		\(\mathcal B_\times^0([4])\).  For this
		same vector, one labelled array \cref{eq:cert-q6-two-copy-array} is
		\[
		(t,2,-2,-2,2,t),
		\]
		where \(t\ne0\); its six nonzero entries put it outside that class.  Thus
		the one-loop conditions alone cannot imply
		the conclusion of \cref{thm:cert-high-charge-q6-interface}.
	\end{remark}
	
	\subsubsection{Eight-port matching families}
	\label{proofsubsec:cert-q8-matching-atlas}
	
	Let \(V=\{0,1,\ldots,7\}\), and for distinct labels \(a,b\) write
	\(ab\) for the unordered pair \(\{a,b\}\).  For every pair
	\(e\in\binom V2\), let
	\(P_e\) be a perfect matching of \(V\setminus e\).  For disjoint
	\(e,r\in\binom V2\), write \(r=\{u,v\}\) and define \(\kappa_r(P_e)\) as follows.  If
	\(r\in P_e\), delete that edge.  Otherwise let \(u',v'\) be the mates of
	\(u,v\), remove \(uu'\) and \(vv'\), add \(u'v'\), and retain the unique
	untouched edge.  In both cases \(\kappa_r(P_e)\) is a two-edge matching on
	\(V\setminus(e\cup r)\).
	
	\begin{proposition}[Eight-port matching classification]
		\label{prop:cert-q8-collision-free-atlas}
		Suppose
		\begin{equation}
			\kappa_r(P_e)=\kappa_e(P_r)
			\qquad(e\cap r=\varnothing)
			\label{eq:cert-q8-commuting-cards}
		\end{equation}
		for every pair of disjoint deleted pairs.  After relabeling so that
		\(P_{01}=\{23,45,67\}\), exactly nine labelled families satisfy these
		constraints.  Seven are induced by a global perfect matching \(M\) on eight
		ports, meaning \(P_e=\kappa_e(M)\) under the same deletion-and-splicing
		rule.  Each of these seven families has \(P_A\cap P_B\ne\varnothing\)
		for some distinct pairs \(A,B\) with \(A\cap B\ne\varnothing\).
		With labels \(0,\ldots,7\) interpreted as three-bit vectors, one remaining
		family is given by the matchings below and the other is obtained by swapping
		\(6,7\), where
		the two-character string \(ab\) abbreviates the unordered pair \(\{a,b\}\):
		\begin{equation*}
			\begin{array}{cc}
				\toprule
				h&\{\{t,t+h\}:t\in\mathbb F_2^3\}\\ \midrule
				001&01,23,45,67\\
				010&02,13,46,57\\
				011&03,12,47,56\\
				100&04,15,26,37\\
				101&05,14,27,36\\
				110&06,17,24,35\\
				111&07,16,25,34.\\
				\bottomrule
			\end{array}
			\tag{Q8.14b}
			\label{eq:cert-q8-translation-atlas}
		\end{equation*}
		For a deleted pair \(p,q\), removing the edge \(pq\) from the row
		\(h=p+q\) gives \(P_{pq}\).
	\end{proposition}
	
	\begin{proof}
		Once \(P_{01}\) is fixed, each of the other twenty-seven port pairs has the
		fifteen perfect matchings on its six residual ports as candidates.
		An exhaustive backtracking search imposes
		\cref{eq:cert-q8-commuting-cards} whenever both sides become known.  The
		search has \(236\) search nodes and exactly nine leaves.  Direct
		inspection of each leaf gives the stated \(7+2\) split and the table in
		\cref{eq:cert-q8-translation-atlas}.
		
		The scripts \texttt{verify\_high\_charge\_q8\_atlas.py} and
		\texttt{verify\_high\_charge\_q8\_normalized.py} reconstruct the same nine
		leaves from the displayed definitions, and the latter also checks the
		\(236\)-node count and the linear systems for the original tensor's
		seventy weight-four coefficients and twenty-eight contraction scales.
		The two programs reuse the same matching-generation logic and are
		therefore not independent implementations. This proposition is only
		the finite matching classification.  The
		direct double-contraction
		identity that yields \cref{eq:cert-q8-commuting-cards}, and the reduction
		that excludes families with \(P_A\cap P_B\ne\varnothing\) for distinct
		\(A,B\) satisfying \(A\cap B\ne\varnothing\), remain in the main proof.
	\end{proof}
	
	\begin{theorem}[Translation form of the matching family]
		\label{thm:cert-q8-collision-free-interface}
		Let \(V\) be an eight-element port set and, for every
		\(e\in\binom V2\), let \(P_e\) be a perfect matching of \(V\setminus e\).
		Suppose
		\[
		\kappa_r(P_e)=\kappa_e(P_r)
		\qquad(e\cap r=\varnothing).
		\]
		If \(P_A\cap P_B=\varnothing\) whenever \(A\ne B\) and
		\(A\cap B\ne\varnothing\), then there is a bijection
		\(\iota:V\to\mathbb F_2^3\) such that, after relabelling, for \(p\ne q\)
		and \(h=p+q\),
		\[
		P_{pq}
		=
		\bigl\{\{t,t+h\}:t\in\mathbb F_2^3\setminus\{p,q\}\bigr\},
		\]
		where duplicate unordered pairs are suppressed.
	\end{theorem}
	
	\begin{proof}
		Relabel \(V\) so that \(P_{01}=\{23,45,67\}\).  By
		\cref{prop:cert-q8-collision-free-atlas}, seven of the nine labelled
		families violate the condition on \(P_A\cap P_B\).  The remaining two
		differ only by a relabelling and give the displayed matching family.
	\end{proof}
	
	\begin{proposition}[Coefficients on affine planes]
		\label{prop:cert-q8-canonical-support}
		Identify \(V=\mathbb F_2^3\), with the labels \(0,\ldots,7\) interpreted as
		their three-bit binary expansions, and for
		\(S\in\binom V4\) put
		\(\sigma(S)=\bigoplus_{s\in S}s\).  Let \(c_S\in\mathbb C\), and let
		\(\lambda_{pq}\in\mathbb C^\times\) for every unordered pair
		\(\{p,q\}\in\binom V2\).  Suppose that, whenever \(p\in S\), \(q\notin S\),
		\(h=p+q\), and \(S'=S\mathbin\triangle\{p,q\}\),
		\begin{equation*}
			c_S+c_{S'}=
			\begin{cases}
				\lambda_{pq},&\sigma(S)\in\{0,h\},\\
				0,&\sigma(S)\notin\{0,h\}.
			\end{cases}
			\tag{Q8.19a}\label{eq:cert-q8-canonical-swap-system}
		\end{equation*}
		Then the four-sets with nonzero coefficients are exactly the fourteen
		affine planes of \(V\), and there is one \(c\in\mathbb C^\times\) such that
		\begin{equation*}
			c_S=
			\begin{cases}
				c,&\sigma(S)=0,\\
				0,&\sigma(S)\ne0.
			\end{cases}
			\tag{Q8.19b}\label{eq:cert-q8-canonical-coefficients}
		\end{equation*}
		Call a four-subset a \emph{nonplane} when it is not an affine plane.
		Let \(G\) be the graph on the fifty-six nonplanes, joining \(S\) to
		\(S\mathbin\triangle\{p,q\}\) whenever \(p\in S\), \(q\notin S\), and
		the latter four-set is also a nonplane.  Let \(H\) be the bipartite
		incidence graph between planes and unordered port pairs, with an edge
		between \(S\) and \(\{p,q\}\) exactly when
		\(\lvert S\cap\{p,q\}\rvert=1\).  Then \(G\) is connected and contains
		an odd cycle, and \(H\) is connected.
	\end{proposition}
	
	\begin{proof}
		A four-set has xor zero exactly when it is an affine plane: xor zero makes
		the fourth point the affine sum of the other three, and the seven nonzero
		linear functionals on \(\mathbb F_2^3\) each have two affine level-set
		planes.  Thus there are fourteen such sets and fifty-six nonplanes.
		For a nonplane \(S\), the exchanged set \(S'=S\mathbin\triangle\{p,q\}\)
		is again a nonplane exactly when \(\sigma(S)\notin\{0,p+q\}\).
		Lift \(t\in V\) to
		\(\widehat t=(1,t)\in\mathbb F_2^4\), and let \(M\) be the represented
		matroid on ground set \(V\) with representing vectors \(\widehat t\).  A four-set is a nonplane exactly
		when its four lifted points form a basis.  Thus \(G\) is the
		basis-exchange graph of the rank-four binary matroid \(M\), whose
		vertices are bases and whose edges exchange one element while preserving
		a basis.  This graph is connected.
		
		In binary labelling it contains the triangle below, where, for example,
		\(0124\) abbreviates the four-set \(\{0,1,2,4\}\):
		\begin{equation*}
			\mathtt{0124}\;\text{---}\;\mathtt{0125}\;\text{---}\;
			\mathtt{0245}\;\text{---}\;\mathtt{0124}.
			\tag{Q8.20}\label{eq:cert-q8-bad-swap-cycle}
		\end{equation*}
		The three xor sums are \(7,6,3\), while the corresponding swap directions
		are \(1,5,4\); hence all three edges belong to \(G\).  The second
		line of \cref{eq:cert-q8-canonical-swap-system} gives, on every edge of \(G\),
		\begin{equation*}
			c_{S'}=-c_S.
		\end{equation*}
		Alternating around
		\cref{eq:cert-q8-bad-swap-cycle} and using characteristic zero forces the
		three coefficients on the triangle to vanish.  Connectivity then forces
		\(c_S=0\) for every nonplane.
		
		Let \(S\) now be a plane and let \(p\in S,q\notin S\).  Then
		\(S'=S\mathbin\triangle\{p,q\}\) has xor \(p+q\ne0\), so it is a nonplane.
		The first line of \cref{eq:cert-q8-canonical-swap-system} therefore gives
		\begin{equation}
			c_S=\lambda_{pq}\ne0.
			\label{eq:cert-q8-plane-pair-value}
		\end{equation}
		It remains only to show that these values agree.  Regard a plane as
		incident with every unordered pair having one endpoint in the plane and
		one outside it.  Two nonparallel affine planes meet in two points and their
		union has six points; a pair from their intersection to the complement of
		their union is incident with both planes in \(H\).  Two distinct
		parallel planes are complementary and hence have the same incident port
		pairs.  Thus \(H\) is connected.  Every pair \(\{p,q\}\) is incident
		with a plane:
		choose a linear functional \(\ell\) with \(\ell(p+q)=1\) and take
		\(\{x:\ell(x)=\ell(p)\}\).
		\Cref{eq:cert-q8-plane-pair-value} makes all fourteen plane coefficients
		and all incident \(\lambda_{pq}\)'s equal to one common nonzero \(c\),
		proving \cref{eq:cert-q8-canonical-coefficients}.
		
		The script \texttt{verify\_high\_charge\_q8\_atlas.py} also checks the
		\(14+56\) split, the connectivity of \(G\), the displayed odd cycle,
		and the connectivity of \(H\). These are supplementary checks of
		the argument above.
	\end{proof}
	
	\begin{theorem}[Affine-plane coefficients]
		\label{thm:cert-q8-canonical-support-interface}
		Identify \(V=\mathbb F_2^3\).  Let
		\(c_S\in\mathbb C\) for \(S\in\binom V4\), and let
		\(\lambda_{pq}\in\mathbb C^\times\) for every unordered pair
		\(\{p,q\}\in\binom V2\).  Suppose that, whenever \(p\in S\), \(q\notin S\),
		\(h=p+q\), and \(S'=S\mathbin\triangle\{p,q\}\),
		\[
		c_S+c_{S'}=
		\begin{cases}
			\lambda_{pq},&\displaystyle\bigoplus_{s\in S}s\in\{0,h\},\\
			0,&\displaystyle\bigoplus_{s\in S}s\notin\{0,h\}.
		\end{cases}
		\]
		Then there is one \(c\in\mathbb C^\times\) such that
		\[
		c_S=
		\begin{cases}
			c,&\displaystyle\bigoplus_{s\in S}s=0,\\
			0,&\displaystyle\bigoplus_{s\in S}s\ne0.
		\end{cases}
		\]
		Equivalently, the support consists exactly of the fourteen affine planes of
		\(\mathbb F_2^3\), and all supported coefficients are equal.
	\end{theorem}
	
	\begin{proof}
		This follows from
		\cref{prop:cert-q8-canonical-support}.
	\end{proof}
	
	\bibliographystyle{alpha}
	\bibliography{paper2-full-holant-dichotomy}
	
\end{document}
